Adaptive regularization with cubics on manifolds
%!TEX root = ./arc_on_manifolds.tex

%\input{responsetoreviewsround1}

\maketitle

\begin{abstract}
Adaptive regularization with cubics (ARC) is an algorithm for unconstrained, non-convex optimization. Akin to the trust-region method, its iterations can be thought of as approximate, safe-guarded Newton steps. For cost functions with Lipschitz continuous Hessian, ARC has optimal iteration complexity, in the sense that it produces an iterate with gradient smaller than $\varepsilon$ in $O(1/\varepsilon^{1.5})$ iterations. For the same price, it can also guarantee a Hessian with smallest eigenvalue larger than $-\sqrt{\varepsilon}$. In this paper, we study a generalization of ARC to optimization on Riemannian manifolds. In particular, we generalize the iteration complexity results to this richer framework. Our central contribution lies in the identification of appropriate manifold-specific assumptions that allow us to secure these complexity guarantees both when using the exponential map and when using a general retraction. A substantial part of the paper is devoted to studying these assumptions---relevant beyond ARC---and providing user-friendly sufficient conditions for them.
%These may prove useful to study various other optimization algorithms on manifolds.
Numerical experiments are encouraging. %, though the Riemannian trust-region method (whose worst-case complexity is worse) still has the upper hand in practice.
\keywords{Optimization on manifolds \and Complexity \and Lipschitz regularity \and Cubic regularization \and Newton's method}
\end{abstract}

%\TODO{Should we say more about~\cite{ring2012optimization}?} % They also talk about uniform Lipschitz assumptions on pullbacks, and they also talk about disentangling those conditions between $f$ and the retraction; they specifically make that statement (without proof) for first-order and Exp map: see their Remark 2. In their Props.~5 and~7, they have a uniform Lipschitz condition on the Hessian of the pullback. Note: their analysis allows for infinite dimensional manifolds. Remark 3 deals with second order unentangling. Read that Remark 3 again, very carefully.

\section{Introduction} \label{sec:intro}

%\TODO{Reviewer 3 from NIPS states it is easy to get the impression, because of our writing, that this paper is just a generalization to manifolds for the sake of generalizing to manifolds, without clear merits. Let's try to act on that comment by reconsidering our claims to fame and stressing them better.}

Adaptive regularization with cubics (ARC) is an iterative algorithm used to solve unconstrained optimization problems of the form
\begin{align*}
	\min_{x \in \Rn} \ f(x),
\end{align*}
where $f \colon \Rn \to \reals$ is twice continuously differentiable~\citep{griewank1981modification}. Given any initial iterate $x_0 \in \Rn$, assuming $f$ is lower-bounded and has a Lipschitz continuous Hessian, ARC produces an iterate $x_k$ with small gradient, namely, $\|\nabla f(x_k)\| \leq \varepsilon$, in at most $O(1/\varepsilon^{1.5})$ iterations~\citep{nesterov2006cubic,cartis2011adaptivecubic,bcgmt}. This improves upon the worst-case iteration complexity of steepest descent and classical trust-region methods.
In fact, this iteration complexity is optimal under those assumptions~\citep{carmon2017lower},
contributing to renewed interest in this method.

In this paper, we study a generalization of ARC to optimization on manifolds, that is, % problems of the form
\begin{align}
	\min_{x \in \calM} \ f(x),
	\label{eq:P}
	\tag{P}
\end{align}
where $\calM$ is a given Riemannian manifold
% I hesitated to add here Hausdorff and second-countable "to be safe", but it's probably silly: authors usually define Riemannian manifolds as smooth manifolds + a metric, and they usually define smooth manifolds as Hausdorff spaces with a complete atlas (see for example O'Neill, p3). So Hausdorff is included already. As for second-countability, it seems to be included as well: see the accepeted answer: /connected/ Hausdorff Riemannian manifolds second countable: https://math.stackexchange.com/questions/2131530/why-is-important-for-a-manifold-to-have-countable-basis. Note that here the writer specifies Hausdorff and Riemannian, but the whole question is specifically about these provisions, so perhaps they want to be explicit. In any case: it seems rather typical (I think?) for writers to omit this assumption. If we were general enough not to need these conditions, we would likely say so explicitly. But see also O'Neill pp21-22: they say explicitly there that they assume second-countability.
and $f \colon \calM \to \reals$ is a (sufficiently smooth) cost function.
The practical interest in optimization on manifolds stems from its ubiquity: it comes up naturally in numerical linear algebra (spectral decompositions, low-rank Lyapunov equations), signal and image processing (shape analysis, diffusion tensor imaging, community detection on graphs, rotational video stabilization), statistics and machine learning (matrix/tensor completion, metric learning, Gaussian mixtures, activity recognition, independent component analysis), robotics and computer vision (simultaneous localization and mapping, structure from motion, pose estimation) and various other fields.
The theoretical interest comes from the fact that Riemannian geometry is arguably the ``right'' setting for \emph{unconstrained} optimization---indeed, it is the minimal mathematical structure required to have comfortable notions of gradients and Hessians, which are the basic building blocks of smooth, unconstrained optimization algorithms.
See for example~\citep{AMS08} and~\citep{boumal2020intromanifolds} for book-length introductions to this topic.
See \emph{related work} below for further references.


Building upon the existing literature for the Euclidean case, we generalize the worst-case iteration complexity analysis of ARC to manifolds, obtaining essentially the same guarantees but with a wider application range: see numerical experiments in Section~\ref{sec:XP} for some examples.

In particular, with the appropriate assumptions discussed in Sections~\ref{sec:firstorderanalysisexp} and~\ref{sec:firstorderanalysisgeneral}, we find that $\varepsilon$-critical points of $f$ on $\calM$ can be computed in $O(1/\varepsilon^{1.5})$ iterations. We also show an iteration complexity bound for the computation of approximate second-order critical points in Section~\ref{sec:secondorder}. Key differences with the Euclidean setting lie in the particular assumptions we make. We further study these assumptions in Sections~\ref{sec:regularity} and~\ref{sec:Dretr}. A subproblem solver---necessary to run ARC---is detailed in Section~\ref{sec:subproblem}. Our algorithm is implemented in the Manopt framework~\citep{manopt} and distributed as part of that toolbox. In Section~\ref{sec:XP}, we close with numerical comparisons to existing solvers, in particular the related Riemannian trust-region method (RTR)~\citep{genrtr}.

\subsection*{Main results}

An important ingredient of ARC on manifolds (Algorithm~\ref{algo:ARC}) is the \emph{retraction} $\Retr$, which allows one to move around the manifold by following tangent vectors. This notion is defined in Section~\ref{sec:arc}. Our results depend on the choice of retraction.

For a twice continuously differentiable cost function $f \colon \calM \to \reals$, the first- and second-order necessary optimality conditions at $x$ read~\citep{yang2012optimality}:
\begin{align*}
	\|\grad f(x)\|_x & = 0, & & \lambdamin(\Hess f(x)) \geq 0,
\end{align*}
where $\grad f$ and $\Hess f$ are the Riemannian gradient and Hessian of $f$---see Section~\ref{sec:firstorderanalysisexp} for definitions; $\|\cdot\|_x$ is the Riemannian norm at $x$ and $\lambdamin$ extracts the smallest eigenvalue of a symmetric operator.

Our first main result applies to \emph{complete} Riemannian manifolds, for which we can use the so-called \emph{exponential map} as retraction $\Retr$. The statement below summarizes more explicit results of Sections~\ref{sec:firstorderanalysisexp} and~\ref{sec:secondorder}, stating iterates of ARC eventually satisfy the necessary optimality conditions up to some tolerance, with a bound on the number of iterations this may require.
\begin{theorem} \label{thm:informalEXP}
%	Let $f$ be a cost function on a complete Riemannian manifold $\calM$. If
	Consider a cost function $f$ on a complete Riemannian manifold $\calM$. If
	\begin{enumerate}
		\item[a)] $f$ is lower bounded (\aref{assu:lowerbound}), and
		\item[b)] the Riemannian Hessian of $f$ is Lipschitz continuous (\aref{assu:firstorderregularityscalar}, \aref{assu:firstorderregularityvector}),
	\end{enumerate}
	then, for any $x_0\in\calM$ and $\varepsilon > 0$, Algorithm~\ref{algo:ARC} with the exponential retraction produces an iterate $x_k \in \calM$ such that $f(x_k) \leq f(x_0)$, $\|\grad f(x_k)\|_{x_k} \leq \varepsilon$ and (if condition~\eqref{eq:secondorderprogress} is enforced) $\lambdamin(\Hess f(x_k)) \geq -\sqrt{\varepsilon}$, with $k = \tilde O(1/\varepsilon^{1.5})$. %, where
	%$\grad f$ is the Riemannian gradient of $f$, $\Hess f$ is its Riemannian Hessian, and
%	$\|\cdot\|$ is the Riemannian norm.
	The bound is dimension- and curvature-free.
\end{theorem}


Our second main result is an extension of the above which allows us to use other retractions, the main motivation being that the exponential map may be unavailable or expensive to compute. We state it as a summary of results in Sections~\ref{sec:firstorderanalysisgeneral} and~\ref{sec:secondorder}.
\begin{theorem} \label{thm:informalRETR}
	Consider a cost function $f$ on a Riemannian manifold $\calM$ equipped with a retraction $\Retr$. If
	\begin{enumerate}
		\item[a)] $f$ is lower bounded (\aref{assu:lowerbound}),
		\item[b)] the pullbacks $f \circ \Retr_x$ satisfy a type of second-order Lipschitz condition (\aref{assu:firstorderregularityscalar}, \aref{assu:firstorderregularityvectorretraction}), and
		\item[c)] the differential of the retraction is well behaved (\aref{assu:DRetr}),
	\end{enumerate}
	then, for any $x_0\in\calM$ and sufficiently small $\varepsilon > 0$, Algorithm~\ref{algo:ARC} with retraction $\Retr$ produces an iterate $x_k \in \calM$ such that $f(x_k) \leq f(x_0)$, $\|\grad f(x_k)\|_{x_k} \leq \varepsilon$ and (if condition~\eqref{eq:secondorderprogress} is enforced and $\Retr$ is second order) $\lambdamin(\Hess f(x_k)) \geq -\sqrt{\varepsilon}$, with $k = \tilde O(1/\varepsilon^{1.5})$. %, where
	%$\grad f$ is the Riemannian gradient of $f$, $\Hess f$ is its Riemannian Hessian, and
	%	$\|\cdot\|$ is the Riemannian norm.
\end{theorem}
We further provide sufficient conditions for the assumptions on the pullbacks and the retraction to be satisfied, in Sections~\ref{sec:regularity} and~\ref{sec:Dretr} respectively. For example, \aref{assu:DRetr} is satisfied if the sublevel set of $x_0$ is compact.


\subsection*{Related work}


Numerous algorithms for unconstrained optimization have been generalized to Riemannian manifolds~\citep{luenberger1972gradient,Gabay1982,smith1994optimization,edelman1998geometry,AMS08}, among them gradient descent, nonlinear conjugate gradients, stochastic gradients \citep{bonnabel2013stochastic,zhang2016riemannian}, BFGS~\citep{ring2012optimization}, Newton's method \citep{adler2002spine} and trust-regions~\citep{genrtr}. See these references and also our numerical experiments in Section~\ref{sec:XP} for a discussion of numerous applications.

ARC in particular was extended to manifolds in the PhD thesis of~\citet{qi2011thesis}. There, under a different set of regularity assumptions, asymptotic convergence analyses are proposed, in the same spirit as the analyses presented in the aforementioned references for other methods. Qi also presents local convergence analyses, showing superlinear local convergence under some assumptions.

In contrast, we here favor a global convergence analysis with explicit bounds on iteration complexity to reach approximate criticality.
Such bounds are standard in optimization on Euclidean spaces. Around the same time, they have been generalized to Riemannian gradient descent and other algorithms by \citet{zhang2016complexitygeodesicallyconvex} (focusing on geodesic convexity), by \citet{bento2017iterationcomplexity} (looking also at proximal point methods), and by \citet{boumal2016globalrates} (also analyzing RTR). In the first two works, the regularity assumptions on the cost function are close in spirit to those we lay out in Section~\ref{sec:firstorderanalysisexp}, whereas in the third work the assumptions are closer to our Section~\ref{sec:firstorderanalysisgeneral}.

Closest to our work, \citet{zhang2018cubicregmanifold} recently proposed a convergence analysis of a cubically regularized method on manifolds, also establishing an $O(1/\varepsilon^{1.5})$ iteration complexity. Their analysis (independent from ours: early versions of our results appeared on public repositories around the same time, theirs two weeks before ours) focuses on compact submanifolds of a Euclidean space and uses a fixed regularization parameter (which must be set properly by the user). Subproblems are assumed to be solved to global optimality, though it appears this could be relaxed within their framework. We improve on these points as follows: our analysis is intrinsic (no embedding space is ever referenced), we do not need $\calM$ to be compact, our regularization parameter $\varsigma_k$ is dynamically adapted (which is both easier for the user and more efficient), and the subproblem solver only needs to meet weak requirements to reach approximate criticality. These improvements lead to implementable, competitive algorithms. \citet{zhang2018cubicregmanifold} also study superlinear local convergence rates, in line with~\citet{qi2011thesis} but with different assumptions.

The work by \citet{zhang2018cubicregmanifold} is also related to adaptive \emph{quadratic} regularization on embedded submanifolds of Euclidean space recently studied by~\citet{hu2017adaptive}, where the quadratic model is written in terms of the Euclidean gradient and Hessian.

More recently, two independent papers generalize work by \citet{jin2019escape} to provide iteration complexity bounds for a Riemannian version of perturbed gradient descent, allowing to reach approximate second-order criticality without looking at the Hessian, and with logarithmic dependence in the dimension of the manifold. \citet{sun2019prgd} provide an analysis based on regularity assumptions akin to the ones we lay out in Section~\ref{sec:firstorderanalysisexp}, while \citet{criscitiello2019escapingsaddles} base their analysis on regularity assumptions closer to the ones we lay out in Section~\ref{sec:firstorderanalysisgeneral}.

Our complexity analysis builds on prior work for the Euclidean case by~\citet{cartis2011adaptivecubic} and \citet{bcgmt}. Complexity lower bounds given by \citet{cartis2018worst} and \citet{carmon2017lower} show that the bounds in \citep{nesterov2006cubic,cartis2011adaptivecubic,bcgmt} are optimal in $\varepsilon$-dependency for the appropriate class of functions. A variant of ARC that is closely related to trust-region methods was presented in \citep{dussault2018arcq}.

Recently, various works have focused on efficiently solving the ARC subproblem (that is, minimizing $m_k$ as defined in~\eqref{eq:mk}) in the Euclidean setting. \citet{agarwal2017finding} propose an efficient method to solve the subproblem leading to fast algorithms for converging to second-order local minima in the Euclidean setting. \citet{carmon2016gradient} and \citet{tripuraneni2017stochastic} propose gradient descent--based methods to solve the subproblem.
Several recent papers consider the effect of subsampling on the subproblem~\citep{tripuraneni2017stochastic,kohler2017sub,zhou2018stochastic,zhang2018adaptive,wang2018stochastic}.

In the Riemannian case, the subproblem is posed on a tangent space, which is a linear subspace. Hence, all of the above methods are applicable in the Riemannian setting as well. In particular, we use the Krylov subspace method originally proposed in~\citep{cartis2011adaptivecubic}. Recently, \citet{carmon2018analysis} and~\citet{gould2019regularizedquadratic} provided a bound on the amount of work this method may require to provide sufficient progress (see also Remark~\ref{rem:carmonduchisubproblem}).
%, based on the target accuracy (independent of the underlying dimension). They hint towards using the method in the ARC framework to provide an overall dimension-free running time bound in terms of gradient and Hessian-vector computations to reach approximate critical points within some tolerance. We expect that, building upon the framework we develop in this paper,
%Such a bound may apply to the Riemannian setting as well.
%see also Remark~\ref{rem:carmonduchisubproblem} below.
%\TODO{See how to reduce space devoted to this.}


%\TODO{Cite~\cite{qi2011thesis} in related work.}

%\TODO{For NB: Re-consider~\citep{ring2012optimization} Lemma 6 and Remark 3 --- there's a lot in common with the discussions we have about how we can disentangle conditions on $f$ and $\Retr$, and second-order conditions on $\Retr$. They also talk about what one can expect from the exponential map.}

On a technical note, in Definition~\ref{def:retrscndnice} we formulate second-order assumptions on the retraction to disentangle the requirements on $f$ from those on the retraction. These are related to (but differ from) the assumptions and discussions in~\citep{ring2012optimization}, specifically Lemma 6, Propositions 5 and 7, and Remarks 2 and 3 in that reference.

%\TODO{About hu2017adaptive: see their eq. (3.3).}


\section{ARC on manifolds} \label{sec:arc}


\begin{algorithm}[t]
	\caption{Riemannian adaptive regularization with cubics (ARC)}
	\label{algo:ARC}
	\begin{algorithmic}[1]
		\State \textbf{Parameters: } $\theta > 0$, $\varsigma_{\min} > 0$, $0 < \eta_1 \leq \eta_2 < 1$, $0 < \gamma_1 < 1 < \gamma_2 < \gamma_3$ % and $\varepsilon > 0$
		\State \textbf{Input:} $x_0 \in \calM$, $\varsigma_0 \geq \varsigma_{\min}$
%		\State \textbf{Init:} $k \leftarrow 0$
%		\Statex 
%		\While{\textbf{true}}
		\For{$k = 0, 1, 2\ldots$}
%		\Statex \TODO{THINK ABOUT WHAT HAPPENS IF GRADIENT IS ZERO, for $x_0$ or for any other iterate: Seems to me if gradient is ever zero, then A3 is impossible to satisfy. Formally, I suppose we could add in A3 that, if this ever happens, then $s_k = 0$ and that we define $\rho_k = 1$ (to disambiguate the 0/0). This does come with some questions related to: do we really consider those steps as successful? Or as unsuccessful? Does that mess up the way we count things? Also, when we use A3 in proofs, we don't want to have to branch between two cases. Perhaps it's better to check, in the algorithm, if gradient ever becomes 0, and terminate in that event. Then, need to make sure all proofs are fine with that.}
%		\State \TODO{If $\grad f(x_k) = 0$, terminate.} \TODO{New; need to change this in second-order analysis too.}
		\State Consider the pullback $\hat f_k = f \circ \Retr_{x_k} \colon \T_{x_k}\calM \to \reals$. Define the model $m_k$ on $\T_{x_k}\calM$: % on the tangent space at $x_k$:
		\begin{align}
			m_k(s) & = \hat f_k(0) + \innersmall{s}{\nabla \hat f_k(0)} + \frac{1}{2}\innersmall{s}{\nabla^2 \hat f_k(0)[s]} + \frac{\varsigma_k}{3} \|s\|^3.
			\label{eq:mk}
		\end{align}
		\State Compute a step $s_k \in \T_{x_k}\calM$ satisfying first-order progress conditions (see Section~\ref{sec:subproblem}):
		\begin{align}
			m_k(s_k) & \leq m_k(0), & \textrm{ and } & & \|\nabla m_k(s_k)\| & \leq \theta \|s_k\|^2.
			\label{eq:firstorderprogress}
		\end{align}
		\Statex \hspace{4.5mm} Optionally, if second-order criticality is targeted, $s_k$ must also satisfy this condition:
		\begin{align}
			\lambdamin(\nabla^2 m_k(s_k)) \geq -\theta \|s_k\|,
			\label{eq:secondorderprogress}
		\end{align}
		\Statex \hspace{4.5mm} where $\lambdamin$ extracts the smallest eigenvalue of a symmetric operator.
%		\Statex \hspace{4.5mm} (This is satisfied in particular if the model Hessian is positive semidefinite at $s_k$.) \TODO{Do we need to define $\lambdamin$?}
		
		\State If $s_k = 0$, terminate (see Lemma~\ref{lem:termination}).
		%\TODO{Discussion about this termination condition. It's related to how we prove the theorems. We could remove this, never terminate, and set in the next step that $\rho_k = 1$ by convention when $s_k = 0$. This creates an infinite stagnating sequence, with $\varsigma_k$ going to $\varsigma_{\min}$ (or staying constant) but it's irrelevant in any case. Could then simplify all talk of 'returning' and finite sequences etc, and move all to one remark about 'all that'.}
		\State Compute the regularized ratio of actual improvement over model improvement:
		\begin{align}
			\rho_k & = \frac{f(x_k) - f(\Retr_{x_k}(s_k))}{m_k(0) - m_k(s_k) + \frac{\varsigma_k}{3} \|s_k\|^3}.
			\label{eq:rhok}
		\end{align}
		\State If $\rho_k \geq \eta_1$, accept the step: $x_{k+1} = \Retr_{x_k}(s_k)$. Otherwise, reject it: $x_{k+1} = x_k$.
%		\Statex
		\State Update the regularization parameter:
		\begin{align}
			\varsigma_{k+1} \in
			\begin{cases}
				\begin{aligned}
					& [\max(\varsigma_{\min}, \gamma_1 \varsigma_k), \varsigma_k] & & \textrm{ if } \rho_k \geq \eta_2 & & \textrm{ (very successful),} \\
					& [\varsigma_k, \gamma_2 \varsigma_k] & & \textrm{ if } \rho_k \in [\eta_1, \eta_2) & & \textrm{ (successful),} \\
					& [\gamma_2 \varsigma_k, \gamma_3 \varsigma_k] & & \textrm{ if } \rho_k < \eta_1 & & \textrm{ (unsuccessful).}
				\end{aligned}
			\end{cases}
			\label{eq:varsigmaupdate}
		\end{align}
%		\Statex 
%		\State $k \leftarrow k+1$
		\EndFor
%		\State \textbf{Output:} $x_k \in \calM$ such that $f(x_k) \leq f(x_0)$ and $\|\grad f(x_k)\| \leq \varepsilon$ (under assumptions.)
	\end{algorithmic}
\end{algorithm}

ARC on manifolds is listed as Algorithm~\ref{algo:ARC}. It is a direct adaptation from~\citep{cartis2011adaptivecubic,bcgmt}. Like many other optimization algorithms, its generalization to manifolds relies on a chosen {retraction}~\citep{shub1986some,AMS08}. For some $x\in\calM$, let $\T_x\calM$ denote the tangent space at $x$: this is a linear space. Intuitively, a retraction $\Retr$ on a manifold provides a means to move away from $x$ along a tangent direction $s\in\T_x\calM$ while remaining on the manifold, producing $\Retr_x(s) \in \calM$.
For a formal definition, we use the \emph{tangent bundle},
\begin{align*}
	\T\calM & = \{ (x, s) : x \in \calM \textrm{ and } s \in \T_x\calM \},
\end{align*}
which is itself a smooth manifold.
\begin{definition}[{\protect{Retraction~\citep[Def.~4.1.1]{AMS08}}}] \label{def:retraction}
	A \emph{retraction} on a manifold $\calM$ is a smooth mapping $\Retr$ from the tangent bundle $\T\calM$ to $\calM$ with the following properties. Let $\Retr_x \colon \T_x\calM \to \calM$ denote the restriction of $\Retr$ to $\T_x\calM$ through $\Retr_x(s) = \Retr(x, s)$. Then,
	\begin{enumerate}
		\item[(i)] $\Retr_x(0) = x$, where $0$ is the zero vector in $\T_x\calM$; and
		\item[(ii)] The differential of $\Retr_x$ at $0$, $\D\Retr_x(0)$, is the identity map on $\T_x\calM$.
	\end{enumerate}
\end{definition}
%Around $t=0$, retraction curves $t \mapsto \Retr_x(ts)$ agree up to first order with geodesics passing through $x$ with velocity $s$.
In other words: retraction curves $c(t) = \Retr_x(ts)$ are smooth and pass through $c(0) = x$ with velocity $c'(0) = \D\Retr_x(0)[s] = s$.
For the special case where $\calM$ is a linear space, the canonical retraction is $\Retr_x(s) = x+s$. For the unit sphere, a typical retraction is $\Retr_x(s) = \frac{x+s}{\|x+s\|}$.

Importantly, the retraction $\Retr$ chosen to optimize over a particular manifold $\calM$ is part of the algorithm specification. For a given cost function $f$ and a specified retraction $\Retr$, at iterate $x_k$, we define the \emph{pullback} of the cost function to the tangent space $\T_{x_k}\calM$:
\begin{align}
	\hat f_k & = f \circ \Retr_{x_k} \colon \T_{x_k}\calM \to \reals.
	\label{eq:fhatk}
\end{align}
This operation lifts $f$ to a linear space.
We then define a model $m_k \colon \T_{x_k}\calM \to \reals$, obtained as a truncated second-order Taylor expansion of the pullback with cubic regularization: see~\eqref{eq:mk}. We use the notation $\inner{\cdot}{\cdot}_x$ to denote the Riemannian metric on $\T_x\calM$, and we usually simplify this notation to $\inner{\cdot}{\cdot}$ when the base point is clear from context. Likewise, $\|s\|_x = \sqrt{\inner{s}{s}_x}$ is the norm of $s \in \T_x\calM$ induced by the Riemannian metric, and we usually omit the subscript, writing $\|s\|$.
%When it is useful to include a reminder, we write $\inner{\cdot}{\cdot}_x$ and $\|\cdot \|_x$ to indicate we work on the tangent space at $x$.
Furthermore, for real functions on linear spaces (such as $\hat f_k$ and $m_k$), we let $\nabla$ and $\nabla^2$ denote the (usual) gradient and Hessian operators.

At iteration $k$, a \emph{subproblem solver} is used to approximately minimize the model $m_k$, producing a \emph{trial step} $s_k$: specific requirements are listed as~\eqref{eq:firstorderprogress} and~\eqref{eq:secondorderprogress}; for the first-order condition, we follow the lead of~\citet{bcgmt}. Section~\ref{sec:subproblem} discusses a practical algorithm.

The quality of the trial step $s_k$ is evaluated by computing $\rho_k$~\eqref{eq:rhok}: the \emph{regularized} ratio of actual to anticipated cost improvement, also following~\citet{bcgmt}. Note that the denominator of $\rho_k$ is equal to the difference between $\hat f_k(0)$ and the second-order Taylor expansion of $\hat f_k$ around 0 evaluated at $s_k$.
If $\rho_k \geq \eta_1$, we accept the trial step and set $x_{k+1} = \Retr_{x_k}(s_k)$: such steps are called \emph{successful}.
Among them, we further identify \emph{very successful} steps, for which $\rho_k \geq \eta_2$; for those, not only is the step accepted, but the regularization parameter $\varsigma_k$ is (usually) decreased.
Otherwise, we reject the step and set $x_{k+1} = x_k$: these steps are \emph{unsuccessful}, and we necessarily increase $\varsigma_k$.

%\TODO{Don't think we need these, although we do need to define the words successful, unsuccessful and very successful iteration.}
%We denote the set of successful iterations and its restriction to the first $\bar k$ iterations by
%\begin{align}
%	\calS & = \{ k : \rho_k \geq \eta_1 \}, & \textrm{ and } & & \calS_{\bar k} & = \calS \cap \{0, \ldots, \bar k-1\}.
%	\label{eq:successful}
%\end{align}
%In contrast, $\calU_{\bar k}$ indexes unsuccessful iterations: it is the complement of $\calS_{\bar k}$ in $0, \ldots, \bar k-1$. Within $\calS$, 
%
%\vspace{2cm}

We expect Algorithm~\ref{algo:ARC} to produce an infinite sequence of iterates. In practice of course, one would terminate the algorithm as soon as approximate criticality is achieved within some prescribed tolerance. This paper bounds the number of iterations this may require. In the unlikely event that the subproblem solver produces the trivial step $s_k = 0$, the algorithm cannot proceed. Fortunately, this only happens if we reached exact criticality of appropriate order. Proofs are in Appendix~\ref{app:arc}.

%Under our requirements for the subproblem solver (conditions \eqref{eq:firstorderprogress} and (optionally) \eqref{eq:secondorderprogress}), Algorithm~\ref{algo:ARC} may terminate at $x_k$ only if it is an exact critical point of appropriate order: see Lemma~\ref{lem:termination}. In that case, $\rho_k$ is not defined, and $k$ is neither successful nor unsuccessful; all results still hold, with limit statements interpreted as referring to $x_k$.
%Barring this unlikely event, the algorithm defines an infinite sequence of iterates. %\TODO{Part of the termination conversation.}
%\TODO{Need to check statements and proofs for the case of finite list of iterates, to make sure they hold, and make a comment about it here. In particular, need to say something about the limit statements.}
%In order to do so, it is convenient to allow the infinite sequence to be defined: this is why we do not include an explicit stopping criterion in Algorithm~\ref{algo:ARC}. \TODO{rework this; it's confusing because there is a stopping criterion; just, not what one would expect, and it's unlikely ever to trigger.} \TODO{Consolidate with lemmas below. Perhaps before getting into all of this, re-evaluate what it would take to make the algorithm stop at approximate stationary points? Might simplify the story.} \TODO{Seems like a good idea to rephrase in terms of: the subproblem solver is \emph{allowed} to return $s_k = 0$ iff ... Barring these exceptional circumstances, the algorithms runs ad infinitum, generating an infinite sequence we set out to study.}
%
\begin{lemma}\label{lem:termination}
	If second-order progress~\eqref{eq:secondorderprogress} is not enforced, then the first-order condition~\eqref{eq:firstorderprogress} allows the subproblem solver to return $s_k = 0$ if and only if $\grad f(x_k) = 0$.
	If both~\eqref{eq:firstorderprogress} and~\eqref{eq:secondorderprogress} are enforced, the subproblem solver is allowed to return $s_k = 0$ if and only if $\grad f(x_k) = 0$ and $\Hess f(x_k)$ is positive semidefinite.
\end{lemma}

%\vspace{2cm}
%
%
%%We sometimes tacitly assume the algorithm produces an infinite sequence of iterates (as is typical); all arguments are easily adapted to the finite case.
%
%\TODO{SAY SOMETHING ABOUT SUBPROBLEM SOLVER}
%The subproblem is posed in a linear space (the tangent space at $x_k$). Thus, any standard technique from the Euclidean case carries over in principle. We give a brief discussion in Section~\ref{sec:subproblem}.
%The subproblem solver we describe in Section~\ref{sec:subproblem} necessarily succeeds in producing a step $s_k$ which satisfies both~\aref{assu:firstorderprogress} and~\aref{assu:secondorderprogress} in a finite number of operations.
%Indeed, since the subproblem consists in minimizing $m_k$ without constraints on a linear space, any global optimum $s$ satisfies both $\nabla m_k(s) = 0$ and $\lambda_{\min}(\nabla^2 m_k(s)) \geq 0$. Since the described subproblem solver produces a global optimum to the subproblem in a finite number of operations (in exact arithmetic), it a fortiori produces a step $s_k$ which satisfies the assumptions. Notwithstanding, we stress here that it is not necessary to solve the subproblem to global optimality. %: in practice, we expect to produce cheaper acceptable steps.
%\TODO{END SUBPROBLEM SOLVER BIT}


%\TODO{Ordering issues here?}

We now introduce two basic assumptions about the cost function $f$, affording us two supporting lemmas.
The first common assumption is that the cost function $f$ is lower bounded. %This is necessary to ensure existence of points on $\calM$ with arbitrarily small gradient (though it does not guarantee existence of a point where the gradient is exactly zero.)
\begin{assumption} \label{assu:lowerbound}
	There exists a finite $\flow$ such that $f(x) \geq \flow$ for all $x \in \calM$.
\end{assumption}
The second assumption is that $f$ is sufficiently differentiable so that the models $m_k$ are well defined, and that second-order Taylor expansions of $\hat f_k$ in the tangent space at $x_k$ are sufficiently accurate. In the Euclidean case, the latter follows from a Lipschitz condition on the Hessian of $f$. In the next two sections, we discuss how this generalizes to manifolds.
\begin{assumption} \label{assu:firstorderregularityscalar}
	The cost function $f$ is twice continuously differentiable.
	% We might as well include the "continuously", because the following conditions actually require it. Indeed, if the Hessian was discontinuous at some point, then we could violate the conditions below even on R^n.
	Furthermore, there exists a constant $L$ such that, at each iteration $k$, for the trial step $s_k$ selected by the subproblem solver, the pullback $\hat f_k = f \circ \Retr_{x_k}$ satisfies
	\begin{align}
		\hat f_k(s_k) - \left[ \hat f_k(0) + \innersmall{s_k}{\nabla \hat f_k(0)} + \frac{1}{2}\innersmall{s_k}{\nabla^2 \hat f_k(0)[s_k]} \right] & \leq \frac{L}{6} \|s_k\|^3.
	\end{align}
%	\TODO{Call this one $L$ and call the other one $L'$ so we can track where they end up. Basically, $L$ comes up only in $\varsigma_{\max}$, ``linearly''. The latter comes up in a log-term, but it ALSO comes up in $L'/2 + \theta + \varsigma_{\max}$, so, it matters.}
\end{assumption}
The two supporting lemmas below follow the standard Euclidean analysis. The first lemma establishes that the regularization parameter $\varsigma_k$ does not grow unbounded.
\begin{lemma}[{\protect\citet[Lem.~2.2]{bcgmt}}] \label{lem:varsigmamax}
	Under~\aref{assu:firstorderregularityscalar}, the regularization parameter remains bounded: for all $k$, it holds that $\varsigma_k \leq \varsigma_{\max}$, with
	\begin{align}
	%\forall k, \quad \varsigma_k \leq
	\varsigma_{\max} = \max\left( \varsigma_0, \frac{L\gamma_3}{2(1-\eta_2)} \right).
	\label{eq:varsigmamax}
	\end{align}
\end{lemma}
Conditioned on the conclusions of this lemma, the next lemma states that among the first $\bar k$ iterations of ARC, a certain number are sure to be successful.
\begin{lemma}[{\protect\citet[Thm.~2.1]{cartis2011adaptivecubic}}] \label{lem:boundKsuccessfulsteps}
	If $\varsigma_k \leq \varsigma_{\max}$ for all $k$ (as provided by Lemma~\ref{lem:varsigmamax}), then the number $K$ of successful iterations among $0, \ldots, \bar k-1$ satisfies
	\begin{align*}
		\bar k \leq \left(1 + \frac{|\log(\gamma_1)|}{\log(\gamma_2)}\right) K + \frac{1}{\log(\gamma_2)} \log\left(\frac{\varsigma_{\max}}{\varsigma_0}\right).
	\end{align*}
%	\TODO{Rephrase so that it says: for any integer $K \geq 0$, after $...$ iterations, at least $K$ are successful (to rework). Might make it easier to summon it with $K = K_1(\varepsilon)$ for example.}
\end{lemma}
In other words, in order to bound the \emph{total} number of iterations ARC may require to attain a certain goal, it is sufficient to bound the number of \emph{successful} iterations that goal may require. The following proposition (extracted from the main proof in~\citep{bcgmt}) further states that this can be done by showing successful steps are not too short.
\begin{proposition} \label{prop:sumnormskcube}
	Let $\{ (x_0, s_0), (x_1, s_1), \ldots \}$ be the set of iterates and trial steps generated by Algorithm~\ref{algo:ARC}. If~\aref{assu:lowerbound} holds, we have
	\begin{align*}
	\sum_{k\in\calS} \|s_k\|^3 \leq \frac{3(f(x_0) - \flow)}{\eta_1 \varsigma_{\min}},
	\end{align*}
	where $\calS$ is the set of successful iterations.
\end{proposition}
\begin{proof}
	By definition, if iteration $k$ is successful, then $\rho_k \geq \eta_1$~\eqref{eq:rhok}. Combining with the first part of the first-order progress condition~\eqref{eq:firstorderprogress} yields
	\begin{align*}
	f(x_k) - f(x_{k+1}) \geq \eta_1\!\left(m_k(0) - m_k(s_k) + \frac{\varsigma_k}{3} \|s_k\|^3 \right)  \geq \frac{\eta_1\varsigma_{\min}}{3} \|s_k\|^3. % \geq \frac{\eta_1\varsigma_k}{3} \|s_k\|^3
	\end{align*}
	On the other hand, for unsuccessful iterations, $x_{k+1} = x_k$ and the cost does not change.
	Using~\aref{assu:lowerbound}, a telescoping sum yields:
	\begin{align*}
	f(x_0) - \flow & \geq \sum_{k=0}^\infty f(x_k) - f(x_{k+1}) = \sum_{k\in\calS} f(x_k) - f(x_{k+1}) \geq \frac{\eta_1 \varsigma_{\min}}{3} \sum_{k\in\calS} \|s_k\|^3,
	\end{align*}
	as announced.
\end{proof}
%If successful steps have norm at least $\alpha$, then the left-hand side in this proposition is at least $\alpha^3 K$, where $K$ ... yeah, it's messy to start explaining this, because it'll be the successful steps among a subset. Drop it.


\section{First-order analysis with the exponential map}
\label{sec:firstorderanalysisexp}


In this section, we provide a first-order analysis of Algorithm~\ref{algo:ARC} for the case where $\calM$ is a complete manifold and we use the exponential retraction $\Retr = \Exp$---we define these terms momentarily. This notably encompasses the Euclidean case where $\calM = \Rn$, with $\Exp_x(s) = x+s$, as well as all compact or Hadamard manifolds. As such, the results in this section offer a strict generalization of the Euclidean analysis proposed in~\citep{bcgmt} under the assumption of Lipschitz continuous Hessian. An in-depth reference for the Riemannian geometry tools we use is the monograph by~\citet{lee2018riemannian}, while \citet{AMS08} offer an optimization-focused treatment.

On a complete Riemannian manifold $\calM$, for any point $x$ and tangent vector $v \in \T_x\calM$, there exists a unique smooth curve $\gamma_v \colon \reals \to \calM$ such that $\gamma_v(0) = x$, $\gamma_v'(0) = v$ and, for $t < t'$ close enough, $\gamma_v|_{[t, t']}$ is the shortest path connecting $\gamma_v(t)$ to $\gamma_v(t')$. This curve is called a \emph{geodesic}. The \emph{exponential map} is built from these geodesics as the map
\begin{align*}
	\Exp \colon \T\calM \to \calM \colon (x, v) \mapsto \Exp_x(v) = \gamma_v(1).
\end{align*}
This is a smooth map. Because $\Exp_x(tv) = \gamma_{tv}(1) = \gamma_v(t)$, we also find that $\Exp_x(0) = x$ and $\D\Exp_x(0)[v] = v$, so that the exponential map is indeed a retraction (Definition~\ref{def:retraction}).
If the manifold is not complete, then $\Exp$ is only defined on an open subset of $\T\calM$: when we need $\calM$ to be complete, we say so explicitly.
%: to handle this scenario, one would need to add safeguards to Algorithm~\ref{algo:ARC}.)

The Riemannian gradient of $f \colon \calM \to \reals$, denoted by $\grad f$, is the vector field on $\calM$ such that $\D f(x)[s] = \inner{\grad f(x)}{s}$, where $\D f(x)[s]$ is the directional derivative of $f$ at $x$ along the tangent direction $s$. One can show that
\begin{align}
	\grad f(x) & = \nabla(f \circ \Exp_x)(0) = \nabla \hat f_x(0),
	\label{eq:gradfpullback}
\end{align}
so that the Riemannian gradient of $f$ at $x$ is nothing but the Euclidean gradient of the pullback $\hat f_x = f \circ \Exp_x$ at the origin of the tangent space $\T_x\calM$ (see also Lemma~\ref{lem:derivativespullback} for a similar statement with retractions).

The Riemannian Hessian of $f$ is the covariant derivative of the gradient vector field, with respect to the Riemannian connection. Denoted by $\Hess f$, it defines a tensor field as follows: $\Hess f(x)$ is a linear operator from $\T_x\calM$ into itself, self-adjoint with respect to the Riemannian metric on that tangent space. Analogously to~\eqref{eq:gradfpullback}, one can show that
\begin{align}
	\Hess f(x) & = \nabla^2 \hat f_x(0),
	\label{eq:HessfHessfhatExp}
\end{align}
which expresses the Riemannian Hessian of $f$ at $x$ as the Euclidean Hessian of the pullback $\hat f_x$ at the origin of $\T_x\calM$. (Here too, see Lemma~\ref{lem:derivativespullback} below.)

These two statements show that the model $m_k$~\eqref{eq:mk} can be written equivalently as
\begin{align}
	m_k(s) & = f(x_k) + \inner{\grad f(x_k)}{s} + \frac{1}{2} \inner{\Hess f(x_k)[s]}{s} + \frac{\varsigma_k}{3} \|s\|^3
	\label{eq:mkExp}
\end{align}
and that~\aref{assu:firstorderregularityscalar} requires
\begin{align}
	f(\Exp_{x_k}(s_k)) - \left[ f(x_k) + \innersmall{s_k}{\grad f(x_k)} + \frac{1}{2}\innersmall{s_k}{\Hess f(x_k)[s_k]} \right] & \leq \frac{L}{6} \|s_k\|^3
	\label{eq:A2storyExp}
\end{align}
for each $(x_k, s_k)$ produced by Algorithm~\ref{algo:ARC}.

In particular, if $\calM$ is a Euclidean space with the exponential map $\Exp_x(s) = x+s$, it is well known that we can secure~\eqref{eq:A2storyExp} if we assume that the Hessian of $f$ is $L$-Lipschitz continuous. This can be written as
\begin{align*}
	\opnorm{\nabla^2 f(x) - \nabla^2 f(y)} & \leq L \|x-y\|,
\end{align*}
where the norm on the left-hand side is the operator norm.
Generalizing this to the Riemannian setting, we face the issue that $\Hess f(x)$ and $\Hess f(y)$ are linear operators defined on distinct tangent spaces (if $x \neq y$): in order to compare them, we need one more tool to compare tangent vectors in distinct tangent spaces.

Given a smooth curve $c \colon [0, 1] \to \calM$ connecting $c(0) = x$ to $c(1) = y$, consider a tangent vector $v \in \T_x\calM$ and a smooth vector field $Z \colon [0, 1] \to \T\calM$ along $c$---that is, $Z(t) \in \T_{c(t)}\calM$---such that $Z(0) = v$.
If the covariant derivative of $Z$ with respect to the Riemannian connection vanishes identically, then we say that $Z$ is a \emph{parallel vector field} along $c$. (For example, the velocity vector field $\gamma'$ of a geodesic $\gamma$ is parallel.) This vector field exists and is unique. We call $Z(1)$ the \emph{parallel transport} of $v$ from $x$ to $y$ along $c$. Parallel transports are linear isometries with respect to the Riemannian metric, and they depend on the chosen path.

Using parallel transports, we can formulate a standard notion of Lipschitz continuity for Riemannian Hessians.
We use the following notation often: given a tangent vector $s \in \T_x\calM$, 
\begin{align}
	P_s \colon \T_x\calM \to \T_{\Exp_x(s)}\calM
	\label{eq:Ps}
\end{align}
denotes parallel transport along the geodesic $\gamma(t) = \Exp_x(ts)$ from $t = 0$ to $t = 1$.
\begin{definition} \label{def:LipschitzHessRiemann}
	A function $f \colon \calM \to \reals$ on a Riemannian manifold $\calM$ has an \emph{$L$-Lipschitz continuous Hessian} if it is twice differentiable and if, for all $(x, s)$ in the domain of $\Exp$,
	\begin{align*}
		%\forall (x, s) \in \T\calM, &&
		\opnorm{ P_s^{-1} \circ \Hess f(\Exp_x(s)) \circ P_s - \Hess f(x)} & \leq L \|s\|. %,
	\end{align*}
%	where $P_s \colon \T_x\calM \to \T_{\Exp_x(s)}\calM$ denotes parallel transport along $\gamma(t) = \Exp_x(ts)$ from $\gamma(0)$ to $\gamma(1)$.
\end{definition}
%\TODO{need complete for this? no: no need; see my book}
For $f$ three times continuously differentiable, this property holds if and only if the covariant derivative of the Riemannian Hessian is uniformly bounded by $L$ (we omit a proof). In particular, this holds with some $L$ for any smooth function on a compact manifold.
In the Euclidean case, parallel transports are identity maps (independent of the transport curve), so that this is equivalent to the usual definition.

Crucially, for cost functions with Lipschitz Hessian, we recover familiar-looking bounds on Taylor expansions of both $f$ itself and, as will be instrumental momentarily, of $\grad f$. Results of this nature are standard: they appear frequently in complexity analyses for Riemannian optimization, see for example~\citep{ferreira2002kantorovichnewton,bento2017iterationcomplexity,sun2019prgd}. Proofs are in Appendix~\ref{app:firstorderexp}.
\begin{proposition} \label{prop:LipschitzHessianBounds}
	Let $f \colon \calM \to \reals$ be twice differentiable on a Riemannian manifold $\calM$. Given $(x, s)$ in the domain of $\Exp$, assume there exists $L \geq 0$ such that, for all $t \in [0, 1]$,
	\begin{align*}
		\opnorm{P_{ts}^{-1} \circ \Hess f(\Exp_x(ts)) \circ P_{ts} - \Hess f(x)} & \leq L \|ts\|. %,
	\end{align*}
	%where $P_{ts}$ denotes parallel transport along $\gamma(t) = \Exp_x(ts)$ from $\gamma(0)$ to $\gamma(t)$.
	Then, the two following inequalities hold:
	\begin{align*}
		\left| f(\Exp_x(s)) - f(x) - \inner{s}{\grad f(x)} - \frac{1}{2} \inner{s}{\Hess f(x)[s]} \right| & \leq \frac{L}{6} \|s\|^3, \textrm{ and }\\
		\left\| P_{s}^{-1} \grad f(\Exp_x(s)) - \grad f(x) - \Hess f(x)[s] \right\| & \leq \frac{L}{2} \|s\|^2.
	\end{align*}
\end{proposition}
This justifies the introduction of the following assumption, still with notation~\eqref{eq:Ps} for $P_{s_k}$.
\begin{assumption} \label{assu:firstorderregularityvector}
	There exists a constant $L'$ such that, at each successful iteration $k$, for the step $s_k$ selected by the subproblem solver, we have
	\begin{align}
		\left\| P_{s_k}^{-1} \grad f(\Exp_{x_k}(s_k)) - \grad f(x_k) - \Hess f(x_k)[s_k] \right\| & \leq \frac{L'}{2} \|s_k\|^2. %,
	\end{align}
%	where $P_{s_k}$ denotes parallel transport along $\gamma(t) = \Exp_{x_k}(ts_k)$ from $\gamma(0)$ to $\gamma(1)$.
\end{assumption}
\begin{corollary}
	If $f \colon \calM \to \reals$ has $L$-Lipschitz continuous Hessian, then, for any sequence $\{(x_k, s_k)\}_{k=0,1,2,\ldots}$ in the domain of $\Exp$, $f$ satisfies~\aref{assu:firstorderregularityscalar} with the same $L$ (and $\Retr = \Exp$), and it satisfies~\aref{assu:firstorderregularityvector} with $L' = L$. In particular, if $f$ is smooth and $\calM$ is compact, assumptions \aref{assu:lowerbound}, \aref{assu:firstorderregularityscalar} and \aref{assu:firstorderregularityvector} hold. %\TODO{Would this also hold for compact sublevel set of $x_0$? --- maybe, but it's not obvious at all, because the claim for compact manifolds is related to the proof-omitted claim that Lipschitz Hessian is equivalent to bounded third derivative, but this may depend on some path, which, if we restrict ourselves to a subset of the manifold, it's not so clear what happens. Drop it.}
\end{corollary}
We can now state our main result regarding the complexity of running Algorithm~\ref{algo:ARC} on a complete manifold with the exponential retraction, for the purpose of computing approximate first-order critical points. %\TODO{This is without~\eqref{eq:secondorderprogress}; is this clear enough?}
We require $\calM$ to be complete so that $\Exp$ is indeed a retraction, defined on the whole tangent bundle.
The proof follows~\citep[Thm.~2.5]{bcgmt} up to the fact that we bound the \emph{total number} of successful iterations that map to points with large gradient (as opposed to bounding the number of such iterations among the first $\bar k$), more in the spirit of~\citep{cartis2012complexity}: this enables us to make a statement about the limit of $\|\grad f(x_k)\|$.
Recall that $\varsigma_{\max}$ is provided by Lemma~\ref{lem:varsigmamax}.
\begin{theorem} \label{thm:masterboundSLipschitzHessian}
	Let $\calM$ be a complete Riemannian manifold and let $\Retr = \Exp$.
	Under~\aref{assu:lowerbound}, \aref{assu:firstorderregularityscalar} and \aref{assu:firstorderregularityvector}, for an arbitrary $x_0 \in \calM$, let $x_0, x_1, x_2\ldots$ be the iterates produced by Algorithm~\ref{algo:ARC}.
	For any $\varepsilon > 0$, the \emph{total number} of \emph{successful} iterations $k$ such that $\|\grad f(x_{k+1})\| > \varepsilon$ is bounded above by
	\begin{align*}
		K_1(\varepsilon) & \triangleq \frac{3(f(x_0) - \flow)}{\eta_1 \varsigma_{\min}} \left(\frac{L'}{2}+\theta+\varsigma_{\max}\right)^{1.5} \frac{1}{\varepsilon^{1.5}}.
	\end{align*}
	Furthermore, $\lim_{k \to \infty} \|\grad f(x_k)\| = 0$ (that is, limit points are critical).
\end{theorem}
\begin{proof}
	If iteration $k$ is successful, we have $x_{k+1} = \Exp_{x_k}(s_k)$.
%	Let $s_k \in \T_{x_k}\calM$ be the trial step generated by the algorithm at \TODO{a successful? if so, remove ``trial''} iteration $k$.
	The gradient of the model $m_k$~\eqref{eq:mkExp} at $s_k$ (in the tangent space at $x_k$) is given by
	\begin{align*}
	\nabla m_k(s_k) & = \grad f(x_k) + \Hess f(x_k)[s_k] + \varsigma_k \|s_k\| s_k \\
	& = P_{s_k}^{-1}\grad f(x_{k+1}) \\ & \quad + \left( \grad f(x_k) + \Hess f(x_k)[s_k] - P_{s_k}^{-1}\grad f(x_{k+1}) \right) + \varsigma_k \|s_k\| s_k,
	\end{align*}
	with $P_{s_k}$ as defined by~\eqref{eq:Ps}. % and $x_k^+ = \Exp_{x_k}(s_k)$.
	Owing to the first-order progress condition~\eqref{eq:firstorderprogress}, by the triangle inequality and also using~\aref{assu:firstorderregularityvector}, we find
	\begin{align*}
		\theta \|s_k\|^2 \geq \|\nabla m_k(s_k)\| \geq \|P_{s_k}^{-1}\grad f(x_{k+1})\| - \frac{L'}{2} \|s_k\|^2 - \varsigma_k \|s_k\|^2.
	\end{align*}
	Rearranging and using that $P_{s_k}$ is an isometry, we get for all successful $k$ that
	\begin{align}
		\|\grad f(x_{k+1})\| = \|P_{s_k}^{-1}\grad f(x_{k+1})\| \leq \left(\frac{L'}{2} + \theta + \varsigma_{\max}\right) \|s_k\|^2,
		\label{eq:ineqgradnormxkplusLipschitzHessian}
	\end{align}
	where we also called upon Lemma~\ref{lem:varsigmamax} to claim $\varsigma_k \leq \varsigma_{\max}$.
	Define a subset of the successful steps based on the tolerance $\varepsilon$:
	\begin{align*}
		\calS_\varepsilon & = \{ k : \rho_k \geq \eta_1 \textrm{ and } \|\grad f(x_{k+1})\| > \varepsilon \}.
	\end{align*}
	For $k \in \calS_\varepsilon$, we can lower-bound $\|s_k\|^3$ using~\eqref{eq:ineqgradnormxkplusLipschitzHessian} since $\|\grad f(x_{k+1})\| > \varepsilon$. Then, calling upon Proposition~\ref{prop:sumnormskcube}, we find
	\begin{align*}
		\frac{3(f(x_0) - \flow)}{\eta_1 \varsigma_{\min}} & \geq \sum_{k \in \calS_\varepsilon} \|s_k\|^3 \geq  \frac{\varepsilon^{1.5}}{\left(\frac{L'}{2} + \theta + \varsigma_{\max}\right)^{1.5}} |\calS_\varepsilon|.
	\end{align*}
	This proves the main claim. The claim regarding limit points is proved in Appendix~\ref{app:firstorderexp}.
\end{proof}
(Above, it is natural to consider the sequence $\{x_{k+1}\}$ for successful iterations $k$, as this enumerates each distinct point in the whole sequence once.) % Notice that we can't loop back to a previously produced point, since the algorithm is monotonically decreasing, so it's ok to talk about distinct points.
%Recall the definition of $\calS_{\bar k}$~\eqref{eq:successful}.
A key consequence of Theorem~\ref{thm:masterboundSLipschitzHessian} is that, if the number of \emph{successful} iterations among $0, \ldots, \bar k-1$ strictly exceeds $K_1(\varepsilon)$, then it must be that $\|\grad f(x_k)\| \leq \varepsilon$ for some $k$ in $0, \ldots, \bar k$. Combining this with Lemma~\ref{lem:boundKsuccessfulsteps} yields the first main result: a bound on the total number of iterations it may take ARC to produce an approximate critical point on a complete manifold, using the exponential map, and (essentially) assuming a Lipschitz continuous Hessian.
%
% The argument goes like this:
%
%
%  (\bar k -> q for ease of notation)
%
%  Let K be the number of successful iterations among 0, ..., q - 1.
%  Lemma~\ref{lem:boundKsuccessfulsteps} states q <= aK + b, that is, K >= (q - b)/a
%
%  Let K_1 = K_1(eps) be the number provided by Theorem~\ref{thm:masterboundSLipschitzHessian}.
% 
%  If K > K_1, then we win, that is, there exists k in 0, ..., q such that grad f(k) is <= eps.
%
%  In particular, if (q - b)/a > K_1, then K >= (q - b)/a > K_1, and we win.
%
%  In other words, if q > aK_1 + b, we win.
%
%  To generate x_0, ..., x_q, we need q iterations. Hence, we need the number of iterations to strictly exceed aK_1 + b.
%
%  Birgin et al. don't have the +1, and they have a "round down", but their definition of S_k is different from ours: the go all the way up to k, not only up to k-1. It's super annoying: just add the +1.
%
%
\begin{corollary} \label{cor:mastercorollary}
	Under the assumptions of Theorem~\ref{thm:masterboundSLipschitzHessian}, Algorithm~\ref{algo:ARC} produces a point $x_k \in \calM$ such that $f(x_k) \leq f(x_0)$ and $\|\grad f(x_k)\| \leq \varepsilon$ in at most
	\begin{align*}
		\left(1 + \frac{|\log(\gamma_1)|}{\log(\gamma_2)}\right) K_1(\varepsilon) + \frac{1}{\log(\gamma_2)} \log\left(\frac{\varsigma_{\max}}{\varsigma_0}\right) + 1
	\end{align*}
	iterations.
\end{corollary}
In the Euclidean case, this recovers the result of~\citep{bcgmt} exactly. Note also that this complexity result is unaffected by the curvature of the manifold. Moreover, if $L$ is known and $L' = L$ (which holds under the Lipschitz Hessian assumption), then we can set $\varsigma_0 = \varsigma_{\min} = \frac{1}{2}L$ (so that $\varsigma_{\max} = \frac{\gamma_3}{2(1-\eta_2)}L$) and $\theta = \frac{1}{2}L$. With those choices, we find that
\begin{align}
	K_1(\varepsilon) & = \frac{6}{\eta_1} \left( 1 + \frac{\gamma_3}{2(1-\eta_2)} \right)^{1.5} (f(x_0) - \flow) \sqrt{L} \frac{1}{\varepsilon^{1.5}}.
\end{align}
This exhibits a complexity scaling with $\sqrt{L}$ when $L$ is known, as in~\citep{nesterov2006cubic} and in the lower bound discussed in~\citep{carmon2017lower}.

%\TODO{have a discussion about what we used that is so specific to Exp, and isometry of PT, and what will happen next. Or defer to next section; it's already long, try to keep it simple.}

\section{First-order analysis with a general retraction}
\label{sec:firstorderanalysisgeneral}

The results of the previous section provide a strict, lossless generalization of a known result in the Euclidean case. However, we note two practical shortcomings:
\begin{itemize}
	\item[--] For the analysis to apply, the algorithm must compute the exponential map. % (and $\calM$ needs to be complete, although... that's a weak point because I don't have good examples of globally defined retractions on incomplete manifolds. It's not even that clear how to optimize over incomplete manifolds to begin with (at least, in theory).)
	\item[--] The Lipschitz condition (Definition~\ref{def:LipschitzHessRiemann}) may be difficult to assess as it involves parallel transports or bounding the covariant derivative of the Riemannian Hessian.
\end{itemize}
Regarding the first point, we organized proofs in Appendix~\ref{app:firstorderexp} to highlight why it is not clear how to generalize Proposition~\ref{prop:LipschitzHessianBounds} to general retractions. In a nutshell, it is because parallel transports and geodesics interact particularly nicely through the fact that the velocity vector field of a geodesic is a parallel vector field.

To address both points, we propose alternate regularity conditions which (a) allow for any retraction, and (b) involve conceptually simpler objects. We do this by focusing on the pullbacks $\hat f_x = f \circ \Retr_x$, which have the merit of being scalar functions on linear spaces---this is in the spirit of prior work~\citep{boumal2016globalrates}. We offer justification for these assumptions below, and in section~\ref{sec:regularity}.

The first regularity assumption, \aref{assu:firstorderregularityscalar}, is readily phrased in terms of the pullback. We focus on providing a replacement for the second condition: \aref{assu:firstorderregularityvector}. Translating this condition to pullbacks by analogy, we aim to bound the difference between $\nabla \hat f_x(s)$---which, conveniently, is a vector tangent at $x$---and a classical truncated Taylor expansion for it: $\nabla \hat f_x(0) + \nabla^2 \hat f_x(0)[s]$. In so doing, it is useful to note that $\nabla \hat f_x(s)$ is related to $\grad f(\Retr_x(s))$ by a linear operator, as follows:
\begin{align}
	\nabla \hat f_x(s) & = T_s^* \grad f(\Retr_x(s)), & \textrm{ with } & & T_s & = \D\Retr_x(s) \colon \T_x\calM \to \T_{\Retr_x(s)}\calM,
	\label{eq:nablahatfTgradf}
\end{align}
where the star indicates the adjoint with respect to the Riemannian metric. Indeed,
\begin{align}
	\forall s, \dot s \in \T_x\calM, \qquad \innersmall{\nabla \hat f_x(s)}{\dot s}_x & = \D \hat f_x(s)[\dot s] \nonumber\\
		& = \D f(\Retr_{x}(s))[ \D\Retr_{x}(s)[\dot s] ] \nonumber\\
		& = \inner{\grad f(\Retr_{x}(s))}{\D\Retr_{x}(s)[\dot s]}_{\Retr_{x}(s)} \nonumber\\
		& = \inner{\left(\D\Retr_{x}(s)\right)^*[\grad f(\Retr_{x}(s))]}{\dot s}_x.
		\label{eq:DRetrnablagrad}
\end{align}
(This also plays a role in~\citep[p599]{ring2012optimization}.)


Considering for a moment how an estimate for $\nabla \hat f_x(s)$ might look like if we use the exponential retraction and under the Lipschitz continuous Hessian assumption~\aref{assu:firstorderregularityvector} as above, we find by triangular inequality that
\begin{align*}
	\|\nabla \hat f_x(s) - \nabla \hat f_x(0) - \nabla^2 \hat f_x(0)[s]\| &
	\leq \|\nabla \hat f_x(s) - P_s^{-1} \grad f(\Exp_x(s))\| \\ & \quad + \|P_s^{-1} \grad f(\Exp_x(s)) - \grad f(x) - \Hess f(x)[s]\| \\
	& \leq \opnorm{T_s^* - P_s^{-1}} \|\grad f(\Exp_x(s))\| + \frac{L'}{2} \|s\|^2,
\end{align*}
using~\eqref{eq:DRetrnablagrad} with $T_s = \D\Exp_x(s)$.
%%%%%%%%%%%%%%%
Since parallel transport $P_s$~\eqref{eq:Ps} is an isometry, $P_s^{-1} = P_s^*$ and $\opnorm{T_s^* - P_s^{-1}} = \opnorm{T_s - P_s}$.
For small $s$, we expect $T_s$ (the differential of the exponential map) and $P_s$ (parallel transport) to be nearly the same. Indeed, $\opnorm{T_s - P_s}$ is a continuous function of $s$ and $T_0 = P_0 = \Id$. How much they differ for nonzero $s$ is related to the curvature of the manifold. As a result, we conclude that
\begin{align}
	\left\| \nabla \hat f_x(s) - \nabla \hat f_x(0) - \nabla^2 \hat f_x(0)[s] \right\| & \leq \frac{L'}{2} \|s\|^2 + q(\|s\|) \cdot \|\grad f(\Exp_x(s))\|
	\label{eq:ineqnablahatfsmotivation}
\end{align}
for some continuous function $q \colon \reals^+ \to \reals^+$ such that $q(0) = 0$. %We expect $q$ to depend on the curvature of the manifold.

As an illustration, consider the special case where $\calM$ has constant sectional curvature $C$. In this case, it can be shown using Jacobi fields (see the proof of Lemma~\ref{lemma:exponentialbound} in the appendix) that
\begin{align}
	T_s \dot s & = \D \Exp_x(s)[\dot s] = P_s \dot s + h(\|s\|) P_s\!\left(\dot s - \frac{\inner{s}{\dot s}}{\|s\|^2} s\right),
	\label{eq:DExpConstantCurvature}
\end{align}
where
\begin{align*}
	h(\|s\|) & = \begin{cases}
			0 & \textrm{ if } C = 0, \\
			\frac{\sin(\|s\|/R)}{\|s\|/R} - 1 & \textrm{ if } C = \frac{1}{R^2} > 0, \\
			\frac{\sinh(\|s\|/R)}{\|s\|/R} - 1 & \textrm{ if } C = -\frac{1}{R^2} < 0.
		\end{cases}
\end{align*}
Thus, for manifolds with constant sectional curvature, inequality~\eqref{eq:ineqnablahatfsmotivation} holds with $q(\|s\|) = |h(\|s\|)|$, independent of $x$. This function behaves as $\frac{1}{6}\!\left(\frac{\|s\|}{R}\right)^2 = \frac{C}{6} \|s\|^2$ for small $\|s\|$. This derivation generalizes to manifolds with sectional curvature bounded both from above and from below~\citep{tripuraneni2018averaging}, \citep[Thm.~A.2.9]{waldmann2012geometric}.

Returning to retractions in general, the above motivates us to introduce the following assumption on the pullbacks, meant to replace \aref{assu:firstorderregularityvector}.
\begin{assumption} \label{assu:firstorderregularityvectorretraction}
	There exists a constant $L'$ such that, at each successful iteration $k$, for the step $s_k$ selected by the subproblem solver, the pullback $\hat f_k = f \circ \Retr_{x_k}$ obeys
	\begin{align}
		\left\| \nabla \hat f_k(s_k) - \nabla \hat f_k(0) - \nabla^2 \hat f_k(0)[s_k] \right\| & \leq \frac{L'}{2} \|s_k\|^2 + q(\|s_k\|) \|\grad f(\Retr_{x_k}(s_k))\|,
	\end{align}
	where $q \colon \reals^+ \to \reals^+$ is some continuous function satisfying $q(0) = 0$.
\end{assumption}
Notice how this assumption involves simple tools compared to~\aref{assu:firstorderregularityvector}, which relies on the exponential map and parallel transports. Furthermore, if we strengthen the condition by forcing $q \equiv 0$, we get a Lipschitz-type condition on the pullback: see Section~\ref{sec:regularity}.

%%%%%%%%%%%%%%%%%%%%%%%%%%%%%
Looking at the proof of Theorem~\ref{thm:masterboundSLipschitzHessian}, specifically equation~\eqref{eq:ineqgradnormxkplusLipschitzHessian}, we anticipate the need to lower-bound the norm of $\nabla \hat f_k(s_k)$. Owing to~\eqref{eq:nablahatfTgradf}, it holds that
\begin{align}
	\|\nabla \hat f_k(s)\| \geq \sigmamin\!\left(\D\Retr_{x_k}(s)\right) \|\grad f(\Retr_{x_k}(s))\|,
	\label{eq:nablahatfgradf}
\end{align}
where $\sigmamin$ extracts the smallest singular value of an operator.
For our purpose, it is important that this least singular value remains % Brian vouched for remains with s
bounded away from zero. This is only a concern for small steps (as large successful steps provide sufficient improvement for other reasons.) Providentially, for $s = 0$, Definition~\ref{def:retraction} ensures $\sigmamin(\D\Retr_{x_k}(0)) = 1$, so that by continuity we expect that it should be possible to meet this requirement. We summarize this discussion in the following assumption.
\begin{assumption} \label{assu:DRetr}
	There exist constants $a > 0$ and $b > 0$ such that, at each successful iteration $k$,
	\begin{align}
		\textrm{if } \|s_k\| \leq a, \textrm{ then } \sigmamin(\D\Retr_{x_k}(s_k)) & \geq b.
	\end{align}
	(The constant $a$ is allowed to be $+\infty$, while $b$ is necessarily at most 1.)
\end{assumption}
In the Euclidean case with $\Retr_x(s) = x+s$, $\D\Retr_x(s)$ is an isometry and one can set $a = +\infty$ and $b = 1$. We secure~\aref{assu:DRetr} in Section~\ref{sec:Dretr} for a large family of manifolds and retractions.


With these new assumptions, we can adapt Theorem~\ref{thm:masterboundSLipschitzHessian} to general retractions. The main change in the proof consists in treating short and long steps separately. This induces a condition that $\varepsilon$ must be small enough for the rate $O(1/\varepsilon^{1.5})$ to materialize. We stress that it is not necessary to know $L, L', q, a$ and $b$ as they appear in \aref{assu:firstorderregularityscalar}, \aref{assu:firstorderregularityvectorretraction} and \aref{assu:DRetr} to run Algorithm~\ref{algo:ARC} in practice: they are only used for the analysis.
Recall that $\varsigma_{\max}$ is provided by Lemma~\ref{lem:varsigmamax}.
\begin{theorem} \label{thm:masterboundSLipschitzHessianretraction}
	Let $\calM$ be a Riemannian manifold equipped with a retraction $\Retr$.
	Under~\aref{assu:lowerbound}, \aref{assu:firstorderregularityscalar}, \aref{assu:firstorderregularityvectorretraction} and \aref{assu:DRetr}, for an arbitrary $x_0 \in \calM$, let $x_0, x_1, x_2\ldots$ be the iterates produced by Algorithm~\ref{algo:ARC}.
	For any $\varepsilon > 0$, the \emph{total number} of \emph{successful} iterations $k$ such that $\|\grad f(x_{k+1})\| > \varepsilon$ is bounded above by
	\begin{align*}
		K_1(\varepsilon) & \triangleq \frac{3(f(x_0) - \flow)}{\eta_1 \varsigma_{\min}} \max\left( \left( \frac{ \frac{L'}{2} + \theta + \varsigma_{\max} }{ b - q(r) } \right)^{1.5} \frac{1}{\varepsilon^{1.5}}, \frac{1}{r^3} \right)
	\end{align*}
	for any $r \in (0, a]$ such that $q(r) < b$.
	Furthermore, $\lim_{k \to \infty} \|\grad f(x_k)\| = 0$.
\end{theorem}
\begin{proof}
	If iteration $k$ is successful, we have $x_{k+1} = \Retr_{x_k}(s_k)$.
	The gradient of the model $m_k$~\eqref{eq:mk} at $s_k$ is given by
	\begin{align*}
		\nabla m_k(s_k) & = \nabla \hat f_k(0) + \nabla^2 \hat f_k(0)[s_k] + \varsigma_k \|s_k\| s_k \\
						& = \nabla \hat f_k(s_k) + \left( \nabla \hat f_k(0) + \nabla^2 \hat f_k(0)[s_k] - \nabla \hat f_k(s_k) \right) + \varsigma_k \|s_k\| s_k.
	\end{align*}
%	Recall that $\nabla \hat f_k(s_k) = T_{s_k}^* \grad f(x_{k+1})$, with $T_{s_k} = \D\Retr_{x_k}(s_k)$~\eqref{eq:nablahatfTgradf}.
	Owing to the first-order progress condition~\eqref{eq:firstorderprogress}, using~\aref{assu:firstorderregularityvectorretraction} and~\eqref{eq:nablahatfgradf} with $T_{s_k} = \D\Retr_{x_k}(s_k)$,
	\begin{multline*}
		\theta \|s_k\|^2 \geq \|\nabla m_k(s_k)\| \geq \sigmamin(T_{s_k}) \|\grad f(x_{k+1})\| \\ - \frac{L'}{2} \|s_k\|^2 - q(\|s_k\|) \cdot \|\grad f(x_{k+1})\| - \varsigma_k \|s_k\|^2.
	\end{multline*}
	Rearranging and calling upon Lemma~\ref{lem:varsigmamax}, we get for all successful iterations $k$ that
	\begin{align}
		\big( \sigmamin(T_{s_k}) - q(\|s_k\|) \big) \cdot \|\grad f(x_{k+1})\| \leq \left(\frac{L'}{2} + \theta + \varsigma_{\max}\right) \|s_k\|^2.
	\label{eq:ineqgradnormxkplus}
	\end{align}	
	If $k$ is a successful step and $\|s_k\| \leq a$, then~\aref{assu:DRetr} guarantees $\sigmamin(T_{s_k}) \geq b > 0$. Additionally, since $q$ is continuous and satisfies $q(0) = 0$ there necessarily exists $r \in (0, a]$ such that $q(r) < b$. This motivates the following. Recall this subset of the successful steps: % based on the tolerance $\varepsilon$:
	\begin{align*}
	\calS_\varepsilon & = \{ k : \rho_k \geq \eta_1 \textrm{ and } \|\grad f(x_{k+1})\| > \varepsilon \}.
	\end{align*}
	Further partition this subset in two, based on step length: \emph{short} steps in $\calS_\short$ and \emph{long} steps in $\calS_\longg$. The partition is based on $r$ as constructed above:
	\begin{align*}
	\calS_\short & = \{ k \in \calS_\varepsilon : \|s_k\| \leq r \}, & \textrm{ and } &  & \calS_\longg & = \calS_\varepsilon \backslash \calS_\short.
	\end{align*}
	
	For $k \in \calS_\short$, we can lower-bound $\|s_k\|^3$ using~\eqref{eq:ineqgradnormxkplus} since $\sigmamin(T_{s_k}) - q(\|s_k\|) \geq b - q(r) > 0$. For $k \in \calS_\longg$, we have $\|s_k\|^3 > r^3$ by definition. Then, calling upon Proposition~\ref{prop:sumnormskcube}, we find
	\begin{align*}
		\frac{3(f(x_0) - \flow)}{\eta_1 \varsigma_{\min}} & \geq \sum_{k \in \calS_\short} \|s_k\|^3 + \sum_{k \in \calS_\longg} \|s_k\|^3 \\ & \geq \frac{ (b - q(r))^{1.5} \varepsilon^{1.5}}{\left(\frac{L'}{2} + \theta + \varsigma_{\max}\right)^{1.5}} |\calS_\short| + r^3 |\calS_\longg| \\
					& \geq \min\left( \left( \frac{ (b - q(r)) \varepsilon }{\frac{L'}{2} + \theta + \varsigma_{\max}} \right)^{1.5}, r^3 \right) |\calS_\varepsilon|.
	\end{align*}
	This proves the main claim. For limit points, see the matching argument in Theorem~\ref{thm:masterboundSLipschitzHessian}.
\end{proof}
A corollary identical to Corollary~\ref{cor:mastercorollary} holds for Theorem~\ref{thm:masterboundSLipschitzHessianretraction} as well.
In the Euclidean case with $\Retr_x(s) = x + s$ and a Lipschitz continuous Hessian, we can set $a = +\infty$, $b = 1$, $q \equiv 0$ and $r = +\infty$, thus also recovering the result of~\citep{bcgmt} exactly.

\TODOF{
\begin{remark}
	\TODO{I tried to find a short-but-correct way to state this, but it didn't quite work out. @Cora, can you suggest something?}
	\TODO{A4 is not needed for *convergence* of ARC (not complexity), even convergence to second order critical points.
		This is actually an important - but not straightforward- comment that we should make in the paper as well. In particular, if we only require say Cauchy decrease for first order criticality (instead of A3) we could prove that grad goes to zero and all limit points are critical. See ARC Part I paper Section 2. It is very similar to Trust region Cauchy-based approaches and so I think that the machinery in the TR complexity on manifold paper would translate here easily.The second-order criticality would stay similar to current ARC approach and it does not require A4.}
\end{remark}
}


\section{Second-order analysis} \label{sec:secondorder}
 
The two previous sections show how to meet first-order necessary optimality conditions approximately. To further satisfy second-order necessary optimality conditions approximately, we also
% Yes, we MUST: otherwise, start at a strict saddle point: without extra requirements, we are allowed to stay there (it's not even a question of deterministic / probabilistic algorithms: it's just that the current requirements are not enough to force escape of strict saddles in general.) -- but it's confusing.
require second-order progress in the subproblem solver, through condition~\eqref{eq:secondorderprogress}.

% This condition is similar to one proposed in~\citep{cartis2017improved} for the same purpose, though the proof of the theorem is different. Specifically, a direct extension of the proof in the referenced paper would involve the Hessian of the pullback at $s_k$ rather than at the origin. As we have seen for gradients, this leads to technical difficulties. 
This condition is similar to one proposed by~\citet{cartis2017improved} for the same purpose in the Euclidean case.
A direct extension of the proof in that reference would involve the Hessian of the pullback at the trial step $s_k$ rather than at the origin. %: $\nabla^2 \hat f_k(s_k)$.
As we have seen for gradients, this leads to technical difficulties.
We provide a proof that achieves the same complexity bound while avoiding such issues.
As a result, there is no need to distinguish between the exponential and the general retraction cases for second-order analysis.



We have the following bound on the total number of successful iterations which can produce points where the Hessian is far from positive semidefinite, akin to Theorems~\ref{thm:masterboundSLipschitzHessian} and~\ref{thm:masterboundSLipschitzHessianretraction}.
\begin{theorem} \label{thm:masterboundSsecond}
	Under~\aref{assu:lowerbound} and~\aref{assu:firstorderregularityscalar}, for an arbitrary $x_0\in\calM$, let $x_0, x_1, x_2\ldots$ be the iterates produced by Algorithm~\ref{algo:ARC} with second-order progress~\eqref{eq:secondorderprogress} enforced. For any $\varepsilon > 0$, the \emph{total number} of \emph{successful} iterations $k$ such that $\lambdamin(\nabla^2 \hat f_k(0)) < -\varepsilon$ is bounded above by
	\begin{align*}
		K_2(\varepsilon) & \triangleq \frac{3(f(x_0) - \flow)}{\eta_1 \varsigma_{\min}} (\theta + 2\varsigma_{\max})^3 \frac{1}{\varepsilon^3}.
	\end{align*}
	Furthermore, $\liminf_{k \to \infty} \lambdamin(\nabla^2 \hat f_k(0)) \geq 0$.
\end{theorem}
\begin{proof}
	The second-order condition~\eqref{eq:secondorderprogress} implies a lower-bound on step-sizes related to the minimal eigenvalue of the Hessian of $\hat f_k = f \circ \Retr_{x_k}$. Indeed, by definition of the model $m_k$~\eqref{eq:mk},
	\begin{align*}
		\forall s, \dot s \in \T_{x_k}\calM, && \nabla^2 m_k(s)[\dot s] & = \nabla^2 \hat f_k(0)[\dot s] + \varsigma_k\left( \|s\| \dot s + \frac{\inner{s}{\dot s}}{\|s\|}s \right).
	\end{align*}
	%(At $s = 0$, $\nabla^2 m_k(0) = \nabla^2 \hat f_k(0)$.)
	It follows that
	\begin{align*}
		\lambdamin(\nabla^2 \hat f_k(0)) & = \min_{\|\dot s\| = 1} \inner{\dot s}{\nabla^2 \hat f_k(0)[\dot s]} \\
			& = \min_{\|\dot s\| = 1} \inner{\dot s}{\nabla^2 m_k(s)[\dot s]} - \varsigma_k \left( \|s\|\|\dot s\|^2 + \frac{\inner{s}{\dot s}^2}{\|s\|} \right) \\ & \geq \lambdamin(\nabla^2 m_k(s)) - 2\varsigma_k \|s\|.
	\end{align*}
	In particular, with $s = s_k$, the second-order progress condition~\eqref{eq:secondorderprogress} and Lemma~\ref{lem:varsigmamax} yield
	\begin{align}
		-\lambdamin(\nabla^2 \hat f_k(0)) & \leq (\theta + 2\varsigma_{\max}) \|s_k\|.
		\label{eq:secondorderstepsizebig}
	\end{align}
	Consider this particular subset of the successful iterations:
	\begin{align*}
		\calS_\varepsilon & = \{  k : \rho_k \geq \eta_1 \textrm{ and } \lambdamin(\nabla^2 \hat f_k(0)) < -\varepsilon \}.
	\end{align*}
	Using Proposition~\ref{prop:sumnormskcube} with~\eqref{eq:secondorderstepsizebig} on this set leads to:
	\begin{align*}
		\frac{3(f(x_0) - \flow)}{\eta_1 \varsigma_{\min}} \geq |\calS_\varepsilon| \left(\frac{\varepsilon}{\theta + 2\varsigma_{\max}}\right)^3,
	\end{align*}
	which is the desired bound on the number of steps in $\calS_\varepsilon$. The limit inferior result follows from an argument similar to that at the end of the proof of Theorem~\ref{thm:masterboundSLipschitzHessian}.
\end{proof}
Here too, if $L$ is known we can set $\varsigma_0 = \varsigma_{\min} = \frac{1}{2}L$ (so that $\varsigma_{\max} = \frac{\gamma_3}{2(1-\eta_2)}L$) and $\theta = \frac{1}{2}L$. With those choices, we find that for a target depending on $L$ we have:
\begin{align}
	K_2(\sqrt{L\varepsilon}) & = \frac{6}{\eta_1} \left(\frac{1}{2} + \frac{\gamma_3}{(1-\eta_2)}\right)^3 (f(x_0) - \flow) \sqrt{L} \frac{1}{\varepsilon^{1.5}}.
\end{align}
This exhibits a complexity scaling with $\sqrt{L}$ when $L$ is known.

Theorem~\ref{thm:masterboundSsecond} is a statement about the Hessian of the pullbacks, $\nabla^2 \hat f_k(0)$, whereas we would more naturally desire a statement about the Hessian of the cost function itself, $\Hess f(x_k)$. For the exponential retraction, these two objects are the same~\eqref{eq:HessfHessfhatExp}. More generally, they are the same for any \emph{second-order retraction} (of which the exponential map is one example)~\citep[\S5]{AMS08}: we defer their (standard) definition to Section~\ref{sec:regularity}. For now, accepting the claim that for second-order retractions we have $\Hess f(x_k) = \nabla^2 \hat f_k(0)$, we get a more directly useful corollary: a complexity result for the computation of approximate second-order critical points on manifolds. The proof is in Appendix~\ref{app:secondorder}.
\begin{corollary} \label{cor:mastercorollarysecond}
	Under~\aref{assu:lowerbound} and~\aref{assu:firstorderregularityscalar},
	for an arbitrary $x_0\in\calM$
	and for any $\varepsilon_{g}, \varepsilon_{H} > 0$,
	if either
	\begin{enumerate}
		\item[(a)] we use the exponential retraction and~\aref{assu:firstorderregularityvector} holds, or
		\item[(b)] we use a second-order retraction and both~\aref{assu:firstorderregularityvectorretraction} and~\aref{assu:DRetr} hold,
	\end{enumerate}
	then
	Algorithm~\ref{algo:ARC} with second-order progress~\eqref{eq:secondorderprogress} enforced
	produces a point $x_k \in \calM$ such that $f(x_k) \leq f(x_0)$, $\|\grad f(x_k)\| \leq \varepsilon_{g}$ and $\lambdamin(\Hess f(x_k)) \geq -\varepsilon_{H}$
	in at most
	\begin{align*}
		\left(1 + \frac{|\log(\gamma_1)|}{\log(\gamma_2)}\right)\left({K_1(\varepsilon_{g}) + K_2(\varepsilon_{H}) + 1}\right) + \frac{1}{\log(\gamma_2)}\log\left(\frac{\varsigma_{\max}}{\varsigma_0}\right) + 1
	\end{align*}
	iterations, with $K_1$ as provided by Theorem~\ref{thm:masterboundSLipschitzHessian} or~\ref{thm:masterboundSLipschitzHessianretraction}, depending on assumptions.
\end{corollary}
If the retraction is not second order, we still get a bound on the eigenvalues of the Riemannian Hessian if the retraction has bounded acceleration at the origin at $x_k$: see~\citep[\S3.5]{boumal2016globalrates} and also Lemma~\ref{lem:derivativespullback} below.






%\clearpage

\section{Regularity assumptions} \label{sec:regularity}

The regularity assumptions~\aref{assu:firstorderregularityscalar} and~\aref{assu:firstorderregularityvectorretraction} pertain to the pullbacks $f \circ \Retr_{x_k}$. As such, they mix the roles of $f$ and $\Retr$. The purpose of this section is to shed some light on these assumptions, specifically in a way that disentangles the roles of $f$ and $\Retr$. In that respect, the main result is Theorem~\ref{thm:regularitycompact} below. Proofs are in Appendix~\ref{app:regularity}.

Each pullback is a function from a Euclidean space $\T_{x_k}\calM$ to $\reals$, so that standard calculus applies. Since the retraction is smooth by definition, pullbacks are as many times differentiable as $f$. This leads to the following simple fact.
\begin{lemma} \label{lem:lipschitzhessianbasic}
	Assume $f \colon \calM \to \reals$ is twice continuously differentiable.
	If there exists $L$ such that, for all $(x, s)$ among the sequence of iterates and trial steps $\{(x_0, s_0), (x_1, s_1), \ldots\}$ produced by Algorithm~\ref{algo:ARC}, with $\hat f = f \circ \Retr_x$, it holds that
	\begin{align}
		\left\| \nabla^2 \hat f(ts) - \nabla^2 \hat f(0) \right\|_{\mathrm{op}} & \leq t L \|s\|
		\label{eq:lipschitzhessianbasic}
	\end{align}
	for all $t \in [0, 1]$, then~\aref{assu:firstorderregularityscalar} holds with this $L$ and~\aref{assu:firstorderregularityvectorretraction} holds with $L' = L$ and $q \equiv 0$.
\end{lemma}

We call this a \emph{Lipschitz-type} assumption on $\nabla^2 \hat f$ because it compares the Hessians at~$ts$ and~0, rather than comparing them at two arbitrary points on the tangent space. On the other hand, we require this to hold on several tangent spaces with the same constant $L$.

In light of Lemma~\ref{lem:lipschitzhessianbasic}, one way to understand our regularity assumptions is to understand the Hessian of the pullback at points which are not the origin. The following lemma provides the necessary identities. The Hessian formula we have not seen elsewhere.
%
%\TODO{We really only need the notion of $c''$, which already appeared in Sec. 3. Later in this section, we also use the notation $\Ddt$. Might be able to move much of this to appendix; not sure if it's better---The statement of the lemma as well as its proof require some tools from Riemannian geometry that we also briefly mentioned in Section~\ref{sec:firstorderanalysisexp}---see~\cite[pp59--67]{oneill}. Specifically, we let $\nabla$ denote the Levi--Civita connection on $\calM$ (not to be confused with $\nabla$ and $\nabla^2$ which denote gradient and Hessian of functions on linear spaces, such as $\hat f$). With this notation, the Riemannian Hessian~\cite[Def.~5.5.1]{AMS08} is defined by $\Hess f = \nabla \grad f$. Furthermore, $\frac{\D}{\dt}$ denotes the covariant derivative of vector fields along curves on $\calM$, induced by $\nabla$. With this notation, given a smooth curve $c \colon \reals \to \calM$, the intrinsic acceleration is defined as $c''(t) = \frac{\D^2}{\dt^2} c(t)$. For example, for a Riemannian submanifold of a Euclidean space, $c''(t)$ is obtained by orthogonal projection of the classical acceleration of $c$ in the embedding space to the tangent space at $c(t)$. Geodesics are those curves which have zero intrinsic acceleration.}
Notation-wise, recall that $T^*$ denotes the adjoint of a linear operator $T$; furthermore, the \emph{intrinsic acceleration} $c''(t)$ of a smooth curve $c(t)$ is the covariant derivative of its velocity vector field $c'(t)$ on the Riemannian manifold $\calM$.
\begin{lemma} \label{lem:derivativespullback}
	Given $f \colon \calM \to \reals$ twice continuously differentiable and $x \in \calM$, the gradient and Hessian of the pullback $\hat f = f \circ \Retr_x$ at $s \in \T_x\calM$ are given by
	\begin{align}
		\nabla \hat f(s) & = T_s^* \grad f(\Retr_x(s)), \label{eq:gradientpullback} \\
		\nabla^2 \hat f(s) & = T_s^* \circ \Hess f(\Retr_x(s)) \circ T_s + W_s,
		\label{eq:hessianpullback}
	\end{align}
	where
	\begin{align}
		T_s & = \D\Retr_x(s) \colon \T_{x}\calM \to \T_{\Retr_x(s)}\calM
		\label{eq:Ts}
	\end{align}
	is linear, and $W_s$ is a symmetric linear operator on $\T_x\calM$ defined through polarization by
	\begin{align}
		\inner{W_s[\dot s]}{\dot s} & = \inner{\grad f(\Retr_x(s))}{c''(0)},
		\label{eq:Whessianpullback}
	\end{align}
	with $c''(0) \in \T_{\Retr_x(s)}\calM$ the intrinsic acceleration on $\calM$ of $c(t) = \Retr_x(s+t\dot s)$ at $t = 0$.
\end{lemma}
The particular case $s = 0$ connects to the comments around Corollary~\ref{cor:mastercorollarysecond}: since $T_0$ is identity by definition of retractions, we find that
\begin{align}
	\nabla^2 \hat f(0) & = \Hess f(x) + W_0,
	\label{eq:nablatwohatfHessfWzero}
\end{align}
where $W_0$ is zero in particular if the initial acceleration of retraction curves is zero (or if $\grad f(x) = 0$). As in~\citep[\S5]{AMS08}, this motivates the definition of second-order retractions, for which $\nabla^2 \hat f(0) = \Hess f(x)$.
\begin{definition}[Second-order retraction]\label{def:retrsecond}
	A retraction $\Retr$ on $\calM$ is \emph{second order} if, for any $x\in\calM$ and $\dot s \in \T_x\calM$, the curve $c(t) = \Retr_x(t\dot s)$ has zero initial acceleration: $c''(0) = 0$.
\end{definition}
Practical second-order retractions are often available~\citep[Ex.~23]{absil2012retractions}.
To prove our main result, we further restrict retractions.
%Combining Lemmas~\ref{lem:lipschitzhessianbasic} and~\ref{lem:derivativespullback} allows to guarantee~\aref{assu:firstorderregularity} holds in particular if $\calM$ is compact, $f$ is smooth and the retraction is nice enough. We formalize this in the definition below and in Theorem~\ref{thm:regularitycompact}. These are sufficient but likely not necessary conditions.
\TODOF{Make a more explicit comment regarding manifolds with negative curvature, since now have full details stating the operator norm of $T_s$ is $\cosh(\|s\|)$ on hyperbolic space, and that goes to infinity (exponentially) with $\|s\|$.---On the other hand, I'm not going to do much about it here, and it's already long. At least, as written now, the door is open for further refinements.}
\begin{definition}[Second-order nice retraction] \label{def:retrscndnice}
	Let $S$ be a subset of the tangent bundle $\T\calM$.
	A retraction $\Retr$ on $\calM$ is \emph{second-order nice on $S$} if there exist constants $c_1, c_2, c_3$ such that, for all $(x, s) \in S$ and for all $\dot s \in \T_x\calM$, all of the following hold:
	\TODOF{Are exponential maps second-order nice? Certainly not on manifolds with sectional curvature upper bounded by a negative number, because then even the sigmamin of DRetr is lower bounded by $\sinh(\|s\|)/\|s\|$ (roughly; check the Jacobi result), and this goes to infinity with $\|s\|$ (and pretty fast, too!) We may have to put in some restrictions on norm of $\|s\|$ (it could work because around 0 that function is pretty flat; though remember: it's the sigmamin, not the sigmamax) ... See notes end of NB36, Nov.\ 21, 2018}
	\begin{enumerate}
		\item $\forall \tauort \in [0, 1]$, $\|T_{ts}\|_{\mathrm{op}} \leq c_1$ where $T_{ts} = \D\Retr_x(ts)$; % is as in~\eqref{eq:Ts};
		\item $\forall \tauort \in [0, 1]$, the covariant derivative of $U(\tauort) = T_{\tauort s} \dot s$ satisfies  $\left\| \frac{\D}{\mathrm{d}\tauort} U(\tauort) \right\| \leq c_2 \|s\|\|\dot s\|$;  and
		\item $\|c''(0)\| \leq c_3 \|s\| \|\dot s\|^2 \textrm{ where } c(t) = \Retr_x(s+t\dot s)$.
	\end{enumerate}
	If this holds for $S = \T\calM$, we say $\Retr$ is \emph{second-order nice}.
\end{definition}
Second-order nice retractions are, in particular, second-order retractions (consider $s = 0$ in the last condition).
For $\calM$ a Euclidean space, the canonical retraction $\Retr_x(s) = x+s$ is second-order nice with $c_1 = 1$ and $c_2 = c_3 = 0$.
The classical retraction on the sphere is also second-order nice, with small constants $c_1, c_2, c_3$. We expect this to be the case for many usual retractions on compact or flat manifolds. % For manifolds with negative curvature, it might be necessary to work with a local version of these properties. % this last comment is because geodesics on Hadamard manifolds curve away from each other, so that T_s cannot have operator norm bounded above independently of s, at least for the exponential map.
For manifolds with negative curvature, the exponential retraction would lead $\opnorm{T_{s}}$ to grow arbitrarily large with $s$ going to infinity, hence it is important to consider restrictions to appropriate subsets, or to use another retraction.
\begin{proposition}\label{prop:sphereretrnice}
	For the unit sphere $\calM = \{ x\in\Rn : \|x\| = 1 \}$ as a Riemannian submanifold of $\Rn$, the retraction $\Retr_x(s) = \frac{x+s}{\|x+s\|}$ is second-order nice with $c_1 = 1, c_2 = \frac{4\sqrt{3}}{9}, c_3 = 2$.
\end{proposition}
\begin{theorem} \label{thm:regularitycompact}
%	\TODO{This claim should probably be replaced by something of the form: if $f$ has Riemannian Lipschitz properties, then we can infer the retraction Lipschitz properties.}
	Let $f \colon \calM \to \reals$ be three times continuously differentiable.
	Assume the retraction is second-order nice on the set $\{(x_0, s_0), (x_1, s_1), \ldots \}$ of points and steps  generated by Algorithm~\ref{algo:ARC} (see Definition~\ref{def:retrscndnice}).
	If the sequence $x_0, x_1, x_2, \ldots$ remains in a compact subset of $\calM$,
%	\begin{enumerate}
%		\item[(a)] $\calM$ is compact, or
%		\item[(b)] the sequence $\{x_0, x_1, \ldots\}$ remains in a compact subset of $\calM$, % and~\aref{assu:firstorderprogress} is satisfied,
%	\end{enumerate}
	then \aref{assu:firstorderregularityscalar} and \aref{assu:firstorderregularityvectorretraction} are satisfied
	with a same $L$ (related to the Lipschitz properties of $f$, $\grad f$ and $\Hess f$: see the proof for an explicit expression) and $q \equiv 0$.
%	(see the proof for a bound on $L$). \TODO{$L_1, L_2$? $q \equiv 0$? ... with $L_1 = L_2 = ...$ and $q \equiv 0$?}
	%
	% Under~\aref{assu:firstorderprogress},
	
	Algorithm~\ref{algo:ARC} is a descent method, so that the last condition holds in particular if the sublevel set $\{x \in \calM : f(x) \leq f(x_0)\}$ is compact, and a fortiori if $\calM$ is compact.
	%This in turn is the case in particular if $\calM$ is compact.
%	\TODO{We can bound the norm of $s_k$ for all $x_k$ using~\aref{assu:firstorderprogress}. Surely we can deduce boundedness of $\calN$ assuming the $x_k$'s are bounded? Can we then also deduce the closure of $\calN$ is compact? If so, we could entirely remove the definition of $\calN$ from the statement, which would clean things up a bit.} \TODO{Careful: see last page of NB36: closed and bounded for a subset of a manifold does NOT imply compact. E.g.: the open disk in $\reals^2$. How about assuming the sequence of $x_k$'s is in a compact set instead? If I take a closed ball of radius $r$ in the tangent space of each $x$ in a compact subset of $\calM$, is the resulting subset of $\T\calM$ compact? Pretty sure it's true---actually, we used that in the paper with PAA. Hard to track down though. There is a discussion about it here: \url{https://math.stackexchange.com/questions/1177264/compactness-of-a-subset-of-a-tangent-bundle}---to finish the argument in that response, further need that closed susbets of compact sets are compact: \url{https://math.stackexchange.com/questions/212181/a-subset-of-a-compact-set-is-compact}, \url{https://planetmath.org/closedsubsetsofacompactsetarecompact}. And it's related to an exercise in Lee's book, ex. 10-19.}
%
%	Consider
%	\begin{align}
%		\calN & = \bigcup_{k} \{ \Retr_{x_k}(ts_k) : t \in [0, 1] \},
%	\end{align}
%	that is, the subset of $\calM$ obtained by collecting all curves generated by retracted steps (both accepted and rejected).
%	If the closure of $\calN$ is a compact subset of $\calM$, then~\aref{assu:firstorderregularity} is satisfied. \TODO{We don't need the closure of $\calN$: we just need $\calN$ to be included in a compact set. ...Do we even need to talk about compactness on the tangent bundle at any point? Yes. But do we then need second-order nice on the containing compact set though? NO we don't: they're separate concerns. The compactness is just to bound some differentials of $f$ owing to them being continuous.}
	% It's not enough to ask for boundedness of calN -- it's true that this would imply the closure of calN is both closed and bounded, but it wouldn't imply compactness because calM is not a Euclidean space: see https://en.wikipedia.org/wiki/Compact_space#Open_cover_definition -- see also right below the definition of compactness of a subset.
	% See also Def. 7 and 8 here: http://www.ms.uky.edu/~ken/ma570/lectures/lecture2/html/compact.htm
%	\TODO{To argue that $\calN$ is in a compact set provided the set of $x$'s is in a compact set, might be able to use that the image of a compact set through a continuous function is a compact set. Then, applying the (continuous) retraction to a compact subset of the tangent bundle must yield a compact subset of the manifold: \url{http://mathonline.wikidot.com/continuous-functions-on-compact-sets-of-metric-spaces}.}
\end{theorem}
%The closure of $\calN$ is compact in particular if $\calM$ is compact%
% See http://planetmath.org/closedsubsetsofacompactsetarecompact
%
%or if the sublevel set of $f(x_0)$ is compact
%, since closed subsets of compact sets are compact. % ~\citep[p68]{berge1963topological}.
% This part about sublevel set was removed because it's tricky: f is decreasing on the iterates, but who knows: it might spike up along the retracted path..... --- This is fixed now, by using that the iterates themselves are always in the sublevel set, and bringing in the fact that the steps are bounded.
%Of course, if $\calM$ is compact, the theorem applies (and in fact, \aref{assu:firstorderprogress} is not required in that case.)


This general result shows existence of (loose) bounds for the Lipschitz constants. For specific optimization problems, it is sometimes easy to derive more accurate constants by direct computation: see for example~\citep[App.~D]{criscitiello2019escapingsaddles} for PCA.


%\clearpage



\section{Controlling the differentiated retraction} \label{sec:Dretr}


In Theorem~\ref{thm:masterboundSLipschitzHessianretraction}, we control the worst-case running time of ARC via the differential of the retraction $\Retr$, through~\aref{assu:DRetr}. This assumption, which does not come up in the Euclidean case, involves constants $a, b$ to control $\sigmamin(\D\Retr_{x_k}(s_k))$. Our first result shows~\aref{assu:DRetr} is satisfied for some $a$ and $b$ for any retraction on any manifold, provided the sublevel set of $f(x_0)$ is compact (see also Theorem~\ref{thm:regularitycompact}). This is mostly a topological argument. Proofs are in Appendix~\ref{sec:proofsDretr}.
\begin{theorem} \label{thm:retractionsgeneralmainnew}
	Let $\Retr$ be a retraction on a Riemannian manifold $\calM$, and let $\calU$ be a nonempty compact subset of $\calM$. For any $b \in (0, 1)$ there exists $a > 0$ such that, for all $x \in \calU$ and $s \in \T_x\calM$ with $\|s\|_x \leq a$, we have $\sigmamin(\D\Retr_x(s)) \geq b$. In particular, \aref{assu:DRetr} is satisfied with such $(a, b)$ provided the iterates $x_0, x_1, x_2, \ldots$ remain in $\calU$.
\end{theorem}

We further quantify the constants $a$ and $b$ in two cases of interest:
\begin{enumerate}
	\item For the Stiefel manifold $\St(n,p) = \{X \in \mathbb{R}^{n \times p} : X\transpose X = I_p\}$ as a Riemannian submanifold of $\Rnp$ with the usual inner product $\inner{A}{B} = \trace(A\transpose B)$, we explicitly control $(a, b)$ for the popular Q-factor retraction ($\Retr_X(S)$ is obtained by Gram--Schmidt orthonormalization of the columns of $X+S$). Special cases include the sphere ($p=1$) and the orthogonal group ($p = n$).
	\item For complete manifolds with bounded sectional curvature, we control $(a, b)$ for the case of the exponential retraction. Important special cases include Euclidean spaces (flat manifolds), manifolds with nonpositive curvature (Hadamard manifolds, including the manifold of positive definite matrices~\citep{moakher2006symmetric,Bhatia}), and compact manifolds \citep[\S9.3, p166]{bishop1964geometry}.
\end{enumerate}

The first result follows from a direct calculation. 
\begin{proposition} \label{prop:stiefelmastercorollary}
	For the Stiefel manifold with the Q-factor retraction,
	for any $a > 0$, define $b  = 1 - 3a - \frac{1}{2}a^2$. If $b$ is positive, then~\aref{assu:DRetr} holds with these $a$ and $b$. Moreover, for the sphere, we have that~\aref{assu:DRetr} is satisfied for any $a > 0$ and $b = \frac{1}{1 + a^2}$.
\end{proposition}

The second result follows from the connection between the differential of the exponential map and certain \emph{Jacobi fields} on $\calM$, together with standard comparison theorems from Riemannian geometry~\citep[Ch.~10, 11]{lee2018riemannian}.
\begin{proposition} \label{prop:mastercorollaryexponetial}
	Let $\calM$, complete,
	% We need complete so that Exp is defined over all of TM, just like all of the retractions we work with.
	have sectional curvature upper bounded by $C$, and let the retraction $\Retr$ be the exponential retraction $\Exp$:
	\begin{itemize}
		\item[] If $C \leq 0$, then~\aref{assu:DRetr} is satisfied for any $a > 0$ and $b = 1$;
		\item[] If $C > 0$, then~\aref{assu:DRetr} is satisfied for any $0 < a < \frac{\pi}{\sqrt{C}}$ and $b = \frac{\sin(a\sqrt{C})}{a\sqrt{C}}$.
	\end{itemize}
\end{proposition}



%\clearpage


\section{Solving the subproblem} \label{sec:subproblem}

%\TODO{Explain Lanczos step better; in particular: the fact there is a random push if need be.}

\TODOF{Parlett, 'A New Look at the Lanczos Algorithm for Solving Symmetric Systems of Linear Equations' has 'old tricks' to do Lanczos right in the face of round-off error; it's rather involved though; it'd be nice to get into that 'at some point'.}

At each iteration, Algorithm~\ref{algo:ARC} requires the approximate minimization of the model $m_k$~\eqref{eq:mk} in the tangent space $\T_{x_k}\calM$. Since the latter is a linear space, this is the same subproblem as in the Euclidean case. In contrast to working simply over $\Rn$ however, one practical difference is that we do not usually have access to a preferred basis for $\T_{x_k}\calM$, so that it is preferable to resort to basis-free solvers. To this end, we describe a Lanczos method as Algorithm~\ref{algo:lanczosoptimized}.

Let us phrase the subproblem in a general context. Given a vector space $\calX$ of dimension $n$ with an inner product $\inner{\cdot}{\cdot}$ (and associated norm $\|\cdot\|$), an element $g \in \calX$, a self-adjoint linear operator $H \colon \calX \rightarrow \calX$ and a real $\varsigma > 0$, define the function $m\colon\mathcal{X} \rightarrow \reals$ as 
\begin{align}
	m(s) = \inner{g}{s} + \frac{1}{2}\inner{s}{H(s)} + \frac{\varsigma}{3}\|s\|^3. 	
	\label{eqn:modelfunctionm}
\end{align}
We wish to compute an element $s \in \calX$ such that
\begin{align}
	m(s) & \leq m(0) & \textrm{ and } & & \|\nabla m(s)\| & \leq \theta \|s\|^2.
	\label{eqn:subsolvercondition}
\end{align}
%and possibly also such that
%\begin{align}
%	\lambdamin(\nabla^2 m(s)) & \geq - \theta \|s\|.
%	\label{eqn:subsolverconditionsecond}
%\end{align}
This corresponds to satisfying condition~\eqref{eq:firstorderprogress}
%These correspond to satisfying condition~\eqref{eq:firstorderprogress} and, optionally, condition~\eqref{eq:secondorderprogress}
at iteration $k$, where $\calX = \T_{x_k}\calM$ is endowed with the Riemannian inner product at $x_k$, $g = \nabla \hat f_k(0)$, $H = \nabla^2 \hat f_k(0)$ and $\varsigma = \varsigma_k$. If $g = 0$, then $s = 0$ satisfies the condition: henceforth, we assume $g \neq 0$.
%the underlying vector space is the tangent space of the manifold and $g$ is the gradient at the current iterate.

%\paragraph{Lanczos Method.}
Certainly, a global minimizer of~\eqref{eqn:modelfunctionm} meets our requirements (it would also satisfy the equivalent of the second-order condition~\eqref{eq:secondorderprogress}). Such a minimizer can be computed, but known procedures for this task involve a diagonalization of $H$, which may be expensive. Instead, we use the Lanczos-based method proposed in~\citep[\S6]{cartis2011adaptivecubicPartI}: %  \citep[Lect.~36]{trefethen1997numerical}
%to produce iteratively growing subspaces within which we minimize the above cubic.
the latter iteratively produces a sequence of orthonormal vectors $\{q_1, \ldots, q_n\}$ and a symmetric tridiagonal matrix $T$ of size $n$ such that~\citep[Lec.~36]{trefethen1997numerical}\footnote{In case of so-called breakdown in the Lanczos iteration at step $k$, we follow the standard procedure which is to generate $q_k$ as a random unit vector orthogonal to $q_1, \ldots, q_{k-1}$, then to proceed as normal. This does not jeopardize the desired properties~\eqref{eqn:lanczosmethod}.\label{footnotelanczos}}
\begin{align}
	q_1 & = \frac{g}{\|g\|} &  \textrm{ and } & & T_{ij} & = \inner{q_i}{H(q_j)} \textrm{ for all } i,j \textrm{ in } 1\ldots n.
	\label{eqn:lanczosmethod}
\end{align}
Let $T_k$ denote the $k \times k$ principal submatrix of $T$: producing $q_1, \ldots, q_k$ and $T_k$ requires exactly $k$ calls to $H$.
%
%The Lanczos method guarantees that, for all $k$, vectors $\{q_1, \ldots, q_k\}$ form an orthonormal basis for a subspace which contains the Krylov subspace $\calK_k$ spanned by $\{ g, H(g), \ldots, H^{k-1}(g)\}$, where $H^k$ represents the $k$-times composition of $H$.\footnote{As is standard in Krylov methods~\citep[Ex.~33.2]{trefethen1997numerical}, if the Krylov space $\calK_k$ has dimension strictly less than $k$ (known as a breakdown), we continue the construction by generating $q_k$ at random.}
%
Consider $m(s)$~\eqref{eqn:modelfunctionm} restricted to the subspace spanned by $q_1, \ldots, q_k$:
\begin{align}
	\forall y \in \Rk, \textrm{ with } s & = \sum_{i = 1}^{k} y_i q_i, &   m(s) & = \inner{g}{y_1q_1} + \frac{1}{2} \sum_{i,j=1}^k y_i y_j \inner{q_i}{H(q_j)} + \frac{\varsigma}{3} \|y\|^3 \nonumber\\
		 & & & = y_1 \|g\| + \frac{1}{2} y\transpose T_ky + \frac{\varsigma}{3} \|y\|^3,
		 \label{eq:modelrestrictedlanczos}
\end{align}
where $\|\cdot\|$ is also the 2-norm over $\Rk$.
Since $T_k$ is tridiagonal, it can be diagonalized efficiently. As a result, it is inexpensive to compute a global minimizer of $m(s)$ restricted to the subspace spanned by $q_1, \ldots, q_k$~\citep[\S6.1]{cartis2011adaptivecubicPartI}.
%
%This cubic in $y$ can be minimized efficiently using the explicit procedure alluded to above. Specifically, because $T_k$ is tridiagonal, it is inexpensive to eigendecompose it.
Furthermore, since the Lanczos basis is constructed incrementally, we can minimize the restricted cubic at $k = 1$, check the stopping criterion~\eqref{eqn:subsolvercondition}, and proceed to $k = 2$ only if necessary, etc. The hope (borne out in experiments) is that the algorithm stops well before $k$ reaches $n$ (at which point it necessarily succeeds.) In this way, we limit the number of calls to $H$, which is typically the most expensive part of the process. This general strategy was first proposed in the context of trust-region subproblems by~\citet{gould1999solving}. Algorithm~\ref{algo:lanczosoptimized} is set up to target first order progress only. In order to ensure satisfaction of second-order progress, one may have to force the execution of $n$ iterations (which is rarely done in practice).

%\TODO{Step $11$ of Algorithm~\ref{algo:lanczosoptimized} is performed via a Newton-based root-finding technique detailed in \cite{cartis2011adaptivecubic}. Note that the tridiagonal structure of $T$ can be leveraged in the computation of various quantities in the inner Newton solver.}

%\paragraph{Checking the stopping criterion.}

%Algorithm~\ref{algo:lanczosoptimized} outlines a procedure to achieve~\eqref{eqn:subsolvercondition} based on the Lanczos method. The key idea at each iteration $k$ of the algorithm is to minimize $m(s)$ over the subspace spanned by $\{q_1, \ldots, q_k\}$. The algorithm continues till condition~\eqref{eqn:subsolvercondition} is met. This necessarily happens if $k = n$. The hope is for it to happen much earlier.

%% This is confusing because we actually build a basis; I made this point at the beginning, let's just remove this.
%Importantly, the whole process is basis free: it involves only linear combinations and inner products of vectors in $\calX$, and calls to $H$ as an operator (as opposed to requiring a matrix representation of $H$, for example). %; that is, it is basis free. %This is particularly well suited for optimization on manifolds, where tangent vectors are usually not represented in particular coordinates, and the Hessian is only available as an operator (rather than a matrix).

We draw attention to a technical point. Upon minimizing~\eqref{eq:modelrestrictedlanczos}, we obtain a vector $y \in \Rk$. To check the stopping criterion~\eqref{eqn:subsolvercondition}, we must compute $\|\nabla m(s)\|$, where $s = \sum_{i=1}^k y_i q_i$. Since
\begin{align}
	\nabla m(s) & = g + H(s) + \varsigma \|s\| s,
\end{align}
one approach involves computing $s$ (that is, form the linear combination of $q_i$'s) and applying $H$ to~$s$: both operations may be expensive in high dimension. An alternative (shown in Algorithm~\ref{algo:lanczosoptimized}) is to recognize that, due to the inner workings of Lanczos iterations, $\nabla m(s)$ lies in the subspace spanned by $\{q_1, \ldots, q_{k+1}\}$ (if $k < n$). Explicitly, %Indeed, using all properties in~\eqref{eqn:lanczosmethod}-----it's not just that; it's also using that H(s) is in the span of q_1, ..., q_{k+1}; don't want to go into those details
\begin{align}
	\nabla m(s) & = \|g\|q_1 + \sum_{i = 1}^{k+1} (T_{1:k+1,1:k}^{} \, y)_i q_i + \varsigma\|y\| \sum_{i=1}^{k}y_iq_i,
	\label{eqn:subsolvernormofgrad}
\end{align}
%where $T_{a:b, c:d}$ is the submatrix of $T$ containing rows $\{a, \ldots, b\}$ and columns $\{c, \ldots, d\}$.
where $T_{1:k+1,1:k}$ is the submatrix of $T$ containing the first $k+1$ rows and first $k$ columns.
%Thus,
%\begin{align}
%	\|\nabla m(s)\| = \big\|\|g\|e_1 + T_{1:k+1,1:k}y + \varsigma\|y\|y\big\|,
%	\label{eqn:subsolvernormofgrad}
%\end{align}
%where $e_1 \in \R^{k+1}$ --- \TODO{actually: $y$ is in Rk, so it would technically need to be extended: let's just not go there.}.
This expression gives a direct way to compute $\|\nabla m(s)\|$ without forming $s$ and without calling $H$,
simply by running the Lanczos iteration one step ahead.
%This modification reduces the number of calls to $H$ by a factor of two, and postpones forming the vector $s$ until the algorithm terminates.

%% Reviewers objected to this.
%We mention here that Lanczos is notoriously numerically unstable. Furthermore, while the stopping criterion in~\aref{assu:firstorderprogress} is appropriate to control worst-case iteration complexity, it may not lead to optimal performance for natural problems. We believe improvements on both these aspects could benefit practical performance.

\begin{remark} \label{rem:carmonduchisubproblem}
%	As per~\aref{assu:firstorderprogress} (and~\aref{assu:secondorderprogress}), it is only necessary to produce an approximate (possibly second-order) critical point of the subproblem: there is no global minimization aspect to these conditions, so that many algorithms can apply. Regarding the Lanczos method in particular,
	\citet{carmon2018analysis} analyze the number of Lanczos iterations that may be required to reach approximate solutions to the subproblem. For the Euclidean case, they conclude that the overall complexity of ARC to compute an $\varepsilon$-critical point in terms of Hessian-vector products (which dominate the number of cost and gradient computations) is $O(\varepsilon^{-7/4})$. Furthermore, the dependence on the dimension of the search space is only logarithmic. (The logarithmic terms are caused by the need to randomize for the so-called hard case---see the reference for important details in that regard.)
	% As the subproblem is posed on a linear space in the Riemannian case as well,
	Their conclusions should extend to the Riemannian setting as well.
	%\TODO{Some unresolved points here about easy/hard case. Also, we talk about this twice: in the ''related work'', could reference this remark for example.}
%	
%	\TODO{Include remark about Carmon \and Duchi regarding overall Hessian-vector product complexity using Lanczos. Check carefully what their assumptions are regarding requirements for the subproblem solve. Can also mention that, really, any method which can compute an approx second order critical point for the subproblem would do.}
%	\TODO{we need to note that we only need a solver that can find an approximate second-order critical point of the cubic model; no global aspect. So any method that can do that will be fine and one can apply their complexity bounds - this is not optimal but will result in dimension-independent bounds under appropriate choice of method. Note that we do not know how to quantify Lanczos except by extending the recent approach of Carmon and Duchi (who do not address the hard case... --- well, they have a randomized approach, no?)}
	\citet{gould2019regularizedquadratic} extend these results to study the decrease of the norm of the gradient specifically, as required here.
\end{remark}


\begin{algorithm}[t]
	\caption{Lanczos-based cubic model subsolver}
	\label{algo:lanczosoptimized}
	\begin{algorithmic}[1]
		\State \textbf{Parameters:} $\theta , \varsigma > 0$, vector space $\calX$ with inner product $\inner{\cdot}{\cdot}$ and norm $\|\cdot\|$
		\State \textbf{Input:} $g \in \calX$ nonzero, a self-adjoint linear operator $H\colon\calX \rightarrow \calX$
		%		\State \textbf{Init:} $k \leftarrow 0$
		\Statex 
		\State{$k \leftarrow 1$}
		\State{Obtain $q_1, T_1$ via a Lanczos iteration \eqref{eqn:lanczosmethod}}
		\State{Solve $y^{(1)} = \argmin{y \in \reals}{\;\;\|g\|y + \frac{1}{2}T_{1}y^2 + \frac{1}{3}\varsigma |y|^3}$ }
		\State{Obtain $q_2, T_2$ via a Lanczos iteration \eqref{eqn:lanczosmethod}}
		\State{Compute $\|\nabla m(s_1)\|$ via~\eqref{eqn:subsolvernormofgrad}}, where $s_1 = y^{(1)} q_1$
		\Statex
		\While{$\|\nabla m(s_k)\| > \theta \|s_k\|^2$}
		\State{$k \leftarrow k + 1$}
		\State{Solve $y^{(k)} = \argmin{y \in \reals^k}{\;\;\|g\| y_1 + \frac{1}{2}y\transpose T_ky + \frac{1}{3}\varsigma \|y\|^3}$ following \citep[\S6.1]{cartis2011adaptivecubicPartI}}
		\State{Obtain $q_{k+1}, T_{k+1}$ via a Lanczos iteration \eqref{eqn:lanczosmethod}}
		\State{Compute $\|\nabla m(s_k)\|$ via~\eqref{eqn:subsolvernormofgrad}}, where $s_k = \sum_{i = 1}^k y^{(k)}_i q_i^{}$
		\EndWhile
		\Statex
		\State \textbf{Output:} $s_k \in \calX$
	\end{algorithmic}
	
\end{algorithm}


%\TODO{Comment somewhere about first-order and second-order conditions \eqref{eq:firstorderprogress} \eqref{eq:secondorderprogress}: if we run $n$ iterations of the Lanczos procedure, in exact arithmetic, we get a global optimizer, so, all good.}

\section{Numerical experiments} \label{sec:XP}


We implement Algorithm~\ref{algo:ARC} within the Manopt framework~\citep{manopt} (our code is part of that toolbox) and compare the performance of our implementation against some existing solvers in that toolbox, namely, the Riemannian trust-region method (RTR)~\citep{genrtr} and the Riemannian conjugate gradients method with Hestenes--Stiefel update formula (CG-HS)~\citep[\S8.3]{AMS08}. All algorithms terminate when $\|\grad f(x_k)\| \leq 10^{-9}$. CG-HS also terminates if it is unable to produce a step of size more than $10^{-10}$. Code to reproduce the experiments is available at \url{https://github.com/NicolasBoumal/arc}. We report results with randomness fixed by \texttt{rng(2019)} within Matlab R2019b.
%Our code will become available as part of this package on the corresponding GitHub repository.\footnote{\url{https://github.com/NicolasBoumal/manopt}} % -- suppressed to improve anonymity

We consider a suite of six Riemannian optimization problems:
\begin{enumerate}
	\item Dominant invariant subspace: $\max_{X \in \Gr(n, p)} \frac{1}{2} \Trace(X\transpose A X)$, where $A \in \Rnn$ is symmetric (randomly generated from i.i.d.\ Gaussian entries) and $\Gr(n, p)$ is the Grassmann manifold of subspaces of dimension $p$ in $\Rn$, represented by orthonormal matrices in $\Rnp$. Optima correspond to dominant invariant subspaces of $A$~\citep{edelman1998geometry}.
	\item Truncated SVD: $\max_{U\in\St(m, p), V\in\St(n, p)} \Trace(U\transpose A V N)$, where $\St(n, p)$ is the set of matrices in $\Rnp$ with orthonormal columns, $A \in \Rmn$ has i.i.d.\ random Gaussian entries and $N = \diag(p, p-1, \ldots, 1)$. Global optima correspond to the $p$ dominant left and right singular vectors of $A$~\citep{sato2013riemannianSVD}. For this and the previous problem, the random matrices have small eigen or singular value gap, which makes them challenging.
	\item Low-rank matrix completion via optimization on one Grassmann manifold, as in~\citep{boumal2011rtrmc}. The target matrix $A \in \Rmn$ has rank $r$: it is fully specified by $r(m+n-r)$ parameters. We observe this many entries of $A$ picked uniformly at random, times an oversampling factor (osf). The task is to recover $A$ from those samples. $A$ is generated from two random Gaussian factors as $A = LR\transpose$ to have rank $r$ exactly. The variable is $U \in \Gr(m, r)$, and the cost function minimizes the sum of squared errors between $UW_U$ and the observed entries of $A$, where $W_U \in \reals^{r \times n}$ is the optimal matrix for that purpose (which has an explicit expression once $U$ is fixed, efficiently computable).
	\item Max-cut: given the adjacency matrix $A \in \Rnn$ of a graph with $n$ nodes, we solve the semidefinite relaxation of the Max-Cut graph partitioning problem via the Burer--Monteiro formulation~\citep{burer2005local} on the oblique manifold~\citep{journee2010low}: $\min_{X\in\mathrm{OB}(n, p)} \frac{1}{2}\Trace(X\transpose A X)$, where $\mathrm{OB}(n, p)$ is the set of matrices in $\Rnp$ with unit-norm rows. Here, we pick graph \#22 from a collection of graphs called Gset (see any of the references): it has $n = 2000$ nodes and $19990$ edges, and $p$ is set close to $\sqrt{2n}$ as justified in~\citep{boumal2018deterministicbm}.
	\item Synchronization of rotations: $m$ rotation matrices $Q_1, \ldots, Q_m$ in the special orthogonal group $\mathrm{SO}(d)$ are estimated from noisy relative measurements $H_{ij} \approx Q_i^{}Q_j\transpose$ for an Erd\H{o}s--R\'enyi random set of pairs $(i, j)$ following a maximum likelihood formulation, as in~\citep{boumal2013MLE}. The specific distribution of the measurements and the corresponding cost function are described in the reference. All algorithms are initialized with the technique proposed in the reference, to avoid convergence to a poor local optimum.
	\item ShapeFit: least-squares formulation of the problem of recovering a rigid structure of $n$ points $x_1, \ldots, x_n$ in $\Rd$ from noisy measurements of some of the pairwise directions $\frac{x_i - x_j}{\|x_i - x_j\|}$ picked uniformly at random, following~\citep{hand2018shapefit}. The set of points is centered and obeys one extra linear constraint to fix scaling ambiguity, so that the search space is effectively a linear subspace of $\reals^{n\times d}$: this is the manifold $\calM$. The cost function as spelled out in the reference makes this a structured linear least-squares problem.
\end{enumerate}
%We do not claim that these are the best ways to address those problems: they serve merely as a test bed for optimization algorithms. In particular, the formulation of truncated SVD suffers from poor conditioning due to the small singular value gap of Gaussian matrices, resulting in many Hessian-vector products for both RTR and ARC. We also note that these problems have rather benign non-convexity: all solvers tend to converge to global optima in the proposed settings, so that we only investigate their convergence rate.

For each problem, we generate one instance and one random initial guess (except for problem 5 which is initialized deterministically). Then, we run each algorithm from that same initial guess on that same instance. Figure~\ref{fig:XPtimegradnorm} displays the progress of each algorithm on each problem as the gradient norm of iterates (on a log scale) as a function of elapsed computation time to reach each iterate (in seconds) on a laptop from 2016. For the same run, Figure~\ref{fig:XPhesscallsgradnorm} reports the number of gradient calls and Hessian-vector products (summed) issued by all algorithms along the way. Figure~\ref{fig:XPouteriterations} reports the number of outer iterations for ARC and RTR, that is, excluding work done by subsolvers.


For ARC, we report results with $\theta = 0.25$ and $\theta = 2$; this is used in the stopping criterion for subproblem solves following~\eqref{eq:firstorderprogress} (we do not check~\eqref{eq:secondorderprogress}). To initialize $\varsigma_0$, we use 100 divided by the initial trust-region radius of RTR ($\Delta_0$) chosen by Manopt. Other parameters of ARC are set as follows: $\varsigma_{\min} = 10^{-10}, \eta_1 = 0.1, \eta_2 = 0.9, \gamma_1 = 0.1, \gamma_2 = \gamma_3 = 2$, with update rule:
\begin{align}
\varsigma_{k+1} =
\begin{cases}
\begin{aligned}
& \max(\varsigma_{\min}, \gamma_1 \varsigma_k) & & \textrm{ if } \rho_k \geq \eta_2 & & \textrm{ (very successful),} \\
& \varsigma_k & & \textrm{ if } \rho_k \in [\eta_1, \eta_2) & & \textrm{ (successful),} \\
& \gamma_2 \varsigma_k & & \textrm{ if } \rho_k < \eta_1 & & \textrm{ (unsuccessful).}
\end{aligned}
\end{cases}
\end{align}
Standard safeguards to account for numerical round-off errors are included in the code (not described here). Parameters for the other methods have the default values given by Manopt.

We experiment with two subproblem solvers for ARC. The first one is the Lanczos-based method as described in Section~\ref{sec:subproblem} (ARC Lanczos). The second one is a (Euclidean) nonlinear, nonnegative-Polak--Ribi\`ere conjugate gradients method run on the model $m_k$~\eqref{eq:mk} in the tangent space $\T_{x_k}\calM$ (ARC NLCG), implemented by Bryan Zhu~\citep{zhu2019arcthesis}. This solver uses the initialization recommended by~\citet{carmon2016gradient} for gradient descent (and from which they proved convergence to a global optimizer, despite non-convexity of the model), and exact line-search.
%While theoretical guarantees for this subproblem solver are not as clear as for the Lanczos-based one,
We find that this subproblem solver performs well in practice. It is simpler to implement, and uses less memory than the Lanczos method.








\begin{figure}[p]
	\centering
	\includegraphics[width=1\linewidth]{Experiments/compare_solvers_Oct_15_2019_174010_time_gradnorm}
	\caption{Gradient norm at each iterate for the three competing solvers on the six benchmark problems, as a function of computation time needed by those solvers to reach those iterates (in seconds). Our algorithm is ARC (tested with two different subproblem solvers, each with two parameter settings); RTR is Riemannian trust-regions; CG-HS is Riemannian conjugate gradients with Hestenes--Stiefel step selection.}
	\label{fig:XPtimegradnorm}
\end{figure}
\begin{figure}[p]
	\centering
	\includegraphics[width=1\linewidth]{Experiments/compare_solvers_Oct_15_2019_174010_gradhesscalls_gradnorm}
	\caption{Gradient norm at each iterate, as a function of the number of gradient calls and Hessian-vector calls (the sum of both) issued by those solvers to reach those iterates.}
	\label{fig:XPhesscallsgradnorm}
\end{figure}
\begin{figure}[p]
	\centering
	\includegraphics[width=1\linewidth]{Experiments/compare_solvers_Oct_15_2019_174010_iter_gradnorm}
	\caption{Gradient norm at each iterate, as a function of the number of outer iterations for ARC and RTR: both of these solvers rely on a subproblem solver. This plot compares the behavior of the algorithms separately from their subproblem solvers' work. As a result, this hides effects related to how stringent the stopping criterion of the subproblem solver is, hence of how costly the subproblem solves are. For example, one can (usually) reduce the number of outer iterations of ARC by reducing $\theta$. As the subproblems of RTR and ARC are similar, we expect that (in principle) it should be possible to solve them equally well in about the same time.}
	\label{fig:XPouteriterations}
\end{figure}



We find that ARC's performance is in the same ballpark as RTR's, with the caveat that ARC's best performance requires tuning (choosing the right subproblem solver and $\theta$ for the problem class), whereas RTR is more robust. Since RTR's code has been refined over many years, we expect that further work can help reduce the gap. For example, we expect that the performance of ARC could be improved with further tuning of the regularization parameter update rule. In particular, we find that it is important to reduce regularization fast when close to convergence (but not earlier), to allow ARC to make steps similar to Newton's method. Work by~\citet{gould2012updatingregularization} could be a good starting point for such exploration.

%\TODO{Rewrite this---From the results, it is clear that ARC's performance is sensitive to the choice of parameter $\theta$, which suggests that it may be beneficial to explore alternative stopping criteria for the subproblem solver in practice. Furthermore, the well-known numerical instabilities of Lanczos lead to performance degradation when many Hessian-vector calls are needed to appropriately solve a subproblem. Effects seem particularly severe for the synchronization of rotations problem. Overall, either RTR or CG-HS, whose implementations have been refined over many years, has the upper hand in each test. It is nevertheless encouraging to see that this initial implementation of ARC has comparable performance on a number of tests. We hope that by refining the subproblem solver strategy (including stopping criteria) and the regularization parameter update rule, we can further improve the competitiveness of Riemannian ARC.}




%\section{Perspectives}



%\TODO{Re-evaluate this, starting from the existing write-up. Include something about curvature, and how the reasoning here also applies to gradient descent and trust regions \citep{boumal2016globalrates}. Or should I make that point in the intro and remove this section? } \TODO{If keep, might want to talk about the issue of globally defined retractions on incomplete manifolds.}

% \begin{footnotesize}

\paragraph{Acknowledgments}
We thank Pierre-Antoine Absil for numerous insightful and technical discussions, Stephen McKeown for directing us to, and guiding us through the relevance of Jacobi fields for our study of~\aref{assu:DRetr}, Chris Criscitiello and Eitan Levin for many discussions regarding regularity assumptions on manifolds, and Bryan Zhu for contributing his nonlinear CG subproblem solver to Manopt, and related discussions.


%\clearpage

\bibliographystyle{abbrvnat}
\bibliography{refs}
% \end{footnotesize}

\clearpage

\appendix 

\section{Proofs from Section~\ref{sec:arc}: mechanical lemmas} \label{app:arc}

Lemma~\ref{lem:termination} characterizes the conditions under which the subproblem solver is allowed to return $s_k = 0$ at iteration $k$.
\begin{proof}[Proof of Lemma~\ref{lem:termination}]
	By definition of the model $m_k$~\eqref{eq:mk} and by properties of retractions~\eqref{eq:nablahatfTgradf},
	\begin{align*}
	\nabla m_k(0) & = \nabla \hat f_k(0) = \grad f(x_k),
	\end{align*}
	where $\hat f_k = f \circ \Retr_{x_k}$.
	Thus, if $\grad f(x_k) = 0$, the first-order condition~\eqref{eq:firstorderprogress} allows $s_k = 0$. The other way around, if $s_k = 0$ is allowed, then $\|\nabla m_k(0)\| = 0$, so that $\grad f(x_k) = 0$.
	
	Now assume the second-order condition~\eqref{eq:secondorderprogress} is enforced. If $s_k = 0$ is allowed, then we already know that $\grad f(x_k) = 0$. Combined with~\eqref{eq:nablatwohatfHessfWzero}, we deduce that
	\begin{align*}
	\nabla^2 m_k(0) = \nabla^2 \hat f_k(0) = \Hess f(x_k),
	\end{align*}
	for any retraction. Then, condition~\eqref{eq:secondorderprogress} at $s_k = 0$ indicates $\nabla^2 m_k(0)$ is positive semidefinite, hence $\Hess f(x_k)$ is positive semidefinite. The other way around, if $\grad f(x_k) = 0$ and $\Hess f(x_k)$ is positive semidefinite, then $\nabla m_k(0) = \grad f(x_k)$ and $\nabla^2 m_k(0) = \Hess f(x_k)$, so that indeed $s_k = 0$ is allowed.
\end{proof}

The two supporting lemmas presented in Section~\ref{sec:arc} follow from the regularization parameter update mechanism of Algorithm~\ref{algo:ARC}. The standard proofs are not affected by the fact we here work on a manifold. We provide them for the sake of completeness.
\begin{proof}[Proof of Lemma~\ref{lem:varsigmamax}]
	 Using the definition of $\rho_k$~\eqref{eq:rhok}, $m_k(0) = f(x_k)$~\eqref{eq:mk} and $m_k(0) - m_k(s_k) \geq 0$ by condition~\eqref{eq:firstorderprogress}: %\TODO{This proof tacitly assumes nonzero steps, so it's good to have the algorithm terminate if that ever happens.}
	\begin{align*}
		1 - \rho_k & = 1 - \frac{f(x_k) - f(\Retr_{x_k}(s_k))}{m_k(0) - m_k(s_k) + \frac{\varsigma_k}{3} \|s_k\|^3} \leq  \frac{f(\Retr_{x_k}(s_k)) - m_k(s_k) + \frac{\varsigma_k}{3} \|s_k\|^3}{\frac{\varsigma_k}{3} \|s_k\|^3}.
	\end{align*}
	Owing to~\aref{assu:firstorderregularityscalar}, the numerator is upper bounded by $(L/6)\|s_k\|^3$. Hence, $1 - \rho_k \leq \frac{L}{2\varsigma_k}$. If $\varsigma_k \geq \frac{L}{2(1-\eta_2)}$, then $1-\rho_k \leq 1-\eta_2$ so that $\rho_k \geq \eta_2$, meaning step $k$ is very successful. The regularization mechanism~\eqref{eq:varsigmaupdate} then ensures $\varsigma_{k+1} \leq \varsigma_k$.
	Thus, $\varsigma_{k+1}$ may exceed $\varsigma_k$ only if $\varsigma_k < \frac{L}{2(1-\eta_2)}$, in which case it can grow at most to $\frac{L\gamma_3}{2(1-\eta_2)}$, but cannot grow beyond that level in later iterations.
%	
%	Thus, the only possibility for $\varsigma_k$ to exceed $\frac{3L}{2(1-\eta_2)}$ is if $\varsigma_0$ does so, or $\varsigma_{k-1} < \frac{3L}{2(1-\eta_2)}$ and step $k-1$ is unsuccessful, leading to an increase of the regularization parameter by at most a factor of $\gamma_3$.
\end{proof}

\begin{proof}[Proof of Lemma~\ref{lem:boundKsuccessfulsteps}]
	Partition iterations $0, \ldots, \bar k - 1$ into successful or very successful ($\calS_{\bar k}$) and unsuccessful ($\calU_{\bar k}$) ones.
	Following the update mechanism~\eqref{eq:varsigmaupdate}, for $k \in \calS_{\bar k}$, $\varsigma_{k+1} \geq \gamma_1 \varsigma_k$, while for $k \in \calU_{\bar k}$, $\varsigma_{k+1} \geq \gamma_2 \varsigma_k$. Thus, by induction, $\varsigma_{\bar k} \geq \varsigma_0 \gamma_1^{|\calS_{\bar k}|} \gamma_2^{|\calU_{\bar k}|}$. By assumption, $\varsigma_{\bar k} \leq \varsigma_{\max}$ so that
	\begin{align*}
		\log\left(\frac{\varsigma_{\max}}{\varsigma_0}\right) \geq |\calS_{\bar k}| \log(\gamma_1) + |\calU_{\bar k}| \log(\gamma_2) = |\calS_{\bar k}|\left[\log(\gamma_1) - \log(\gamma_2)\right] + \bar k \log(\gamma_2),
	\end{align*}
	where we also used $|\calS_{\bar k}| + |\calU_{\bar k}| = \bar k$. Isolating $\bar k$ using $\gamma_2 > 1 > \gamma_1$ %yields
%	\begin{align*}
%		\bar k & \leq \frac{1}{\log(\gamma_2)} \left[ |\calS_{\bar k}|\log\left(\frac{\gamma_2}{\gamma_1}\right) + \log\left(\frac{\varsigma_{\max}}{\varsigma_0}\right) \right],
%	\end{align*}
%	which
	allows to conclude.
\end{proof}



\section{Proofs from Section~\ref{sec:firstorderanalysisexp}: first-order analysis, exponentials} \label{app:firstorderexp}

Certain tools from Riemannian geometry are useful throughout the appendices---see for example \cite[pp59--67]{oneill}. To fix notation, let $\nabla$ denote the Riemannian connection on $\calM$ (not to be confused with $\nabla$ and $\nabla^2$ which denote gradient and Hessian of functions on linear spaces, such as pullbacks). With this notation, the Riemannian Hessian~\cite[Def.~5.5.1]{AMS08} is defined by $\Hess f = \nabla \grad f$. Furthermore, $\frac{\D}{\dt}$ denotes the covariant derivative of vector fields along curves on $\calM$, induced by $\nabla$. With this notation, given a smooth curve $c \colon \reals \to \calM$, the intrinsic acceleration is defined as $c''(t) = \frac{\D^2}{\dt^2} c(t)$. For example, for a Riemannian submanifold of a Euclidean space, $c''(t)$ is obtained by orthogonal projection of the classical acceleration of $c$ in the embedding space to the tangent space at $c(t)$. Geodesics are those curves which have zero intrinsic acceleration.

%\TODO{See notes Oct. 3, 2019, based on square (not star): could do the full proof for a general curve $c(t)$, and at the end show specifically which properties of geodesics lead to drastic simplifications. Could then refer to this proof in the body of the text to argue why it's not so easy to get a similar statement with a general retraction, thus justifying our proposal.}

We first state and prove a partial version of Proposition~\ref{prop:LipschitzHessianBounds} which applies for general retractions. Right after this, we prove Proposition~\ref{prop:LipschitzHessianBounds}. The purpose of this detour is to highlight how crucial properties of geodesics and of their interaction with parallel transports allow for the more direct guarantees of Section~\ref{sec:firstorderanalysisexp}. In turn, this serves as motivation for the developments in Section~\ref{sec:firstorderanalysisgeneral}.
\begin{proposition} \label{prop:LipschitzHessianBoundsretraction}
	Let $f \colon \calM \to \reals$ be twice differentiable on a Riemannian manifold $\calM$ equipped with a retraction $\Retr$. Given $(x, s) \in \T\calM$, assume there exists $L \geq 0$ such that, for all $t \in [0, 1]$,
	\begin{align*}
%		\left\| \left(P_{t s}^{-1} \circ \Hess f(c(t)) \circ T_{t s} - \Hess f(x)\right)\![s] \right\| & \leq L \cdot \ell(c|_{[0, t]}) \cdot \|s\|,
		\left\| P_{t s}^{-1}\!\left(\Hess f(c(t))[c'(t)]\right) - \Hess f(x)[s] \right\| & \leq L \|s\| \cdot \ell(c|_{[0, t]}),
	\end{align*}
	where $P_{ts}$ is parallel transport along $c(t) = \Retr_x(ts)$ from $c(0)$ to $c(t)$ (note the retraction instead of the exponential) %,
	%$T_{ts} = \D\Retr_x(ts)$,
	and $\ell(c|_{[0, t]}) = \int_{0}^{t} \|c'(\tau)\| \dtau$ is the length of $c$ restricted to the interval $[0, t]$. Then,
	\begin{align*}
		\left\| P_{s}^{-1}\grad f(\Retr_x(s)) - \grad f(x) - \Hess f(x)[s] \right\| & \leq L \|s\| \int_{0}^{1} \ell(c|_{[0, t]}) \,\dt.
	\end{align*}	
\end{proposition}
\begin{proof}
	Pick a basis $v_1, \ldots, v_d$ for $\T_x\calM$, and define the parallel vector fields $V_i(t) = P_{ts}(v_i)$ along $c(t)$. Since parallel transport is an isometry, $V_1(t), \ldots, V_d(t)$ form a basis for $\T_{c(t)}\calM$ for each $t \in [0, 1]$. As a result, we can express the gradient of $f$ along $c(t)$ in these bases,
	\begin{align}
		\grad f(c(t)) & = \sum_{i = 1}^{d} \alpha_i(t) V_i(t),
		\label{eq:gradfctexpansion}
	\end{align}
	with $\alpha_1(t), \ldots, \alpha_d(t)$ differentiable. Using properties of the Riemannian connection $\nabla$ and its associated covariant derivative $\Ddt$~\cite[pp59--67]{oneill}, we find on one hand that
	\begin{align*}
		\Ddt \grad f(c(t)) & = \nabla_{c'(t)} \grad f = \Hess f(c(t))[c'(t)],
	\end{align*}
	and on the other hand that
	\begin{align*}
		\Ddt \sum_{i = 1}^{d} \alpha_i(t) V_i(t) & = \sum_{i = 1}^{d} \alpha_i'(t) V_i(t) = P_{ts} \sum_{i = 1}^{d} \alpha_i'(t) v_i,
	\end{align*}
	where we used that $\Ddt V_i(t) = 0$, by definition of parallel transport.
	Furthermore,
	\begin{align*}
		c'(t) = \D\Retr_x(ts)[s] = T_{ts}(s),
	\end{align*}
	where $T_{ts} = \D\Retr_x(ts)$ is a linear operator from the tangent space at $x$ to the tangent space at $c(t)$---just like $P_{ts}$.
	Combining, we deduce that
	\begin{align*}
		\sum_{i = 1}^d \alpha_i'(t) v_i & = \left(P_{ts}^{-1} \circ \Hess f(c(t)) \circ T_{ts}\right)\![s].
	\end{align*}
	Going back to~\eqref{eq:gradfctexpansion}, we also see that
	\begin{align*}
		G(t) & \triangleq P_{ts}^{-1} \grad f(c(t)) = \sum_{i = 1}^{d} \alpha_i(t) v_i
	\end{align*}
	is a map from (a subset of) $\reals$ to $\T_x\calM$---two linear spaces---so that we can differentiate it in the usual way:
	\begin{align*}
		G'(t) = \sum_{i = 1}^{d} \alpha_i'(t) v_i.
	\end{align*}
	We conclude that
	\begin{align}
		G'(t) = \ddt\!\left[ P_{ts}^{-1} \grad f(c(t)) \right] & = \left(P_{ts}^{-1} \circ \Hess f(c(t)) \circ T_{ts}\right)\![s].
		\label{eq:ddtptsinvgradfct}
	\end{align}
	Since $G'$ is continuous,
	\begin{align*}
		P_{ts}^{-1}\grad f(c(t)) = G(t) & = G(0) + \int_{0}^{t} G'(\tau) \dtau \\
										& = \grad f(x) + \int_{0}^{t} \left(P_{\tau s}^{-1} \circ \Hess f(c(\tau)) \circ T_{\tau s}\right)\![s]  \dtau.
	\end{align*}
	Moving $\grad f(x)$ to the left-hand side and subtracting $\Hess f(x)[ts]$ on both sides, we find
	\begin{align*}
		P_{ts}^{-1}\grad f(c(t)) - \grad f(x) - \Hess f(x)[ts] & = \int_{0}^{t} \left(P_{\tau s}^{-1} \circ \Hess f(c(\tau)) \circ T_{\tau s} - \Hess f(x)\right)\![s]  \dtau.
	\end{align*}
	Using the main assumption on $\Hess f$ along $c$, it easily follows that
	\begin{align}
		\left\| P_{ts}^{-1}\grad f(c(t)) - \grad f(x) - \Hess f(x)[ts] \right\| & \leq \|s\| L \int_{0}^{t} \ell(c|_{[0, \tau]}) \dtau.
		\label{eq:step2normboundLipschitzHessiangeneral}
	\end{align}
	For $t = 1$, this is the announced inequality.
\end{proof}

\begin{proof}[Proof of Proposition~\ref{prop:LipschitzHessianBounds}]
	In this proposition we work with the exponential retraction, so that instead of a general retraction curve $c(t)$ we work along a geodesic $\gamma(t) = \Exp_x(ts)$. By definition, the velocity vector field $\gamma'(t)$ of a geodesic $\gamma(t)$ is parallel, meaning
	\begin{align}
		\gamma'(t) = P_{ts}(\gamma'(0)) = P_{ts}(s).
		\label{eq:gammaprimePt}
	\end{align}
	This elegant interplay of geodesics and parallel transport is crucial.
	In particular, 
	\begin{align*}
		 \ell(\gamma|_{[0, t]}) = \int_{0}^{t} \|\gamma'(\tau)\| \dtau = t \|s\|,
	\end{align*}
	and the condition in Proposition~\ref{prop:LipschitzHessianBoundsretraction} becomes
	\begin{align*}
		\left\| P_{t s}^{-1}\!\left(\Hess f(\gamma(t))[P_{ts}(s)]\right) - \Hess f(x)[s] \right\| & \leq t L \|s\|^2,
	\end{align*}
	which is indeed guaranteed by our own assumptions. We deduce that~\eqref{eq:step2normboundLipschitzHessiangeneral} holds:
	\begin{align}
		\left\| P_{ts}^{-1}\grad f(\gamma(t)) - \grad f(x) - \Hess f(x)[ts] \right\| & \leq \|s\| L \int_{0}^{t} \ell(\gamma|_{[0, \tau]}) \dtau = \frac{L}{2} \|s\|^2 t^2.
		\label{eq:step2normboundLipschitzHessian}
	\end{align}
	The relation~\eqref{eq:gammaprimePt} also yields the scalar inequality. Indeed, since $f \circ \gamma \colon [0, 1] \to \reals$ is continuously differentiable,
	\begin{align*}
	f(\Exp_x(s)) = f(\gamma(1)) & = f(\gamma(0)) + \int_{0}^{1} (f \circ \gamma)'(t) \dt \\
							& = f(x) + \int_{0}^{1} \inner{\grad f(\gamma(t))}{\gamma'(t)} \dt \\
							& = f(x) + \int_{0}^{1} \inner{P_{ts}^{-1}\grad f(\gamma(t))}{s} \dt,
	\end{align*}
	where on the last line we used~\eqref{eq:gammaprimePt} and the fact that $P_{ts}$ is an isometry.
	For a general retraction curve $c(t)$, instead of $s$ as the right-most term we would find $P_{ts}^{-1}(c'(t))$ which may vary with $t$: this would make the next step significantly more difficult.
	Move $f(x)$ to the left-hand side and subtract terms on both sides to get
	\begin{multline*}
		f(\Exp_x(s)) - f(x) - \inner{\grad f(x)}{s} - \frac{1}{2} \inner{s}{\Hess f(x)[s]} \\ = \int_{0}^{1} \inner{P_{ts}^{-1}\grad f(\gamma(t)) - \grad f(x) - \Hess f(x)[ts]}{s} \dt.
	\end{multline*}
	Using~\eqref{eq:step2normboundLipschitzHessian} and Cauchy--Schwarz, it follows immediately that
	\begin{align*}
	\left| f(\Exp_x(s)) - f(x) - \inner{s}{\grad f(x)} - \frac{1}{2} \inner{s}{\Hess f(x)[s]} \right| \leq \int_{0}^{1} \frac{L}{2} \|s\|^3 t^2 \dt = \frac{L}{6} \|s\|^3,
	\end{align*}
	as announced.
\end{proof}

Next, we provide an argument for the last claim in Theorem~\ref{thm:masterboundSLipschitzHessian}.
\begin{proof}[Proof of Theorem~\ref{thm:masterboundSLipschitzHessian}]
	We argue that $\lim_{k \to \infty} \|\grad f(x_k)\| = 0$. The first claim of the theorem states that, for every $\varepsilon > 0$, there is a finite number of successful steps $k$ such that $x_{k+1}$ has gradient larger than $\varepsilon$. Thus, for any $\varepsilon > 0$, there exists $K$: the last successful step such that $x_{K+1}$ has gradient larger than $\varepsilon$. Furthermore, there is a finite number of unsuccessful steps directly after $K+1$. Indeed, $\varsigma_{K+1} \geq \varsigma_{\min}$, and failures increase $\varsigma$ exponentially; additionally, $\varsigma$ cannot outgrow $\varsigma_{\max}$ by Lemma~\ref{lem:varsigmamax}. Thus, after a finite number of failures, a new success arises, necessarily producing an iterate with gradient norm at most $\varepsilon$ since $K$ was the last successful step to produce a larger gradient. By the same argument, all subsequent iterates have gradient norm at most $\varepsilon$. In other words: for any $\varepsilon > 0$, there exists $K'$ finite such that for all $k \geq K'$, $\|\grad f(x_k)\| \leq \varepsilon$, that is: $\lim_{k\to\infty} \|\grad f(x_k)\| = 0$.
\end{proof}



%\section{Proofs from Section~\ref{sec:firstorderanalysisgeneral}: first-order analysis, retractions} \label{app:firstordergeneral}

%\TODO{...}



\section{Proofs from Section~\ref{sec:secondorder}: second-order analysis} \label{app:secondorder}


\begin{proof}[Proof of Corollary~\ref{cor:mastercorollarysecond}]
	Consider these subsets of the set of successful iterations $\calS$: % we killed notation $\calS_{\bar k}$, but this use of $\calS$ is fine, because it's defined by the sentence.
	\begin{align*}
		\calS^1 & \defeq \{k \in \calS : \|\grad f(x_{k + 1})\| > \varepsilon_g\}, & & \text{ and } & \calS^2 & \defeq \{k \in \calS : \lambda_{\min}(\Hess f(x_{k})) < - \varepsilon_{{H}}\}.
	\end{align*}
	These sets are finite: for $K_1 = K_1(\varepsilon_g)$ as provided by either Theorem~\ref{thm:masterboundSLipschitzHessian} or Theorem~\ref{thm:masterboundSLipschitzHessianretraction}, and for $K_2 = K_2(\varepsilon_{{H}})$ as provided by Theorem~\ref{thm:masterboundSsecond}, we know that
	\begin{align*}
	|\calS^1| & \leq K_1, & & \textrm{ and } & |\calS^2| & \leq K_2.
	\end{align*}
	Note that successful steps are in one-to-one correspondence with the distinct points in the sequence of iterates $x_1, x_2, x_3, \ldots$\footnote{This is true because the cost function is strictly decreasing when successful, so that any $x_k$ can only be repeated in one contiguous subset of iterates. Hence, if $k$ is a successful iteration, match it to $x_{k+1}$ (this is why we omitted $x_0$ from the list.)} The first inequality states at most $K_1$ of the distinct points in that list have large gradient. The second inequality states at most $K_2$ of the distinct points in that same list have significantly negative Hessian eigenvalues. Thus, if more than $K_1+K_2+1$ distinct points appear among $x_0, x_1, \ldots, x_{\bar k}$ (note the $+1$ as we added $x_0$ to the list),
	% See notes NB Sep 10, 2018
	then at least one of these points has both a small gradient and an almost positive semidefinite Hessian.
	%	\[
	%	|\calS_{\bar k}^1| \leq K_1 \defeq \frac{3(f(x_0) - \flow)}{\eta_1 \varsigma_{\min}} \left(\frac{\frac{L}{2}+\theta+\varsigma_{\max}}{b}\right)^{1.5} \frac{1}{\varepsilon_g^{1.5}},
	%	\]
	%	\[|
	%	\calS_{\bar k}^2| \leq K_2  \triangleq \frac{3(f(x_0) - \flow)}{\eta_1 \varsigma_{\min}} (\theta + 2\varsigma_{\max})^3 \frac{1}{\varepsilon_{{H}}^3}.
	%	\]
	%
	In particular, as long as the number of successful iterations among $0, \ldots, \bar k-1$ exceeds $K_1 + K_2 + 1$ (strictly),
	there must exist $k \in \{0, \ldots, \bar k\}$ such that
	\begin{align*}
	\|\grad f(x_{k})\| & \leq \varepsilon_g & & \text{ and } & \lambda_{\min}(\Hess f(x_{k})) & \geq - \varepsilon_{{H}}.
	\end{align*} 
	Lemma~\ref{lem:boundKsuccessfulsteps} allows to conclude. 
	%
	%
	% \TODO{The argment is essentially this (to check with a skeptical mind): consider all the successful steps: they identify the unique points in the sequence of iterates (because the only way to reach a new point is via a successful step, and there can be no duplicates since the method descends.) Now, the first-order analysis says at most $K_1$ of these points can be excluded due to large gradient, and the second-order analysis says at most $K_2$ of these points can be excluded due to too-negative Hessian. It doesn't matter that first-order speaks of the gradient at the next step and that second order speaks of the Hessian at the current step: it's just about how many unique points can be bad. Worst-case scenario, $K_1 + K_2$ unique points are ruled out. So, if we have strictly more than $K_1 + K_2$ successful steps, at least one of the produced points is good on both counts. Then we round up with Lemma~\ref{lem:boundKsuccessfulsteps}.}
\end{proof}


\section{Proofs from Section~\ref{sec:regularity}: regularity assumptions} \label{app:regularity}
\begin{proof}[Proof of Lemma~\ref{lem:lipschitzhessianbasic}]
%	By the fundamental theorem of calculus applied to $\hat f$ and $\nabla \hat f$,
	Since $\hat f$ is a real function on a linear space, standard calculus applies:
	\begin{align*}
		\hat f(s) - \left[ \hat f(0) + \innersmall{s}{\nabla \hat f(0)} + \frac{1}{2} \innersmall{s}{\nabla^2 \hat f(0)[s]} \right] & = \int_{0}^{1}\!\int_{0}^{1} \tauort_1 \inner{\left[ \nabla^2 \hat f(\tauort_1\tauort_2 s) - \nabla^2 \hat f(0) \right]\![s]}{s}  \mathrm{d}\tauort_1 \mathrm{d}\tauort_2, \\
		\nabla \hat f(s) - \left[\nabla \hat f(0) + \nabla^2 \hat f(0)[s]\right] & = \int_{0}^{1} \left[ \nabla^2 \hat f(\tauort s) - \nabla^2 \hat f(0) \right]\![s] \, \mathrm{d}\tauort.
	\end{align*}
	Taking norms on both sides, by a triangular inequality to pass the norm through the integral and integrating respectively $\tauort_1^2 \tauort_2^{}$ and $\tauort$, we find using our main assumption~\eqref{eq:lipschitzhessianbasic} that
	\begin{align*}
		\left| \hat f(s) - \left[ \hat f(0) + \innersmall{s}{\nabla \hat f(0)} + \frac{1}{2} \innersmall{s}{\nabla^2 \hat f(0)[s]} \right] \right| & \leq \frac{1}{6} L \|s\|^3, \textrm{ and } \\
		\left\| \nabla \hat f(s) - \left[\nabla \hat f(0) + \nabla^2 \hat f(0)[s]\right] \right\| & \leq \frac{1}{2} L \|s\|^2. \qedhere
	\end{align*}
\end{proof}

\begin{proof}[Proof of Lemma~\ref{lem:derivativespullback}]
For an arbitrary $\dot s \in \T_x\calM$, consider the curve $c(t) = \Retr_x(s + t\dot s)$, and let $g = f \circ c \colon \reals \to \reals$. We compute the derivatives of $g$ in two different ways. On the one hand, $g(t) = \hat f(s + t\dot s)$ so that
	\begin{align*}
		g'(t) & = \D\hat f(s+t\dot s)[\dot s] = \innersmall{\nabla \hat f(s+t\dot s)}{\dot s}, \\
		g''(t) & = \inner{ \ddt \nabla \hat f(s+t\dot s)}{\dot s} = \innersmall{\nabla^2 \hat f(s+t\dot s)[\dot s]}{\dot s}.
	\end{align*}
	On the other hand, $g(t) = f(c(t))$ so that, using properties of $\frac{\D}{\mathrm{d}t}$~\cite[pp59--67]{oneill}:
	\begin{align*}
		g'(t)  & = \D f(c(t))[c'(t)] = \innersmall{\grad f(c(t))}{c'(t)}, \\
		g''(t) & = \ddt \inner{(\grad f \circ c)(t))}{c'(t)} \\
		       & = \inner{\nabla_{c'(t)} \grad f}{c'(t)} + \inner{(\grad f \circ c)(t))}{\frac{\D}{\dt} c'(t)} \\
		 	   & = \inner{\Hess f(c(t))[c'(t)]}{c'(t)} + \inner{\grad f(c(t))}{c''(t)}.
	\end{align*}
%	In the above expression, $c''(t)$ is the ``intrinsic'' acceleration of the curve $c$ at time $t$; it is a tangent vector at $c(t)$. For Riemannian submanifolds embedded in a Euclidean space, $c''$ is simply the tangent component of the classical acceleration of $c$ in the embedding space.
%	
	Equating the different identities for $g'(t)$ and $g''(t)$ at $t = 0$ while using $c'(0) = T_s \dot s$, we find for all $\dot s \in \T_x\calM$:
	\begin{align*}
		\innersmall{\nabla \hat f(s)}{\dot s} & = \inner{\grad f(\Retr_x(s))}{T_s \dot s}, \\
		\innersmall{\nabla^2 \hat f(s)[\dot s]}{\dot s} & = \inner{\Hess f(\Retr_x(s))[T_s \dot s]}{T_s \dot s} + \inner{\grad f(\Retr_x(s))}{c''(0)}.
	\end{align*}
	The last term, $\inner{\grad f(\Retr_x(s))}{c''(0)}$, is seen to be the difference of two quadratic forms in $\dot s$, so that it is itself a quadratic form in $\dot s$. This justifies the definition of $W_s$ through polarization.
	The announced identities follow by identification.
\end{proof}


%\TODO{We should be able to do the same proof for the polar retraction on Stiefel, because it has a quite explicit (and similar) expression: $\Retr_X(S) = (X + S)(I_r + S\transpose S)^{-1/2}$. --- See notes June 22 and 25, 2018 and Matlab code on GIT: we have a formula for the derivative of polar factor on Stiefel but it's not so nice because it involves the SVD. On the other hand, the derivative is quite simply on Grassmann because the Sylvester equation is trivial in that case. So perhaps look into it.
\begin{proof}[Proof of Proposition~\ref{prop:sphereretrnice}]
%	We start with the first property.
	With $\Retr_x(s) = \frac{x+s}{\sqrt{1+\|s\|^2}}$, it is easy to derive
	\begin{align}
		T_s \dot s \triangleq \D\Retr_x(s)[\dot s] & = \left[ \frac{1}{\sqrt{1+\|s\|^2}} I_n - \frac{1}{\sqrt{1+\|s\|^2}^3} (x+s)s\transpose \right] \dot s \nonumber \\
		& = \frac{1}{\sqrt{1+\|s\|^2}} \left[ I_n - \Retr_x(s) \Retr_x(s)\transpose \right] \dot s,
		\label{eq:DretrSphere}
	\end{align}
	where we used $x\transpose \dot s = 0$ in between the two steps to replace $s\transpose$ with $(x+s)\transpose$.
	The matrix between brackets is the orthogonal projector from $\Rn$ to $\T_{\Retr_x(s)}\calM$. Thus, its singular values are upper bounded by 1. Since $T_s$ is an operator on $\T_x\calM \subset \Rn$, 
	\begin{align*}
		\opnorm{T_s} & \leq \frac{1}{\sqrt{1+\|s\|^2}} \leq 1.
	\end{align*}
	This secures the first property with $c_1 = 1$.
	%The right hand side is only a function of $\|s\| \geq 0$. It attains its maximum for $\|s\| = 1$, giving a bound $\opnorm{T_s} \leq c_2 \triangleq \frac{3}{2}$.
	
	For the second property, consider $U(t) = T_{ts} \dot s$ and
	\begin{align*}
		 U'(t) \triangleq \frac{\D}{\mathrm{d}t} U(t) = \Proj_{c(t)} \ddt U(t),
	\end{align*}
	where $\Proj_y(v) = v - y (y\transpose v)$ is the orthogonal projector to $\T_y\calM$ and $c(t) = \Retr_x(ts)$. Define $g(t) = \frac{1}{\sqrt{1+t^2\|s\|^2}}$. Then, from~\eqref{eq:DretrSphere}, we have
	\begin{align}
		U(t) & = \left[ g(t) I_n - t g(t)^3 (x+ts)s\transpose \right] \dot s.
	\end{align}
	This is easily differentiated in the embedding space $\Rn$:
	\begin{align*}
		\ddt U(t) & = \left[ g'(t) I_n - (t g(t)^3)' (x+ts)s\transpose - t g(t)^3 ss\transpose \right] \dot s.
	\end{align*}
	The projection at $c(t)$ zeros out the middle term, as it is parallel to $x+ts$. This offers a simple expression for $U'(t)$, where in the last equality we use $g'(t) = -t g(t)^3 \|s\|^2$:
	\begin{align*}
		U'(t) & = \Proj_{c(t)}\left( \left[ g'(t) I_n - t g(t)^3 ss\transpose \right] \dot s \right) = -t g(t)^3 \cdot  \Proj_{c(t)}\left( \left[ \|s\|^2 I_n + ss\transpose \right] \dot s \right).
	\end{align*}
	The norm can only decrease after projection, so that, for $t \in [0, 1]$, % can do global t if careful with sign
	\begin{align*}
		\|U'(t)\| & \leq 2 t g(t)^3 \|s\|^2 \|\dot s\|.
	\end{align*}
	Let $h(t) = 2 t g(t)^3 \|s\|^2 = \frac{2t\|s\|^2}{(1+t^2 \|s\|^2)^{1.5}}$. For $s = 0$, $h$ is identically zero. Otherwise, $h$ attains its maximum $h\left(t = \frac{1}{\sqrt{2} \|s\|}\right) = \frac{4 \sqrt{3}}{9} \|s\|$. It follows that $\|U'(t)\| \leq c_2 \|s\| \|\dot s\|$ for all $t \in [0, 1]$ with $c_2 = \frac{4\sqrt{3}}{9}$.
	
	Finally, we establish the last property. Given $s, \dot s \in \T_x\calM$, consider $c(t) = \Retr_x(s + t\dot s)$. Simple calculations yield:
	\begin{align}
		c'(t) & = \ddt c(t) = \frac{1}{\sqrt{1 + \|s + t\dot s\|^2}} \left[ \dot s - \inner{\dot s}{c(t)} c(t) \right] = \frac{1}{\sqrt{1 + \|s + t\dot s\|^2}} \Proj_{c(t)} \dot s.
		\label{eq:ddtRsphere}
	\end{align}
	This is indeed in the tangent space at $c(t)$. The classical derivative of $c'(t)$ is given by
	\begin{align*}
		\ddt c'(t) & = - \frac{1}{\sqrt{1 + \|s + t\dot s\|^2}} \left[ \inner{\dot s}{c'(t)} c(t) + \inner{\dot s}{c(t)} c'(t) + \frac{\inner{s + t\dot s}{\dot s}}{1 + \|s + t \dot s\|^2} \Proj_{c(t)} \dot s \right] \\
			& = - \frac{1}{\sqrt{1 + \|s + t\dot s\|^2}} \left[ \inner{\dot s}{c'(t)} c(t) + 2\frac{\inner{s + t\dot s}{\dot s}}{1 + \|s + t \dot s\|^2} \Proj_{c(t)} \dot s \right],
	\end{align*}
	where we used~\eqref{eq:ddtRsphere} and orthogonality of $x$ and $\dot s$ in $\inner{c(t)}{\dot s} = \frac{1}{\sqrt{1+\|s+t\dot s\|^2}} \inner{x + s + t \dot s}{\dot s}$.
	The acceleration of $c$ is $c''(t) = \frac{\D}{\mathrm{d}t} c'(t) = \Proj_{c(t)}\left( \ddt c'(t) \right)$. The first term vanishes after projection, while the second term is unchanged. Overall,
	\begin{align}
		c''(t) & = -\frac{2\inner{s + t\dot s}{\dot s}}{\sqrt{1 + \|s + t \dot s\|^2}^3} \Proj_{c(t)} \dot s =  -\frac{2\inner{c(t)}{\dot s}}{1 + \|s + t \dot s\|^2} \Proj_{c(t)} \dot s.
		% See D2Retr.m and notes NB April 2018 -- this is correct.
	\end{align}
%	Hence, acceleration is bounded as $\|c''(t)\| \leq 2\|\dot s\|^2$ for all $t$: the property holds with $c_1 = 2$. \TODO{NOT ENOUGH -- need a norm of s in there; perhaps only true at $t = 0$.}
	In particular, $c''(0) = -2 \frac{\inner{s}{\dot s}}{\sqrt{1+\|s\|^2}^3} \Proj_{c(0)} \dot s$, so that $\|c''(0)\| \leq 2 \min(\|s\|, 0.4) \|\dot s\|^2$ and the property holds with $c_3 = 2$.
	% https://www.wolframalpha.com/input/?i=derive+f(x)+%3D+x%2Fsqrt(1%2Bx%5E2)%5E3 -- maximum of x/sqrt(1+x^2)^3 is attained at x = 1/sqrt(2) and equals 0.3849...
	(Peculiarly, if $s$ and $\dot s$ are orthogonal, $c''(0) = 0$.)
\end{proof}

In order to prove Theorem~\ref{thm:regularitycompact}, we introduce two supporting lemmas (needed only for the case where $\calM$ is not compact) and one key lemma. The first lemma below is similar in spirit to~\citep[Lem.~2.2]{cartis2011adaptivecubicPartI}.
\begin{lemma}\label{lem:boundessteps}
	Let $f \colon \calM \to \reals$ be twice continuously differentiable. % Fine to talk about 2nd order diff. instead of smooth here because we only involve differentials of the pullback.
	Let $\{ (x_0, s_0), (x_1, s_1), \ldots \}$ be the points and steps generated by Algorithm~\ref{algo:ARC}.
	%Under~\aref{assu:firstorderprogress},
	Each step has norm bounded as:
	\begin{align}
		\|s_k\| & \leq \sqrt{\frac{3 \|\nabla \hat f_k(0)\|}{\varsigma_{\min}}} + \frac{3}{2\varsigma_{\min}} \max\left(0, -\lambdamin(\nabla^2 \hat f_k(0))\right),
	\end{align}
	where $\hat f_k = f \circ \Retr_{x_k}$ is the pullback, as in~\eqref{eq:fhatk}.
\end{lemma}
\begin{proof}
%	Under~\aref{assu:firstorderprogress},
	Owing to the first-order progress condition~\eqref{eq:firstorderprogress}, using Cauchy--Schwarz and the fact that $\varsigma_k \geq \varsigma_{\min}$ for all $k$ by design of the algorithm, we find
	\begin{align*}
		\varsigma_{\min} \|s_k\|^3 \leq \varsigma_{k} \|s_k\|^3 & \leq -3\inner{s_k}{\nabla \hat f_k(0) + \frac{1}{2}\nabla^2 \hat f_k(0)[s_k]} \\
			& \leq 3 \|s_k\| \left( \|\nabla \hat f_k(0)\| + \frac{1}{2} \max\left(0, -\lambdamin(\nabla^2 \hat f_k(0))\right) \|s_k\| \right).
	\end{align*}
	This defines a quadratic inequality in $\|s_k\|$:
	\begin{align*}
		\varsigma_{\min} \|s_k\|^2 - h_k \|s_k\| - g_k \leq 0,
	\end{align*}
	where to simplify notation we let $h_k = \frac{3}{2} \max(0, -\lambdamin(\nabla^2 \hat f_k(0)))$ and $g_k = 3 \|\nabla \hat f_k(0)\|$. Since $\|s_k\|$ must lie between the two roots of this quadratic, we know in particular that
	\begin{align*}
		\|s_k\| & \leq \frac{h_k + \sqrt{h_k^2 + 4\varsigma_{\min}g_k}}{2\varsigma_{\min}} \leq \frac{h_k + \sqrt{\varsigma_{\min} g_k}}{\varsigma_{\min}},
	\end{align*}
	where in the last step we used $\sqrt{u+v} \leq \sqrt{u} + \sqrt{v}$ for any $u, v \geq 0$.
\end{proof}
\begin{lemma}\label{lem:calNsubsetcompact}
	Let $f \colon \calM \to \reals$ be twice continuously differentiable. % Fine to require just what is needed for the helper lemma: {lem:boundessteps} ; this is the only reason we need to assume anything about $f$ here.
	Let $\{ (x_0, s_0), (x_1, s_1), \ldots \}$ be the points and steps generated by Algorithm~\ref{algo:ARC}.
	Consider the following subset of $\calM$, obtained by collecting all curves generated by retracted steps (both accepted and rejected):
	\begin{align}
		\calN & = \bigcup_{k} \left\{ \Retr_{x_k}(ts_k) : t \in [0, 1] \right\}.
		\label{eq:calNretractedsteps}
	\end{align}
	If the sequence $\{x_0, x_1, x_2, \ldots\}$ remains in a compact subset of $\calM$, then $\calN$ is included in a compact subset of $\calM$.
\end{lemma}
\begin{proof}
	If $\calM$ is compact, the claim is clear since $\calN \subseteq \calM$. Otherwise, we use Lemma~\ref{lem:boundessteps}. Specifically, considering the upper bound in that lemma, define
	\begin{align*}
		\alpha(x) & = \sqrt{\frac{3 \|\nabla \hat f_x(0)\|}{\varsigma_{\min}}} + \frac{3}{2\varsigma_{\min}} \max\left(0, -\lambdamin(\nabla^2 \hat f_x(0))\right),
	\end{align*}
	where $\hat f_x = f \circ \Retr_x$. This is a continuous function of $x$, and $\|s_k\| \leq \alpha(x_k)$. Since by assumption $\{x_0, x_1, \ldots\} \subseteq \calK$ with $\calK$ compact, we find that
	\begin{align*}
		\forall k, \quad \|s_k\| \leq \sup_{k'} \alpha(x_{k'}) \leq \max_{x \in \calK} \alpha(x) \triangleq r,
	\end{align*}
	where $r$ is a finite number. Consider the following subset of the tangent bundle $\T\calM$:
	\begin{align*}
		\calK' & = \{ (x, s) \in \T\calM : x \in \calK, \|s\|_x \leq r \}.
	\end{align*}
	Since $\calK$ is compact, $\calK'$ is compact. % See notes early January 2019, NB37
	Furthermore, since the retraction is a continuous map, $\Retr(\calK')$ is compact, and it contains $\calN$.
\end{proof}
\begin{lemma} \label{lem:coreOfThmAboutA2}
	Let $f \colon \calM \to \reals$ be three times continuously differentiable, and consider the points and steps $\{(x_0, s_0), (x_1, s_1), \ldots \}$ generated by Algorithm~\ref{algo:ARC}.
	Assume the retraction is second-order nice on this set (see Definition~\ref{def:retrscndnice}).
	If the set $\calN$ as defined by~\eqref{eq:calNretractedsteps} is contained in a compact set $\calK$, then \aref{assu:firstorderregularityscalar} and~\aref{assu:firstorderregularityvectorretraction} are satisfied.
%	\TODO{Remove assumption that $\calK$ is compact (change name?) and instead request that $f$ have appropriate Lipschitz behavior ``on that set''? What does it mean to be Lipschitz on the set though? It's easy to ask for bounded derivatives on the set...}
\end{lemma}
\begin{proof}
	
%	\TODO{[NO, NOT SO EASY. DROP THIS.] I believe this could be rephrased in terms of Lipschitz gradient and Hessian, which could be acceptable. Then there are these extra assumptions on the retraction (which are not mild...).And there's still the gradient term due to $W_s$, which... not sure how it'd play out as we aim for the scalar inequality, in particular. Might even be able to get a separate claim for scalar and vector inequalities: we'd really have one (weak) statement for the scalar inequality, and two statements for the vector inequality: a weak one, and one that involves a $q$ function. Hmm, no, because for the $q$ part we had this integral of the gradient, so that didn't quite work as intended: see notes NB40, Oct 1, 2019.}
	
%	Let $\{ (x_0, s_0), (x_1, s_1), \ldots \}$ be the sequence of points and steps generated by Algorithm~\ref{algo:ARC}.
	For some $k$ and $\bar t \in [0, 1]$, let $(x, s) = (x_k, \bar t s_k)$ and define the pullback $\hat f = f \circ \Retr_x$. Notice in particular that $\Retr_x(s) \in \calN \subseteq \calK$.
	Combine the expression for the Hessian of the pullback~\eqref{eq:hessianpullback} with~\eqref{eq:lipschitzhessianbasic} to get:
	\begin{align*}
	\left\| \nabla^2 \hat f(s) - \nabla^2 \hat f(0) \right\|_{\mathrm{op}} & \leq \left\| T_{s}^* \circ \Hess f(\Retr_x(s)) \circ T_{s} - \Hess f(x) \right\|_{\mathrm{op}} + \left\|  W_{s} - W_0 \right\|_{\mathrm{op}}.
	\end{align*}
	By definition of $W_s$~\eqref{eq:Whessianpullback}, using the third condition on the retraction, we find that $W_0 = 0$ and
	\begin{align*}
	\|W_s\|_\mathrm{op} & = \max_{\substack{\dot s \in \T_x\calM \\ \|\dot s\| \leq 1}} \left| \inner{W_s[\dot s]}{\dot s} \right| \leq \|\grad f(\Retr_x(s))\| \cdot \max_{\substack{\dot s \in \T_x\calM \\ \|\dot s\| \leq 1}} \|c''(0)\| \leq c_3 G \|s\|,
	\end{align*}
	where $G = \max_{y \in \calK} \|\grad f(y)\|$ is finite by compactness of $\calK$ and continuity of the gradient norm.
	% The argument goes: sup_{N} f <= sup_{closure N} f = max_{closure N} f since closure of N is compact.
	Thus, it remains to show that
	\begin{align*}
	\left\| T_s^* \circ \Hess f(\Retr_x(s)) \circ T_s - \Hess f(x) \right\|_{\mathrm{op}} & \leq c' \|s\|
	\end{align*}
	for some constant $c'$. For an arbitrary $\dot s \in \T_x\calM$, owing to differentiability properties of $f$, %it holds that
	\begin{align}
		\inner{\big[T_s^* \circ \Hess f(\Retr_x(s)) \circ T_s - \Hess f(x)\big][\dot s]}{\dot s} & = \int_{0}^{1} \frac{\mathrm{d}}{\mathrm{d}\tauort}  \inner{T_{\tauort s}^* \circ \Hess f(\Retr_x(\tauort s)) \circ T_{\tauort s} [\dot s]}{\dot s}  \mathrm{d}\tauort.
		\label{eq:integralddth}
	\end{align}
	We aim to upper bound the above by $c' \|s\| \|\dot s\|^2$.
	Consider the curve $c(\tauort) = \Retr_{x}(\tauort s)$ and a tangent vector field $U(\tauort) = T_{\tauort s} \dot s$ along $c$. Then, define
	\begin{align*}
	h(\tauort) & = \inner{T_{\tauort s}^* \circ \Hess f(c(\tauort)) \circ T_{\tauort s} [\dot s]}{\dot s} \\
	& = \inner{\Hess f(c(\tauort))[T_{\tauort s} \dot s]}{T_{\tauort s} \dot s} \\
	& = \inner{\Hess f(c(\tauort))[U(\tauort)]}{U(\tauort)}.
	\end{align*}
	The integrand in~\eqref{eq:integralddth} is the derivative of the real function $h$:
	\begin{align*}
	h'(\tauort) & = \frac{\mathrm{d}}{\mathrm{d}\tauort} \inner{\Hess f(c(\tauort))[U(\tauort)]}{U(\tauort)} \\
	& = \inner{\frac{\D}{\mathrm{d}\tauort}\Big[\Hess f(c(\tauort))[U(\tauort)]\Big]}{U(\tauort)} + \inner{\Hess f(c(\tauort))[U(\tauort)]}{\frac{\D}{\mathrm{d}\tauort} U(\tauort)} \\
	& = \inner{\left( \nabla_{c'(\tauort)} \Hess f \right)[U(\tauort)]}{U(\tauort)} + 2\inner{\Hess f(c(\tauort))[U(\tauort)]}{U'(\tauort)},
	\end{align*}
	% https://en.wikipedia.org/wiki/Second_covariant_derivative -- and see the general notion of covariant derivative of tensor fields; see also notes in this project, outside the paper, to justify the derivative of the tensor field along a curve.
	where $U'(\tauort) \triangleq \frac{\D}{\mathrm{d}\tauort} U(\tauort)$ and we used that the Hessian is symmetric. Here, $\nabla_{c'(\tauort)} \Hess f$ is the Levi--Civita derivative of the Hessian tensor field at $c(\tauort)$ along $c'(\tauort)$---see~\cite[Def.~4.5.7, p102]{dC92}
	% do Carmo is really nicer than O'Neill here (O'Neill would be Prop. 2.13, but it's not explicitly with Levi--Civita yet, and it's only in tensor form (output is a scalar).
	for the notion of derivative of a tensor field. For every $\tauort$, the latter is a symmetric linear operator on the tangent space at $c(\tauort)$. By Cauchy--Schwarz,
	\begin{align*}
	|h'(\tauort)| & \leq \|\nabla_{c'(\tauort)} \Hess f\|_{\mathrm{op}} \|U(\tauort)\|^2 + 2 \|\Hess f(c(\tauort))\|_{\mathrm{op}} \|U(\tauort)\| \|U'(\tauort)\|.
	\end{align*}
	By compactness of $\calK$ and continuity of the Hessian, we can define
	\begin{align*}
		H & = \max_{y \in \calK} \|\Hess f(y)\|_{\mathrm{op}}.
	\end{align*}
	By linearity of the connection $\nabla$, if $c'(\tauort) \neq 0$,
	\begin{align*}
	\nabla_{c'(\tauort)} \Hess f & = \|c'(\tauort)\| \cdot \nabla_{\frac{c'(\tauort)}{\|c'(\tauort)\|}} \Hess f.
	\end{align*}
	Furthermore, $c'(\tauort) = T_{\tauort s} s$ has norm bounded by the first assumption on the retraction: $\|c'(\tauort)\| \leq c_1 \|s\|$. Thus, in all cases, by compactness of $\calK$ and continuity of the function $v \to \nabla_v \Hess f$ on the tangent bundle $\T\calM$, there is a finite $J$ as follows: % \TODO{Here too we use compactness of a set of balls in TM with compact base: see notes early January 2019, NB37}
	%\TODO{Should formally check that this is still a compact set in some appropriate topology where the function is continuous.}
	% For connected manifolds, perhaps a way to check this is: TM is a Riemannian manifold, with a certain metric, and the induced topology (if want to think of it as the topology of a metric space, need the manifold to be connnected) (see https://math.stackexchange.com/questions/980452/topologies-in-a-riemannian-manifold), so the notion of compactness should be "intuitive"; but even so, it'd be nice to write down a proper proof.
	% See also this question: https://math.stackexchange.com/questions/1177264/compactness-of-a-subset-of-a-tangent-bundle
	\begin{align*}
		\|\nabla_{c'(\tauort)} \Hess f \|_{\mathrm{op}} & \leq c_1 \|s\| \cdot \overbrace{\max_{\substack{y \in \calK, v \in \T_y\calM \\ \|v\| \leq 1}} \| \nabla_v \Hess f \|_{\mathrm{op}}}^{J}.
	\end{align*}
	Of course, $\|U(\tauort)\| \leq c_1 \|\dot s\|$. Finally, we bound $\|U'(\tauort)\|$ using the second property of the retraction: $\|U'(\tauort)\| \leq c_2\|s\| \|\dot s\|$. Collecting what we learned about $|h'(t)|$ and injecting in~\eqref{eq:integralddth},
	\begin{align*}
		\left|\inner{\big[T_s^* \circ \Hess f(\Retr_x(s)) \circ T_s - \Hess f(x)\big][\dot s]}{\dot s}\right| & \leq \int_{0}^{1} |h'(\tauort)| \mathrm{d}\tauort \leq \left[ c_1^3 J + 2 c_1 c_2 H \right] \|s\| \|\dot s\|^2.
	\end{align*}
	Finally, it follows from Lemma~\ref{lem:lipschitzhessianbasic} that~\aref{assu:firstorderregularityscalar} and~\aref{assu:firstorderregularityvectorretraction} hold with $L = L' = c_3G + 2 c_1 c_2 H + c_1^3 J$ and $q \equiv 0$. We note in closing that the constants $G, H, J$ can be related to the Lipschitz properties of $f$, $\grad f$ and $\Hess f$, respectively.
\end{proof}
The theorem we wanted to prove now follows as a direct corollary.
\begin{proof}[Proof of Theorem~\ref{thm:regularitycompact}]
	For the main result, simply combine Lemmas~\ref{lem:calNsubsetcompact} and~\ref{lem:coreOfThmAboutA2}. To support the closing statement, it is sufficient to verify that Algorithm~\ref{algo:ARC} is a descent method
	owing to the step acceptance mechanism and the first part of condition~\eqref{eq:firstorderprogress}.
	%owing to~\eqref{eq:firstorderprogress}, that is, $f(x_{k+1}) \leq f(x_k)$ for all $k$. This is argued at the beginning of the proof of Theorem~\ref{thm:masterboundS}.
\end{proof}

\section{Proofs from Section~\ref{sec:Dretr}: differential of retraction} \label{sec:proofsDretr}

\subsection*{Stiefel manifold}

Proposition~\ref{prop:stiefelmastercorollary} regarding the Stiefel manifold is a corollary of the following statement.
\begin{lemma} \label{lemma:stiefelbound}
	For the Stiefel manifold $\calM = \St(n,p)$ with the Q-factor retraction $\Retr$, for all $X \in \calM$ and $S \in \T_X\calM$,
	\begin{align*}
	\sigmamin(\D\Retr_X(S)) \geq 1 - 3\|S\|_{\mathrm{F}} - \frac{1}{2}\|S\|_{\mathrm{F}}^2,
	\end{align*}
	where $\|\cdot\|_{\mathrm{F}}$ denotes the Frobenius norm. Moreover, for the special case $p = 1$ (the unit sphere in $\Rn$), the retraction reduces to $\Retr_x(s) = \frac{x+s}{\|x+s\|}$ and we have for all $x \in \calM, s \in \T_x\calM$:
	\begin{align*}
	\sigmamin(\D\Retr_x(s)) = \frac{1}{1 + \|s\|^2}.
	\end{align*}
\end{lemma}
\begin{proof}
Let $X \in \St(n, p)$ and $S \in \T_X \St(n, p) = \{ \dot X \in \Rnp : \dot X\transpose  X + X\transpose \dot X = 0 \}$ be fixed. Define $Q,R$ as the thin $QR$-decomposition of $X + S$, that is, $Q$ is an $n \times p$ matrix with orthonormal columns and $R$ is a $p \times p$ upper triangular matrix with positive diagonal entries such that $X+S = QR$: this decomposition exists and is unique since $X+S$ has full column rank, as shown below~\eqref{eq:QRfull}. By definition, we have that $\Retr_X(S) = Q$. 

For a matrix $M$, define $\mathrm{tril}(M)$ as the lower triangular portion of the matrix $M$, that is, $\mathrm{tril}(M)_{ij} = M_{ij}$ if $i \geq j$ and 0 otherwise. Further define $\rho_{\mathrm{skew}}(M)$ as  
\[\rho_{\mathrm{skew}}(M) \defeq \mathrm{tril}(M) - \mathrm{tril}(M)\transpose.\]
As derived in~\citep[Ex.~8.1.5]{AMS08} (see also the erratum for the reference) we have a formula for the directional derivative of the retraction along any $Z \in \T_X \St(n, p)$:
\begin{align}
	\D\Retr_X(S)[Z] & = Q\rho_{\mathrm{skew}}(Q\transpose Z R^{-1}) + (I-QQ\transpose) Z R^{-1}.
	\label{eq:DRetrQR}
\end{align}
We first confirm that $R$ is always invertible. To see this, note that $S$ being tangent at $X$ means $S\transpose X + X\transpose S = 0$ and therefore
\begin{align}
	R\transpose R & = \underbrace{(X+S)\transpose (X + S)}_{\textrm{start reading here}} = X\transpose X + X\transpose S + S\transpose X + S\transpose S = I_p + S\transpose S,
	\label{eq:QRfull}
\end{align}
which shows $R$ is invertible. Moreover the above expression also implies that:
\begin{align*}
	\sigma_k(R) = \sigma_k(X + S) = \sqrt{\lambda_k((X+S)\transpose (X+S))} = \sqrt{1 + \lambda_k(S\transpose S)} = \sqrt{1 + \sigma_k(S)^2},
\end{align*}
where $\sigma_k(M)$ represents the $k$th singular value of $M$ and $\lambda_k$ likewise extracts the $k$th eigenvalue (in decreasing order for symmetric matrices).
In particular we have that
\begin{align}
	\sigmamin(R^{-1}) & = \frac{1}{\sqrt{1 + \sigmamax(S)^2}} \geq \frac{1}{\sqrt{1 + \|S\|_{\mathrm{F}}^2}} \geq 1 - \frac{1}{2}\|S\|_{\mathrm{F}}^2, \nonumber\\
	\sigmamax(R^{-1}) & = \frac{1}{\sqrt{1 + \sigmamin(S)^2}} \leq 1. \label{eqn:sigmamaxbound}
\end{align}
Further note that since $QR = X + S$, we have that $Q = (X+S)R^{-1}$ and therefore
\begin{align*}
	Q\transpose Z R^{-1} & = (R^{-1})\transpose (X+S)\transpose Z R^{-1} \nonumber\\
	& = (R^{-1})\transpose X\transpose Z R^{-1} + (R^{-1})\transpose S\transpose Z R^{-1}.
\end{align*}
The first term above is always skew-symmetric since $Z$ is tangent at $X$, so that $X\transpose Z + Z\transpose X = 0$. Furthermore, for any skew-symmetric matrix $M$, $\rho_{\mathrm{skew}}(M) = M$. Therefore, using~\eqref{eq:DRetrQR},
\begin{align}
	\D\Retr_X(S)[Z] & = Q\rho_{\mathrm{skew}}(Q\transpose Z R^{-1}) + (I-QQ\transpose) Z R^{-1} \nonumber \\
	& = Q\left(\rho_{\mathrm{skew}}(Q\transpose Z R^{-1}) - Q\transpose Z R^{-1} \right) + Z R^{-1} \nonumber\\
	& = Q\left(\rho_{\mathrm{skew}}((R^{-1})\transpose S\transpose Z R^{-1}) - (R^{-1})\transpose S\transpose Z R^{-1} \right) + Z R^{-1}, \label{eq:DRetrequalsplit}
\end{align}
where in the last step we used $XR^{-1} - Q = -SR^{-1}$.
Further note that for any matrix $M$ of size $p \times p$,
\begin{align}
	\|Q(\rho_{\mathrm{skew}}(M) - M)\|_{\mathrm{F}} = \|\mathrm{tril}(M) - \mathrm{tril}(M)\transpose - M\|_{\mathrm{F}} \leq 3 \|M\|_{\mathrm{F}}.
\end{align}
Hence, we have that,
\begin{align}
	\| \D\Retr_X(S)[Z] \|_{\mathrm{F}} & \geq \|ZR^{-1}\|_{\mathrm{F}} - 3\| (R^{-1})\transpose S\transpose Z R^{-1} \|_{\mathrm{F}} \nonumber\\
						& \geq \|Z\|_{\mathrm{F}} \left( \sigmamin(R^{-1}) - 3 \sigmamax(R^{-1})^2 \sigmamax(S) \right),
\end{align}
where we have used $\|A\|_{\mathrm{F}}\sigmamin(B) \leq \|AB\|_{\mathrm{F}} \leq \|A\|_{\mathrm{F}} \sigmamax(B)$ multiple times. Using the bounds on the singular values of $R^{-1}$ (derived in \eqref{eqn:sigmamaxbound}) we get that
\[
	\| \D\Retr_X(S)[Z] \|_{\mathrm{F}} \geq \|Z\|_{\mathrm{F}} \left( 1 - \frac{1}{2}\|S\|_{\mathrm{F}}^2 - 3 \|S\|_{\mathrm{F}} \right). % \geq 1 - 3 \|S\|_{\mathrm{F}} - \frac{1}{2}\|S\|_{\mathrm{F}}^2 -- this part missing Z, but we can just as well omit it all.
\]
Since this holds for all tangent vectors $Z$, we get that
\[ \sigmamin(\D\Retr_X(S)) \geq 1 - 3 \|S\|_{\mathrm{F}} - \frac{1}{2}\|S\|_{\mathrm{F}}^2.\]
To prove a better bound for the case of $p=1$ (the sphere), we improve the analysis of the expression derived in~\eqref{eq:DRetrequalsplit}. Note that for $p = 1$, the matrix inside the $\rho_{\mathrm{skew}}$ operator is a scalar, whose skew-symmetric part is necessarily zero. Also note that $Q$ is a single column matrix with value $\frac{x + s}{\|x + s\|}$ and $R = \|x + s\|$. Also, $X\transpose S X\transpose Z = 0$ since $S,Z$ are tangent. Therefore,
\begin{align*}
	\D\Retr_X(S)[Z] & = Z R^{-1} - Q(R^{-1})\transpose S\transpose Z R^{-1} \\
	 				& = \frac{1}{\|x+s\|} \left(z - \frac{s\transpose z}{1+\|s\|^2}(x+s)\right) \\
	 				& = \frac{1}{\|x+s\|} \left(z - \frac{s\transpose z}{1+\|s\|^2} s - \frac{s\transpose z}{1+\|s\|^2} x\right).
\end{align*}
Since $x$ is orthogonal to $s$ and $z$,
\begin{align*}
	\|\D\Retr_x(s)[z]\|^2 & = \frac{1}{1+\|s\|^2} \left( \left\|z - \frac{s\transpose z}{1+\|s\|^2} s \right\|^2 + \left( \frac{s\transpose z}{1+\|s\|^2}\right)^2 \right) \\
					   	    & = \frac{1}{1+\|s\|^2} \left( \|z\|^2 - 2\frac{(s\transpose z)^2}{1+\|s\|^2}  + \left( \frac{s\transpose z}{1+\|s\|^2}\right)^2 (1 + \|s\|^2) \right) \\
					   	    & = \frac{1}{1+\|s\|^2} \left( \|z\|^2 - \frac{(s\transpose z)^2}{1+\|s\|^2}\right) \\
					   	    & \geq \|z\|^2 \frac{1}{1+\|s\|^2} \left( 1 - \frac{\|s\|^2}{1+\|s\|^2} \right) \\
					   	    & = \|z\|^2 \frac{1}{(1+\|s\|^2)^2}.
\end{align*}
The worst-case scenario is achieved when $z$ and $s$ are aligned. Overall, we get
\[
	\|\D\Retr_x(s)[z]\| \geq \|z\| \frac{1}{1+\|s\|^2},
\]
which establishes the bound for the sphere. 
%\TODO{This is tricky business actually --- this is true: $\D\Retr_x(s) = \frac{1}{\|x+s\|} \Proj_{\Retr_x(s)}$, without square on the norm in the fraction. The projector is the restriction to $\T_x\calM$ of the orthogonal projector from all of $\Rn$ to $\T_{\Retr_x(s)}\calM$. It's true that this \emph{restricted} projector doesn't have a unit singular value, unless $s = 0$. That's where we get a square on the multiplier: this restricted projector has all its singular values smaller than 1 when s is not zero.}
\end{proof}

\subsection*{Differential of exponential map for manifolds with bounded curvature}

% \begin{proof}[Proof of Corollary~\ref{cor:mastercorollaryexponetial}]
% The case $C \leq 0$ is a direct consequence of Lemma~\ref{lemma:exponentialbound} and Corollary~\ref{cor:mastercorollary}. For the case when $C > 0$, note that 
% \[\forall x \in [0,\pi], \qquad \frac{\sin(x)}{x} \geq 1 - \frac{x}{\pi}.\] 
% Therefore we invoke Corollary~\ref{cor:mastercorollary} with $a = \pi\sqrt{\frac{b}{C}}$ and $b = 1 - \sqrt{\frac{\varepsilon C}{\pi^2 (\frac{L}{2} + \theta + \varsigma_{\max})}}$. Lemma~\ref{lemma:exponentialbound} implies that Assumption~\ref{assu:DRetr} is satisfied and the condition on $\varepsilon$ can also be checked under these assumptions. Corollary~\ref{cor:mastercorollary} now gives the result. 
% \end{proof}

Proposition~\ref{prop:mastercorollaryexponetial} regarding the differential of the exponential map on complete manifolds with bounded sectional curvature follows as a corollary of the following statement.
\begin{lemma} \label{lemma:exponentialbound}
	Assume all sectional curvatures of $\calM$, complete, are bounded above by $C$:
	\begin{itemize}
		\item[] If $C \leq 0$, then $\sigmamin(\D\Exp_x(s)) = 1$;
		\item[] If $C = \frac{1}{R^2} > 0$ and $\|s\| \leq \pi R$, then $1 \geq \sigmamin(\D\Exp_x(s)) \geq \frac{\sin(\|s\|/R)}{\|s\|/R}$.
	\end{itemize}
	As usual, we use the convention $\sin(t)/t = 1$ at $t = 0$.
\end{lemma}
\begin{proof}
	This results from a combination of few standard facts in Riemannian geometry:
	\begin{enumerate}
		\item \citep[Prop.~10.10]{lee2018riemannian} Given any two tangent vectors $s, \dot s \in \T_x\calM$, $J(t) = \D\Exp_x(ts)[t\dot s]$ is the unique Jacobi field along the geodesic $\gamma(t) = \Exp_x(ts)$ satisfying $J(0) = 0$ and $\Ddt J(0) = \dot s$.
		
		\item In particular, if $\dot s = \alpha s$ for some $\alpha \in \reals$ so that $\dot s$ and $s$ are parallel, then
		\begin{align*}
			J(t) & = \D\Exp_x(ts)][t\dot s] = \left. \ddq \Exp_x(ts + qt\dot s) \right|_{q=0} = \left. \ddq \gamma(t + q\alpha t) \right|_{q=0} = \alpha t \gamma'(t) = t P_{ts}(\dot s),
		\end{align*}
		using $\gamma'(t) = P_{ts}(s)$.
		It remains to understand the case where $\dot s$ is orthogonal to $s$.
		
		\item \citep[Prop.~10.12]{lee2018riemannian} If $\calM$ has constant sectional curvature $C$, $\|s\| = 1$ and $\inner{s}{\dot s} = 0$, the Jacobi field above is given by:
		\begin{align*}
			J(t) & = s_C(t) P_{ts}(\dot s),
		\end{align*}
		where $P_{ts}$ denotes parallel transport along $\gamma$ as in~\eqref{eq:Ps} and
		\begin{align*}
			s_C(t) & = \begin{cases}
				t & \textrm{ if } C = 0, \\
				R \sin(t/R) & \textrm{ if } C = \frac{1}{R^2} > 0, \textrm{ and} \\
				R \sinh(t/R) & \textrm{ if } C = -\frac{1}{R^2}.
			\end{cases}
		\end{align*}
		This can be reparameterized to allow for $\|s\| \neq 1$.
		Evaluating at $t = 1$ and using linearity in $\dot s$, we find for any $s, \dot s \in \T_x\calM$ that
		\begin{align}
			\D\Exp_x(s)[\dot s] & = P_s\!\left(\dot s_\parallel + \frac{s_C(\|s\|)}{\|s\|} \dot s_\perp \right),
		\end{align}
		where $\dot s_\perp$ is the part of $\dot s$ which is orthogonal to $s$ and $\dot s_\parallel$ is the part of $\dot s$ which is parallel to $s$---this corresponds to expression~\eqref{eq:DExpConstantCurvature}. By isometry of parallel transport, it is a simple exercise in linear algebra to deduce that
		\begin{align*}
			\sigmamin(\D\Exp_x(s)) & = \min\left( 1,  \frac{s_C(\|s\|)}{\|s\|} \right).
		\end{align*}
		
		\item \citep[Thm.~11.9(a)]{lee2018riemannian} Consider the case where $\dot s$ is orthogonal to $s$ of unit norm once again: the Jacobi field comparison theorem states that if the sectional curvatures of $\calM$ are upper-bounded by $C$, then $\|J(t)\|$ is at least as large as what it would be if $\calM$ had constant sectional curvature $C$---with the additional condition that $\|s\| \leq \pi R$ if $C = 1/R^2 > 0$. This leads to the conclusion through similar developments as above, using also~\citep[Prop.~10.7]{lee2018riemannian} to separate the components of $J(t)$ that are parallel or orthogonal to $\gamma'(t)$.
		\qedhere
	\end{enumerate}

%
%\TODO{-------------------------------------------}
%	
%The proof relies essentially on Jacobi fields and the Jacobi field comparison theorem, as detailed below. Inequalities hold as equalities if the sectional curvature is constant and equal to $C$ (use~\cite[Lemma~10.8]{lee1997riemannian} \TODO{MAKE ALL REFS TO BE ABOUT 2018 EDITION} instead of the comparison theorem in that case.)
%%
%Of course, for $s = 0$, $\D\Exp_x(0)$ is the identity operator, with all singular values equal to 1. Thus, we focus on $s \neq 0$.
%%Furthermore, because the exponential map is a smooth operator on the tangent bundle, the least singular value depends continuously on $x$ and $s$. Here, under some conditions on the curvature of $\calM$, we seek to bound how much smaller than 1 the least singular value may be as a function of $\|s\|$.
%
%Consider the smooth map $\Gamma \colon I_1 \times I_2 \to \calM$ defined by
%\begin{align*}
%	\Gamma(q, t) & = \Exp_x\left( \frac{t}{\|s\|}\left[ s + q\dot s \right] \right).
%\end{align*}
%(Intervals $I_1, I_2 \subseteq \reals$ are neighborhoods of 0, as large as $\Exp$ permits.)
%Notice that, for any $q$, $\Gamma(q, \cdot)$ defines a geodesic: we say that $\Gamma$ is a \emph{variation through geodesics} of the geodesic $\gamma(t) = \Gamma(0, t)$. Furthermore, $\gamma$ has unit speed since its velocity at $\gamma(0) = x$ is $\dot \gamma(0) = s/\|s\|$.
%%This object relates to our target $\D\Exp_x(s)[\dot s]$ is the following way.
%The following (smooth) vector field along $\gamma$ describes how $\gamma$ varies at every point:
%\begin{align*}
%	J(t) & = \left. \frac{\partial}{\partial q} \Gamma(q, t) \right|_{q = 0} = \frac{t}{\|s\|} \D\Exp_x\left(\frac{t}{\|s\|} s\right)[\dot s].
%\end{align*}
%Clearly, $J(t) \in \T_{\gamma(t)}\calM$ for all $t$.
%In particular,
%\begin{align*}
%	J(\|s\|) & = \D\Exp_x(s)[\dot s], % = \left. \D\Exp_x(s + q\dot s)[\dot s] \right|_{q = 0}§
%\end{align*}
%so that our goal is to understand the norm of $J(\|s\|)$.
%%Indeed, $\Gamma(q, \|s\|) = \Exp_x(s + q\dot s)$, so that the derivative of this curve with respect to $q$ at $q = 0$ describes how the end point $\Exp_x(s)$ moves (infinitesimally) when $s$ is perturbed in the direction $\dot s$. \TODO{A drawing would be better. I might have something. It's Fig 10.9 in Lee's intro to curvate.}
%
%The variation field $J$ of $\gamma$ is called a \emph{Jacobi field} along $\gamma$. A Jacobi field is called \emph{normal} along $\gamma$ if $J(t)$ is orthogonal to $\dot\gamma(t)$ for all $t$. By~\cite[Lemma~10.6]{lee1997riemannian}, this is the case if and only if $J(0)$ and $\dot J(0)$ are orthogonal to $\dot \gamma(0)$. In our case, $\dot \gamma(0) = s/\|s\|$ and
%\begin{align*}
%	J(0) & = 0, & \dot J(0) & = \frac{1}{\|s\|} \dot s,
%\end{align*}
%so that $J$ is normal if and only if $\dot s$ is orthogonal to $s$.
%We now have all the pieces to apply the following theorem (particularized to our situation).
%\begin{theorem}[\protect{Jacobi field comparison theorem,~\cite[Thm.~11.2]{lee1997riemannian}}]
%	Assume all sectional curvatures of $\calM$ are bounded above by $C$. If $J$ is normal (that is, if $\dot s$ is orthogonal to $s$), then
%	\begin{align*}
%		\|J(t)\| & \geq \begin{cases}
%		t \frac{\|\dot s\|}{\|s\|} & \textrm{ for } 0 \leq t, \textrm{ if } C = 0; \\
%		R \sin\left( \frac{t}{R} \right) \frac{\|\dot s\|}{\|s\|} & \textrm{ for } 0 \leq t \leq \pi R, \textrm{ if } C = \frac{1}{R^2} > 0; \\
%		R \sinh\left( \frac{t}{R} \right) \frac{\|\dot s\|}{\|s\|} & \textrm{ for } 0 \leq t, \textrm{ if } C = -\frac{1}{R^2} < 0.
%		\end{cases}
%	\end{align*}
%\end{theorem}
%\noindent In our setting, with $J(\|s\|) = \D\Exp_x(s)[\dot s]$, this means that if $\dot s$ is orthogonal to $s$, then
%\begin{align}
%	\|\D\Exp_x(s)[\dot s]\| & \geq \begin{cases}
%	\|\dot s\| & \textrm{ if } C = 0; \\
%	\frac{\sin\left( \|s\| / R \right)}{\|s\| / R} \|\dot s\| & \textrm{ if } \|s\| \leq \pi R \textrm{ and } C = \frac{1}{R^2} > 0; \\
%	\frac{\sinh\left( \|s\| / R \right)}{\|s\| / R} \|\dot s\| & \textrm{ if } C = -\frac{1}{R^2} < 0,
%	\end{cases}
%	\label{eq:boundsDexpinnerproof}
%\end{align}
%and $\D\Exp_x(s)[\dot s]$ is orthogonal to $\dot \gamma(\|s\|)$ owing to normality of $J$.
%
%On the other hand, $J$ is called \emph{tangential} if $J(t)$ is a multiple of $\dot \gamma(t)$ for all $t$. This is the case if $\dot s = \alpha s$ for some $\alpha \in \reals$. Indeed, in this case,
%\begin{align*}
%	\Gamma(q, t) = \Exp_x\left(\frac{t}{\|s\|}\left[ 1 + \alpha q \right]s \right) = \gamma((1+\alpha q)t).
%\end{align*}
%That is, $\Gamma(q, \cdot)$ is a geodesic with speed $1+\alpha q$. By the chain rule,
%\begin{align*}
%	J(t) & = \left. \frac{\partial}{\partial q} \Gamma(q, t) \right|_{q = 0} = \alpha t \dot \gamma(t).
%\end{align*}
%As announced, $J$ is tangential. Furthermore, $\|\dot \gamma(t)\| = 1$ for all $t$, so that
%\begin{align*}
%	\|\D\Exp_x(s)[\dot s]\| & = \big\| J(\|s\|) \big\| = \alpha \|s\| = \|\dot s\|.
%\end{align*}
%To summarize, we found $\D\Exp_x(s)$ is a linear operator such that $\D\Exp_x(s)[s] = \|s\|\dot \gamma(\|s\|)$ and such that if $\dot s \perp s$, then $\D\Exp_x(s)[\dot s] \perp \D\Exp_x(s)[s]$ (where $u \perp v$ means $u$ is orthogonal to $v$). The following elementary lemma applies.
%\begin{lemma}
%	Consider a linear operator $A \colon \calE \to \calF$ between two Euclidean spaces such that there exists $v \in \calE, u \in \calF$ of unit norm and $\sigma \geq 0$ with $Av = \sigma u$, and such that $Az \perp Av$ for all $z \in \calE$ with $z \perp v$. Then, $(u, v, \sigma)$ is a singular triplet of $A$.
%\end{lemma}
%\begin{proof}
%	We work in coordinates: $A$ is a matrix and $u, v$ are column vectors. Let $U_\perp$ and $V_\perp$ be such that $\begin{bmatrix} u & U_\perp \end{bmatrix}$ and $\begin{bmatrix} v & V_\perp \end{bmatrix}$ are orthogonal matrices. Then, $\begin{bmatrix} u & U_\perp \end{bmatrix}\transpose  A \begin{bmatrix} v & V_\perp \end{bmatrix} = \begin{bmatrix} \sigma & 0 \\ 0 & K \end{bmatrix}$ for some matrix $K$, owing to our assumptions. With the SVD decomposition $K = U\Sigma V\transpose $, we find
%	\begin{align*}
%		A & = \begin{bmatrix} u & U_\perp \end{bmatrix} \begin{bmatrix} 1 & 0 \\ 0 & U \end{bmatrix} \begin{bmatrix} \sigma & 0 \\ 0 & \Sigma \end{bmatrix} \begin{bmatrix} 1 & 0 \\ 0 & V \end{bmatrix}\transpose  \begin{bmatrix} v & V_\perp \end{bmatrix}\transpose  \\
%			& = \begin{bmatrix} u & U_\perp U \end{bmatrix} \begin{bmatrix} \sigma & 0 \\ 0 & \Sigma \end{bmatrix} \begin{bmatrix} v & V_\perp V \end{bmatrix}\transpose.
%	\end{align*}
%	Up to permutations, this is an SVD of $A$, which concludes the proof.
%\end{proof}
%Thus, $\D\Exp_x(s)$ admits 1 as a singular value, and all other singular values are bounded as per~\eqref{eq:boundsDexpinnerproof}. To conclude, observe that $\sinh(t)/t \geq 1$ for all $t$. % To prove it, consider the Taylor expansion of sinh: it's x + all odd powers of x with positive coefficient, so when divide by x you get 1 + all even powers of x, with positive coefficients.

\end{proof}

\subsection*{Extending to general retractions}

In order to prove Theorem~\ref{thm:retractionsgeneralmainnew}, we first introduce a result from topology. We follow~\cite{berge1963topological}, including the blanket assumption that all encountered topological spaces are Hausdorff (page 65 in that reference)---this is the case for us so long as the topology of $\calM$ itself is Hausdorff, which most authors require as part of the definition of a smooth manifold. Products of topological spaces are equipped with the product topology. Neighborhoods are open. A \emph{correspondence} $\Gamma \colon Y \to Z$ maps points in $Y$ to subsets of $Z$.
\begin{definition}[Upper semicontinuous (u.s.c.)\ mapping] \label{def:uscmapping}
	A correspondence $\Gamma \colon Y \rightarrow Z$ between two topological spaces $Y, Z$ is a \emph{u.s.c.\ mapping} if, for all $y$ in $Y$, $\Gamma(y)$ is a compact subset of $Z$ and, for any neighborhood $V$ of $\Gamma(y)$, there exists a neighborhood $U$ of $y$ such that, for all $u \in U$, $\Gamma(u) \subseteq V$. 
\end{definition}
\begin{theorem}[{\protect\citet[Thm.~VI.2, p116]{berge1963topological}}] \label{thm:bergemain}
	If $\phi$ is an upper semicontinuous, real-valued function in $Y \times Z$ and $\Gamma$ is a u.s.c.\ mapping of $Y$ into $Z$ (two topological spaces) such that $\Gamma(y)$ is nonempty for each $y$, then the real-valued function $M$ defined by
	\begin{align*}
	M(y) = \max_{z \in \Gamma(y)} \phi(y,z)
	\end{align*}
	is upper semicontinuous. (Under the assumptions, the maximum is indeed attained.)
	%Here, $Y$ and $Z$ are (Hausdorff) topological spaces.
\end{theorem}

We use the above theorem to establish our result. Manifolds (including tangent bundles) are equipped with the natural topology inherited from their smooth structure.

\begin{proof}[Proof of Theorem~\ref{thm:retractionsgeneralmainnew}]
	It is sufficient to show that the function
	\begin{align}
	t(r) & = \inf_{(x, s) \in \T\calM : x \in \calU, \|s\|_x \leq r} \sigmamin(\D\Retr_x(s))
	\end{align}
	is lower semicontinuous from $\reals^+ = \{r \in \reals : r \geq 0 \}$ to $\reals$, with respect to their usual topologies. Indeed, $t(0) = 1$ owing to the fact that $\D\Retr_x(0)$ is the identity map for all $x$, and $t$ being lower semicontinuous means that it cannot ``jump down''. Explicitly, lower semicontinuity at $r = 0$ implies that, for all $\delta > 0$, there exists $a > 0$ such that for all $r \leq a$ we have $t(r) \geq t(0) - \delta = 1 - \delta \triangleq b$. % inequalities here are fine, see https://en.wikipedia.org/wiki/Semi-continuity#Formal_definition
	
	To this end, consider the correspondence $\Gamma \colon \reals^+ \to \T\calM$ defined by
	\begin{align}
	\Gamma(r) & = \{ (x, s) \in \T\calM : x \in \calU \textrm{ and } \|s\|_x \leq r \}.
	\end{align}
	%where $\T\calM$ is equipped with its usual topology, induced by its manifold structure.
	Further consider the function $\phi \colon \reals^+ \times \T\calM \to \reals$ defined by $\phi(r, (x, s)) = -\sigmamin(\D\Retr_x(s))$. Then, $t(r) = -M(r)$, where
	\begin{align}
	M(r) & = \sup_{(x, s) \in \Gamma(r)} \phi(r, (x, s)).
	\label{eq:Mrsup}
	\end{align}
	Thus, we must show $M$ is upper semicontinuous. By Theorem~\ref{thm:bergemain}, this is the case if
	\begin{enumerate}
		\item $\phi$ is upper semicontinuous,
		\item $\Gamma(r)$ is nonempty and compact for all $r \geq 0$, and
		\item For any $r \geq 0$ and any neighborhood $\calV$ of $\Gamma(r)$ in $\T\calM$, there exists a neighborhood $I$ of $r$ in $\reals^+$ such that, for all $r' \in I$, we have $\Gamma(r') \subseteq \calV$.
	\end{enumerate}
	The first condition holds a fortiori since $\phi$ is continuous, owing to smoothness of $\Retr \colon \T\calM \to \calM$. The second condition holds since $\calU$ is nonempty and compact.
	% (In particular, this shows the supremum in the definition of $M$~\eqref{eq:Mrsup} is attained.)
	For the third condition, we show in Lemma~\ref{lem:neighborhoodballsTM} below that there exists a continuous function $\Delta \colon \calU \to \reals$ (continuous with respect to the subspace topology) such that $\{ (x, s) \in \T\calM : x \in \calU \textrm{ and } \|s\|_x \leq \Delta(x) \} \subseteq \calV$ and $\Delta(x) > r$ for all $x \in \calU$ (if $\calM$ is not connected, apply the lemma to each connected component which intersects with $\calU$). As a result, $\min_{x\in\calU} \Delta(x) = r + \varepsilon$ for some $\varepsilon > 0$ (using $\calU$ compact), and $\Gamma(r+\varepsilon)$ is included in $\calV$. We conclude that $I = [0, r+\varepsilon)$ is a suitable neighborhood of $r$ to verify the condition.
\end{proof}

We now state and prove the last piece of the puzzle, which applies above with $r(x)$ constant ($L = 0$). Although the context is quite different, the first part of the proof is inspired by that of the tubular neighborhood theorem in~\citep[Thm.~5.25]{lee2018riemannian}.

\begin{lemma} \label{lem:neighborhoodballsTM}
	Let $\calU$ be any subset of a connected Riemannian manifold $\calM$ and let $r \colon \calU \to \reals^+$ be $L$-Lipschitz continuous with respect to the Riemannian distance $\dist$ on $\calM$, that is, %\footnote{If $x, x'$ belong to distinct connected components of $\calM$, set $\dist(x, x') = \infty$.}
	\begin{align*}
	\forall x, x' \in \calU, \quad |r(x) - r(x')| \leq L \dist(x, x').
	\end{align*}
	Consider this subset of the tangent bundle:
	\begin{align*}
	\left\{ (x, s) \in \T\calM : x \in \calU \textrm{ and } \|s\|_x \leq r(x) \right\}.
	\end{align*}
	For any neighborhood $\calV$ of this set in $\T\calM$, there exists an $(L+1)$-Lipschitz continuous function $\Delta \colon \calU \to \reals^+$ such that $\Delta(x) > r(x)$ for all $x \in \calU$ and
	\begin{align*}
	\left\{ (x, s) \in \T\calM : x \in \calU \textrm{ and } \|s\|_x \leq \Delta(x) \right\} \subseteq \calV.
	\end{align*}
\end{lemma}
\begin{proof}
	Consider the following open subsets of the tangent bundle, defined for each $x \in \calM$ and $\delta \in \reals$:
	\begin{align*}
	V_\delta(x) & = \left\{ (x', s') \in \T\calM : \dist(x, x') < \delta - r(x) \textrm{ and } \|s'\|_{x'} < \delta \right\}.
	\end{align*}
	Referring to these sets, define the function $\Delta \colon \calU \to \reals$ as:
	\begin{align*}
	\Delta(x) & = \sup\left\{ \delta \in \reals : V_\delta(x) \subseteq \calV \right\}.
	\end{align*}
	This is well defined since $V_{r(x)}(x) = \emptyset$, so that $\Delta(x) \geq r(x)$ for all $x$.
	If $\Delta(x) = \infty$ for some $x$, then $\calV = \T\calM$ and the claim is clear (for example, redefine $\Delta(x) = r(x) + 1$ for all $x$). Thus, we assume $\Delta(x)$ finite for all $x$. The rest of the 	proof is in two parts.
	
	\textbf{Step 1: $\Delta$ is Lipschitz continuous.} Pick $x, x' \in \calU$, arbitrary. We must show
	\begin{align*}
	\Delta(x) - \Delta(x') \leq (L+1) \dist(x, x').
	\end{align*}
	Then, by reversing the roles of $x$ and $x'$, we get $|\Delta(x) - \Delta(x')| \leq (L+1) \dist(x, x')$, as desired. If $\Delta(x) \leq (L+1) \dist(x, x')$, the claim is clear since $\Delta(x') \geq 0$. Thus, we now assume $\Delta(x) > (L+1) \dist(x, x')$. Define $\delta = \Delta(x) - (L+1) \dist(x, x') > 0$. It is sufficient to show that $V_\delta(x') \subseteq \calV$, as this implies $\Delta(x') \geq \delta = \Delta(x) - (L+1) \dist(x, x')$, allowing us to conclude. To this end, we show the first inclusion in:
	\begin{align*}
	V_\delta(x') \subseteq V_{\Delta(x)}(x) \subseteq \calV.
	\end{align*}
	Consider an arbitrary $(x'', s'') \in V_\delta(x')$. This implies two things: first, $\|s''\|_{x''} < \delta \leq \Delta(x)$, and second:
	\begin{align*}
	\dist(x'', x) & \leq \dist(x'', x') + \dist(x', x) \\ & < \delta - r(x') + \dist(x', x) \\ & = \Delta(x) - r(x) + r(x) - r(x') - L \dist(x, x') \\ & \leq \Delta(x) - r(x),
	\end{align*}
	where in the last step we used $r(x) - r(x') \leq L\dist(x, x')$ since $r$ is $L$-Lipschitz continuous on $\calU$.
	As a result, $(x'', s'')$ is in $V_{\Delta(x)}(x)$, which concludes this part of the proof.
	
	
	\textbf{Step 2: $\Delta(x) > r(x)$ for all $x \in \calU$.}
	% WRONG -- The short version of the argument goes like this: by definition, $I_x = \left\{ \delta \in \reals : V_\delta(x) \subseteq \calV \right\}$ is either open of the form $(-\infty, \Delta(x))$, or closed of the form $(-\infty, \Delta(x)]$. Since we know $I_x$ contains $r(x)$, to show $\Delta(x) > r(x)$ we only need to show that $I_x$ is open. Intuitively, this follows from the fact that $V_\delta(x)$ is open for each $\delta$ and $\calV$ is open. We give more details below. \TODO{Not sure if should keep this intuition, because the proof doesn't actually show that $I_x$ is open; it only shows that $\Delta(x) > r(x)$. Might be able to extract it from there, but should be sure. Actually, I think this is WRONG: to simplify, take calV to be the open unit disk in R^2, and take V_delta to be the open disk on radius delta in R^2: the set of deltas such that V_delta is in calV is simply (-\infty, 1], closed. So just drop this intuition: the proof below works out.}
	%
	Pick $x \in \calU$, arbitrary: $\calV$ is a neighborhood of
	\begin{align}
	\left\{ (x, s) \in \T\calM : \|s\|_x \leq r(x) \right\}.
	\label{eq:DeltalargerthanrmanifoldsetA}
	\end{align}
	The claim is that there exists $\varepsilon > 0$ such that
	\begin{align}
	\left\{ (x', s') \in \T\calM : \dist(x, x') \leq \varepsilon \textrm{ and } \|s'\|_{x'} \leq r(x) + \varepsilon \right\}
	\label{eq:DeltalargerthanrmanifoldsetB}
	\end{align}
	is included in $\calV$. Indeed, that would show that $\Delta(x) \geq r(x) + \varepsilon > r(x)$. To show this, we construct special coordinates on $\T\calM$ around $x$.
	
	The (inverse of the) exponential map at $x$ restricted to tangent vectors of norm strictly less than $\inj(x)$ (the injectivity radius at $x$) provides a diffeomorphism $\varphi$ from $\calW \subseteq \calM$ (the open geodesic ball of radius $\inj(x)$ around $x$) to $B(0, \inj(x))$: the open ball centered around the origin in the Euclidean space $\Rd$, where $d = \dim\calM$. Additionally, from the chart $(\calW, \varphi)$, we extract coordinate vector fields on $\calW$: a set of smooth vector fields $W_1, \ldots, W_d$ on $\calW$ such that, at each point in $\calW$, the corresponding tangent vectors form a basis for the tangent space. We further orthonormalize this local frame (see~\citep[Prop.~2.8]{lee2018riemannian}) into a new local frame, $E_1, \ldots, E_d$, so that for each $x' \in \calW$ we have that $E_1(x'), \ldots, E_d(x')$ form an orthonormal basis for $\T_{x'}\calM$ (with respect to the Riemannian metric at $x'$). Then, the map
	\begin{align*}
	\psi(x', s') & = \left( \varphi(x'), \zeta(x', s') \right) & & \textrm{ with } & \zeta(x', s') & = \left( \inner{E_1(x')}{s'}_{x'}, \ldots, \inner{E_d(x')}{s'}_{x'} \right)
	\end{align*}
	establishes a diffeomorphism between $\T\calW$ and $B(0, \inj(x)) \times \Rd$, with the following properties:
	\begin{enumerate}
		\item $\dist(x, x') = \|\varphi(x')\|$ (in particular, $\varphi(x) = 0$), and
		\item For any $s', v' \in \T_{x'}\calM$, it holds $\inner{s'}{v'}_{x'} = \inner{\zeta(x', s')}{\zeta(x', v')}$.
	\end{enumerate}
	(Here, $\inner{\cdot}{\cdot}$ and $\|\cdot\|$ denote the Euclidean inner product and norm in $\Rd$.)
	
	Expressed in these coordinates (that is, mapped through $\psi$), the set in~\eqref{eq:DeltalargerthanrmanifoldsetA} becomes:
	\begin{align*}
	D_0 & = \{ 0 \} \times \bar B(0, r(x)),
	\end{align*}
	where $\bar B(0, r(x))$ denotes the closed Euclidean ball of radius $r(x)$ around the origin in $\Rd$. Of course, $\calV \cap \T\calW$ maps to a neighborhood of $D_0$ in $\Rd \times \Rd$: call it $O$. Similarly, the set in~\eqref{eq:DeltalargerthanrmanifoldsetB} maps to:
	\begin{align*}
	D_\varepsilon & = \bar B(0, \varepsilon) \times \bar B(0, r(x) + \varepsilon).
	\end{align*}
	It remains to show that there exists $\varepsilon > 0$ such that $D_\varepsilon$ is included in $O$. %: this is a simple exercise in topology.
	
	Use this distance on $\Rd \times \Rd$: $\dist((y, z), (y', z')) = \max(\|y - y'\|, \|z - z'\|)$. This distance is compatible with the usual topology. For each $(0, z)$ in $D_0$, there exists $\varepsilon_z > 0$ such that
	\begin{align*}
	C(z, \varepsilon_z) = \left\{ (y', z') \in \Rd \times \Rd : \|y'\| < \varepsilon_z \textrm{ and } \|z - z'\| < \varepsilon_z \right\}
	\end{align*}
	is included in $O$ (this is where we use the fact that $\calV$---hence $O$---is open). The collection of open sets $C(z, \varepsilon_z/2)$ forms an open cover of $D_0$. Since $D_0$ is compact, we may extract a finite subcover, that is, we select $z_1, \ldots, z_n$ such that the sets $C(z_i, \varepsilon_{z_i}/2)$ cover $D_0$. Now, define $\varepsilon = \min_{i = 1, \ldots, n} \varepsilon_{z_i}/2$ (necessarily positive), and consider any point $(y, z) \in D_\varepsilon$. We must show that $(y, z)$ is in $O$. To this end, let $\bar z$ denote the point in $\bar B(0, r(x))$ which is closest to $z$. Since $(0, \bar z)$ is in $D_0$, there exists $i$ such that $(0, \bar z)$ is in $C(z_i, \varepsilon_{z_i}/2)$. As a result,
	\begin{align*}
	\|z - z_i\| & \leq \|z - \bar z\| + \|\bar z - z_i\| < \varepsilon + \varepsilon_{z_i}/2 \leq \varepsilon_{z_i}.
	\end{align*}
	Likewise, $\|y\| \leq \varepsilon \leq \varepsilon_{z_i}/2 < \varepsilon_{z_i}$. Thus, we conclude that $(y, z)$ is in $C(z_i, \varepsilon_{z_i})$, which is included in $O$. This confirms $D_\varepsilon$ is in $O$, so that the set in~\eqref{eq:DeltalargerthanrmanifoldsetB} is in $\calV$ for some $\varepsilon > 0$.
	% see notes August 15, 2019.
\end{proof}


%
% Make this an exercise in the book; I already moved the latex code there somewhere (Oct 9, 2019).
%
%\begin{lemma} \label{lem:compactballbundle}
%	Let $\calU$ be any subset of a Riemannian manifold $\calM$ and let $r \colon \calU \to \reals^+$ % need R^+ rather than just R so that pi(calT) = calU; otherwise, might be missing bits.
%	be continuous with respect to the subspace topology. The set $\calU$ is compact in $\calM$ if and only if the set
%	\begin{align*}
%	\calT & = \left\{ (x, s) \in \T\calM : x \in \calU \textrm{ and } \|s\|_x \leq r(x) \right\}
%	\end{align*}
%	is compact in $\T\calM$.
%\end{lemma}
%\begin{proof}
%	Let $\pi \colon \T\calM \to \calM$ be the natural projection $\pi(x, s) = x$. If $\calT$ is compact, then $\calU = \pi(\calT)$ is compact since $\pi$ is continuous.
%	
%	For the other direction, we now assume $\calU$ is compact. Consider the collection of open sets $B(x, \inj(x)/2)$ (open geodesic balls centered at $x \in \calU$ with radius given by half the injectivity radius of $\calM$ at $x$): these form an open cover of $\calU$. By compactness, we can extract a finite subcover, that is: we select $x_1, \ldots, x_n \in \calU$ for some finite $n$ such that the union of the $n$ balls $B(x_i, \inj(x_i)/2)$, $i = 1, \ldots, n$, contains $\calU$. Define the larger open balls $\calW_i = B(x_i, \inj(x_i))$. For each $x_i$, setup the local coordinates $\psi_i = (\varphi_i, \zeta_i)$ from $\T\calW_i$ to $B(0, \inj(x_i)) \times \Rd$ as in Step 2 of the proof of Lemma~\ref{lem:neighborhoodballsTM}. Further define closed geodesic balls $\calW_i' = \bar B(x_i, \inj(x_i)/2)$. Clearly, $\calU$ is included in the union of these closed balls, so that
%	\begin{align*}
%	\calV & = \bigcup_{i = 1, \ldots, n} \calV_i, & \textrm{ with }  & & \calV_i & = \left\{ (x, s) \in \T\calM : x \in \calW_i' \cap \calU \textrm{ and } \|s\|_x \leq r(x) \right\}.
%	\end{align*}
%	To conclude, it is sufficient to show that each $\calV_i$ is compact in $\T\calM$. This is the case if and only if each $\psi_i(\calV_i)$ is compact in $\Rd \times \Rd$. This is indeed true since
%	\begin{align*}
%	\psi_i(\calV_i) & = \{ (y, z) \in \Rd \times \Rd : y \in \varphi_i(\calW_i' \cap \calU) \textrm{ and } \|z\| \leq r(\varphi_i^{-1}(y)) \}
%	\end{align*}
%	is closed and bounded in $\Rd \times \Rd$. To see this, notice that $\varphi_i(\calW_i' \cap \calU)$ is compact by continuity of $\varphi_i$ and compactness of $\calW_i'$ and $\calU$, and $r \circ \varphi_i^{-1}$ is continuous.
%\end{proof}
